% !TEX program = lualatex
% ---------------------------------------------------------------------------------------------
%  Readback note 21: The seesaw identity for the chain of theta lifts
%
%  Build:  python3 tools/rb.py note <slug>
%
%  One unit of a Lean artifact, rewritten as ordinary mathematics. The prose is hand-written. Lean names,
%  line ranges and counts come from data.tex; quoted Lean statements are copied from the artifact by the build.
%
%  Rules of the text: standard terminology; every symbol is explained where it first appears and listed again
%  in the notation at the end; no letter has two meanings, except as the notation lists; a proof is complete or says what it takes for granted;
%  Lean names appear only in the LEAN lines and in the concordance.
%
%  The signature rule (the build enforces it): a function, relation or predicate that the artifact defines is
%  written only through its notation macro, which the build makes from the spec (`notation`) with one slot for
%  each explicit argument of the Lean definition. No argument is fixed and dropped, also not "for brevity" and
%  not for one section. Every argument and every hypothesis of a stated declaration stands inside its
%  statement, or in a Setting that the statement cites by number; only implicit and instance arguments (the
%  ambient objects) may stand in the conventions.
% ---------------------------------------------------------------------------------------------
\documentclass[11pt,a4paper]{article}
\ifdefined\pdfvariable\pdfvariable suppressoptionalinfo \numexpr 1+2+4+8\relax\fi
\usepackage[a4paper,hmargin=35mm,top=22mm,bottom=27mm,footskip=12mm,nohead]{geometry}

\usepackage{fontspec}
\directlua{
  luaotfload.add_fallback("gfallback", {"STIXTwoMath-Regular.otf:mode=node;", "DejaVuSans.ttf:mode=node;"})
  luaotfload.add_fallback("gmonofallback", {"DejaVuSansMono.ttf:mode=node;", "STIXTwoMath-Regular.otf:mode=node;"})
}
\setmainfont{NewCM10}[Extension=.otf,UprightFont=*-Book,ItalicFont=*-BookItalic,
  BoldFont=*-Bold,BoldItalicFont=*-BoldItalic,SmallCapsFont=*-Book,SmallCapsFeatures={Letters=SmallCaps},
  RawFeature={fallback=gfallback}]
\setsansfont{NewCMSans10}[Extension=.otf,UprightFont=*-Book,ItalicFont=*-BookOblique,
  BoldFont=*-Bold,BoldItalicFont=*-BoldOblique,RawFeature={fallback=gfallback}]
\setmonofont{JetBrainsMono}[Extension=.otf,UprightFont=*-Regular,ItalicFont=*-Italic,
  BoldFont=*-SemiBold,BoldItalicFont=*-SemiBoldItalic,Scale=MatchLowercase,
  RawFeature={-calt;-liga;fallback=gmonofallback}]
\usepackage{amsmath}
\usepackage{amsthm}
\usepackage[math-style=TeX]{unicode-math}
\setmathfont{NewCMMath-Book.otf}
\emergencystretch=2.5em\relax

\usepackage{xcolor}
\usepackage{tikz}
\usetikzlibrary{calc,arrows.meta,positioning,backgrounds}
\usepackage[protrusion=true,expansion=true]{microtype}
\usepackage{enumitem}
\usepackage{needspace}
\usepackage{booktabs}
\usepackage{array}
\usepackage{longtable}
\usepackage{fancyhdr}
\usepackage{lastpage}
\usepackage{multicol}
\usepackage[hidelinks,pdfusetitle]{hyperref}

\definecolor{ink}{HTML}{0F172A}
\definecolor{inksoft}{HTML}{475569}
\definecolor{inkfaint}{HTML}{94A3B8}
\definecolor{rulegrey}{HTML}{CBD5E1}
\definecolor{arith}{HTML}{C2410C}
\definecolor{est}{HTML}{7C3AED}
\definecolor{ana}{HTML}{0891B2}
\definecolor{ok}{HTML}{047857}
\definecolor{leancol}{HTML}{3730A3}

\title{The seesaw identity for the chain of theta lifts}
\author{grh-compact}
\date{}

% ------------------------------------------------------------------------------------- macros
\newcommand{\leanfont}{\ttfamily\fontsize{8.6}{10.5}\selectfont\color{leancol}}
\newcommand{\lean}[1]{\csname LN@#1\endcsname}
\newcommand{\fillin}[1]{{\color{arith}\itshape[#1]}}  % a place of the template that is not written yet; the build reports every one that is left
\newcommand{\leanquote}[1]{\csname LQ@#1\endcsname}
% generated by rb.py note; do not edit
\newcommand{\ConcDecls}{80}
\newcommand{\Pin}{72fead10}
\newcommand{\PinText}{that of commit {\leanfont 72fead10}}
\newcommand{\ArtifactName}{grh-compact}
\expandafter\newcommand\csname LN@fillingfits\endcsname{\href{https://github.com/lsernau/grh-lean/blob/72fead102f663d5c7a407964a13dc0cf9d977207/grh-compact/src/New/Seesaw/BesselGeneric.lean\#L14-L22}{{\leanfont Data.\allowbreak{}filling\_fits}}}
\expandafter\newcommand\csname LN@columnCoords\endcsname{\href{https://github.com/lsernau/grh-lean/blob/72fead102f663d5c7a407964a13dc0cf9d977207/grh-compact/src/New/Seesaw/BesselGeneric.lean\#L27-L28}{{\leanfont Data.\allowbreak{}columnCoords}}}
\expandafter\newcommand\csname LN@goodProjTranslate\endcsname{\href{https://github.com/lsernau/grh-lean/blob/72fead102f663d5c7a407964a13dc0cf9d977207/grh-compact/src/New/Seesaw/BesselGeneric.lean\#L32-L33}{{\leanfont Data.\allowbreak{}goodProjTranslate}}}
\expandafter\newcommand\csname LN@poincareratmul\endcsname{\href{https://github.com/lsernau/grh-lean/blob/72fead102f663d5c7a407964a13dc0cf9d977207/grh-compact/src/New/Seesaw/BesselGeneric.lean\#L40-L41}{{\leanfont Filling.\allowbreak{}Fits.\allowbreak{}poincare\_rat\_mul}}}
\expandafter\newcommand\csname LN@existsbesselFinmulmemofnezero\endcsname{\href{https://github.com/lsernau/grh-lean/blob/72fead102f663d5c7a407964a13dc0cf9d977207/grh-compact/src/New/Seesaw/BesselGeneric.lean\#L46-L50}{{\leanfont Filling.\allowbreak{}Fits.\allowbreak{}exists\_besselFin\_mul\_mem\_of\_ne\_zero}}}
\expandafter\newcommand\csname LN@poincarebesselnormalform\endcsname{\href{https://github.com/lsernau/grh-lean/blob/72fead102f663d5c7a407964a13dc0cf9d977207/grh-compact/src/New/Seesaw/BesselGeneric.lean\#L58-L95}{{\leanfont Filling.\allowbreak{}Fits.\allowbreak{}poincare\_bessel\_normal\_form}}}
\expandafter\newcommand\csname LN@isGenericpoincare\endcsname{\href{https://github.com/lsernau/grh-lean/blob/72fead102f663d5c7a407964a13dc0cf9d977207/grh-compact/src/New/Seesaw/BesselGeneric.lean\#L101-L123}{{\leanfont Filling.\allowbreak{}Fits.\allowbreak{}isGeneric\_poincare}}}
\expandafter\newcommand\csname LN@flagOrder\endcsname{\href{https://github.com/lsernau/grh-lean/blob/72fead102f663d5c7a407964a13dc0cf9d977207/grh-compact/src/New/Seesaw/EqualLinesOrbit.lean\#L18-L18}{{\leanfont flagOrder}}}
\expandafter\newcommand\csname LN@mulapplystringofblockTriangular\endcsname{\href{https://github.com/lsernau/grh-lean/blob/72fead102f663d5c7a407964a13dc0cf9d977207/grh-compact/src/New/Seesaw/EqualLinesOrbit.lean\#L24-L27}{{\leanfont mul\_apply\_string\_of\_blockTriangular}}}
\expandafter\newcommand\csname LN@IsAdaptedBasis\endcsname{\href{https://github.com/lsernau/grh-lean/blob/72fead102f663d5c7a407964a13dc0cf9d977207/grh-compact/src/New/Seesaw/EqualLinesOrbit.lean\#L31-L34}{{\leanfont Grading.\allowbreak{}IsAdaptedBasis}}}
\expandafter\newcommand\csname LN@IsAdaptedAt\endcsname{\href{https://github.com/lsernau/grh-lean/blob/72fead102f663d5c7a407964a13dc0cf9d977207/grh-compact/src/New/Seesaw/EqualLinesOrbit.lean\#L38-L38}{{\leanfont Grading.\allowbreak{}IsAdaptedAt}}}
\expandafter\newcommand\csname LN@HasNearLines\endcsname{\href{https://github.com/lsernau/grh-lean/blob/72fead102f663d5c7a407964a13dc0cf9d977207/grh-compact/src/New/Seesaw/EqualLinesOrbit.lean\#L42-L44}{{\leanfont Grading.\allowbreak{}HasNearLines}}}
\expandafter\newcommand\csname LN@existsisAdaptedAt\endcsname{\href{https://github.com/lsernau/grh-lean/blob/72fead102f663d5c7a407964a13dc0cf9d977207/grh-compact/src/New/Seesaw/EqualLinesOrbit.lean\#L46-L49}{{\leanfont Grading.\allowbreak{}exists\_isAdaptedAt}}}
\expandafter\newcommand\csname LN@gramstringentry\endcsname{\href{https://github.com/lsernau/grh-lean/blob/72fead102f663d5c7a407964a13dc0cf9d977207/grh-compact/src/New/Seesaw/EqualLinesOrbit.lean\#L54-L60}{{\leanfont Grading.\allowbreak{}IsAdaptedAt.\allowbreak{}gram\_string\_entry}}}
\expandafter\newcommand\csname LN@inv\endcsname{\href{https://github.com/lsernau/grh-lean/blob/72fead102f663d5c7a407964a13dc0cf9d977207/grh-compact/src/New/Seesaw/EqualLinesOrbit.lean\#L62-L64}{{\leanfont Grading.\allowbreak{}IsAdaptedBasis.\allowbreak{}inv}}}
\expandafter\newcommand\csname LN@ideleLogValle\endcsname{\href{https://github.com/lsernau/grh-lean/blob/72fead102f663d5c7a407964a13dc0cf9d977207/grh-compact/src/New/Seesaw/EqualLinesOrbit.lean\#L68-L73}{{\leanfont Grading.\allowbreak{}IsAdaptedBasis.\allowbreak{}ideleLogVal\_le}}}
\expandafter\newcommand\csname LN@ideleLogValeq\endcsname{\href{https://github.com/lsernau/grh-lean/blob/72fead102f663d5c7a407964a13dc0cf9d977207/grh-compact/src/New/Seesaw/EqualLinesOrbit.lean\#L75-L77}{{\leanfont Grading.\allowbreak{}IsAdaptedBasis.\allowbreak{}ideleLogVal\_eq}}}
\expandafter\newcommand\csname LN@existsmembesselFin\endcsname{\href{https://github.com/lsernau/grh-lean/blob/72fead102f663d5c7a407964a13dc0cf9d977207/grh-compact/src/New/Seesaw/EqualLinesOrbit.lean\#L84-L106}{{\leanfont Grading.\allowbreak{}IsAdaptedAt.\allowbreak{}exists\_mem\_besselFin}}}
\expandafter\newcommand\csname LN@isAdmissibleEvalpairEvalY\endcsname{\href{https://github.com/lsernau/grh-lean/blob/72fead102f663d5c7a407964a13dc0cf9d977207/grh-compact/src/New/Seesaw/GoodDatumExists.lean\#L21-L38}{{\leanfont Data.\allowbreak{}isAdmissibleEval\_pairEvalY}}}
\expandafter\newcommand\csname LN@existsisAdmissibleCombreducedVacuum\endcsname{\href{https://github.com/lsernau/grh-lean/blob/72fead102f663d5c7a407964a13dc0cf9d977207/grh-compact/src/New/Seesaw/GoodDatumExists.lean\#L43-L82}{{\leanfont Data.\allowbreak{}exists\_isAdmissibleComb\_reducedVacuum}}}
\expandafter\newcommand\csname LN@projOnFinPartpairTranslates\endcsname{\href{https://github.com/lsernau/grh-lean/blob/72fead102f663d5c7a407964a13dc0cf9d977207/grh-compact/src/New/Seesaw/GoodDatumExists.lean\#L84-L97}{{\leanfont Data.\allowbreak{}projOnFinPart\_pairTranslates}}}
\expandafter\newcommand\csname LN@isAdmissibleCombreducedVacuum\endcsname{\href{https://github.com/lsernau/grh-lean/blob/72fead102f663d5c7a407964a13dc0cf9d977207/grh-compact/src/New/Seesaw/GoodDatumExists.lean\#L101-L113}{{\leanfont Data.\allowbreak{}isAdmissibleComb\_reducedVacuum}}}
\expandafter\newcommand\csname LN@existsisGoodDatumofinvariant\endcsname{\href{https://github.com/lsernau/grh-lean/blob/72fead102f663d5c7a407964a13dc0cf9d977207/grh-compact/src/New/Seesaw/GoodDatumExists.lean\#L121-L136}{{\leanfont Data.\allowbreak{}exists\_isGoodDatum\_of\_invariant}}}
\expandafter\newcommand\csname LN@inclY\endcsname{\href{https://github.com/lsernau/grh-lean/blob/72fead102f663d5c7a407964a13dc0cf9d977207/grh-compact/src/New/Seesaw/GradingsOfY.lean\#L23-L23}{{\leanfont Data.\allowbreak{}inclY}}}
\expandafter\newcommand\csname LN@liftTameParteqofmemQ\endcsname{\href{https://github.com/lsernau/grh-lean/blob/72fead102f663d5c7a407964a13dc0cf9d977207/grh-compact/src/New/Seesaw/GradingsOfY.lean\#L25-L26}{{\leanfont Data.\allowbreak{}liftTamePart\_eq\_of\_mem\_Q}}}
\expandafter\newcommand\csname LN@levelY\endcsname{\href{https://github.com/lsernau/grh-lean/blob/72fead102f663d5c7a407964a13dc0cf9d977207/grh-compact/src/New/Seesaw/GradingsOfY.lean\#L30-L30}{{\leanfont Data.\allowbreak{}levelY}}}
\expandafter\newcommand\csname LN@levelYlechosenLevel\endcsname{\href{https://github.com/lsernau/grh-lean/blob/72fead102f663d5c7a407964a13dc0cf9d977207/grh-compact/src/New/Seesaw/GradingsOfY.lean\#L32-L33}{{\leanfont Data.\allowbreak{}levelY\_le\_chosenLevel}}}
\expandafter\newcommand\csname LN@upsilonY\endcsname{\href{https://github.com/lsernau/grh-lean/blob/72fead102f663d5c7a407964a13dc0cf9d977207/grh-compact/src/New/Seesaw/GradingsOfY.lean\#L37-L38}{{\leanfont Data.\allowbreak{}upsilonY}}}
\expandafter\newcommand\csname LN@thetaKernelY\endcsname{\href{https://github.com/lsernau/grh-lean/blob/72fead102f663d5c7a407964a13dc0cf9d977207/grh-compact/src/New/Seesaw/GradingsOfY.lean\#L42-L43}{{\leanfont Data.\allowbreak{}thetaKernelY}}}
\expandafter\newcommand\csname LN@liftTame\endcsname{\href{https://github.com/lsernau/grh-lean/blob/72fead102f663d5c7a407964a13dc0cf9d977207/grh-compact/src/New/Seesaw/GradingsOfY.lean\#L47-L47}{{\leanfont Data.\allowbreak{}liftTame}}}
\expandafter\newcommand\csname LN@thetaKernelYomegaupsilonY\endcsname{\href{https://github.com/lsernau/grh-lean/blob/72fead102f663d5c7a407964a13dc0cf9d977207/grh-compact/src/New/Seesaw/GradingsOfY.lean\#L52-L62}{{\leanfont Data.\allowbreak{}thetaKernelY\_omega\_upsilonY}}}
\expandafter\newcommand\csname LN@blockNumber\endcsname{\href{https://github.com/lsernau/grh-lean/blob/72fead102f663d5c7a407964a13dc0cf9d977207/grh-compact/src/New/Seesaw/GradingsOfY.lean\#L66-L67}{{\leanfont blockNumber}}}
\expandafter\newcommand\csname LN@blockNumbereqzero\endcsname{\href{https://github.com/lsernau/grh-lean/blob/72fead102f663d5c7a407964a13dc0cf9d977207/grh-compact/src/New/Seesaw/GradingsOfY.lean\#L69-L70}{{\leanfont blockNumber\_eq\_zero}}}
\expandafter\newcommand\csname LN@instNonemptyY\endcsname{\href{https://github.com/lsernau/grh-lean/blob/72fead102f663d5c7a407964a13dc0cf9d977207/grh-compact/src/New/Seesaw/GradingsOfY.lean\#L72-L72}{{\leanfont Data.\allowbreak{}instNonemptyY}}}
\expandafter\newcommand\csname LN@FormY\endcsname{\href{https://github.com/lsernau/grh-lean/blob/72fead102f663d5c7a407964a13dc0cf9d977207/grh-compact/src/New/Seesaw/GradingsOfY.lean\#L76-L76}{{\leanfont Data.\allowbreak{}FormY}}}
\expandafter\newcommand\csname LN@datumY\endcsname{\href{https://github.com/lsernau/grh-lean/blob/72fead102f663d5c7a407964a13dc0cf9d977207/grh-compact/src/New/Seesaw/GradingsOfY.lean\#L80-L81}{{\leanfont Data.\allowbreak{}datumY}}}
\expandafter\newcommand\csname LN@isArrowdatumY\endcsname{\href{https://github.com/lsernau/grh-lean/blob/72fead102f663d5c7a407964a13dc0cf9d977207/grh-compact/src/New/Seesaw/GradingsOfY.lean\#L86-L89}{{\leanfont Data.\allowbreak{}isArrow\_datumY}}}
\expandafter\newcommand\csname LN@liftY\endcsname{\href{https://github.com/lsernau/grh-lean/blob/72fead102f663d5c7a407964a13dc0cf9d977207/grh-compact/src/New/Seesaw/GradingsOfY.lean\#L93-L94}{{\leanfont Data.\allowbreak{}liftY}}}
\expandafter\newcommand\csname LN@lastResLevel\endcsname{\href{https://github.com/lsernau/grh-lean/blob/72fead102f663d5c7a407964a13dc0cf9d977207/grh-compact/src/New/Seesaw/GradingsOfY.lean\#L98-L99}{{\leanfont Data.\allowbreak{}lastResLevel}}}
\expandafter\newcommand\csname LN@liftChainY\endcsname{\href{https://github.com/lsernau/grh-lean/blob/72fead102f663d5c7a407964a13dc0cf9d977207/grh-compact/src/New/Seesaw/GradingsOfY.lean\#L103-L104}{{\leanfont Data.\allowbreak{}liftChainY}}}
\expandafter\newcommand\csname LN@decl\endcsname{\href{https://github.com/lsernau/grh-lean/blob/72fead102f663d5c7a407964a13dc0cf9d977207/grh-compact/src/New/Seesaw/GradingsOfY.lean\#L107-L107}{{\leanfont Data.\allowbreak{}ι}}}
\expandafter\newcommand\csname LN@algFormRepiota\endcsname{\href{https://github.com/lsernau/grh-lean/blob/72fead102f663d5c7a407964a13dc0cf9d977207/grh-compact/src/New/Seesaw/GradingsOfY.lean\#L111-L131}{{\leanfont Data.\allowbreak{}algFormRep\_iota}}}
\expandafter\newcommand\csname LN@lastLevelChar\endcsname{\href{https://github.com/lsernau/grh-lean/blob/72fead102f663d5c7a407964a13dc0cf9d977207/grh-compact/src/New/Seesaw/GradingsOfY.lean\#L135-L137}{{\leanfont Data.\allowbreak{}lastLevelChar}}}
\expandafter\newcommand\csname LN@iotaAtLevel\endcsname{\href{https://github.com/lsernau/grh-lean/blob/72fead102f663d5c7a407964a13dc0cf9d977207/grh-compact/src/New/Seesaw/GradingsOfY.lean\#L141-L145}{{\leanfont Data.\allowbreak{}iotaAtLevel}}}
\expandafter\newcommand\csname LN@groupLift\endcsname{\href{https://github.com/lsernau/grh-lean/blob/72fead102f663d5c7a407964a13dc0cf9d977207/grh-compact/src/New/Seesaw/GradingsOfY.lean\#L149-L150}{{\leanfont Data.\allowbreak{}groupLift}}}
\expandafter\newcommand\csname LN@iotaprincipalRep\endcsname{\href{https://github.com/lsernau/grh-lean/blob/72fead102f663d5c7a407964a13dc0cf9d977207/grh-compact/src/New/Seesaw/GradingsOfY.lean\#L154-L161}{{\leanfont Data.\allowbreak{}iota\_principalRep}}}
\expandafter\newcommand\csname LN@cardparabolicblockNumber\endcsname{\href{https://github.com/lsernau/grh-lean/blob/72fead102f663d5c7a407964a13dc0cf9d977207/grh-compact/src/New/Seesaw/GradingsOfY.lean\#L163-L173}{{\leanfont card\_parabolic\_blockNumber}}}
\expandafter\newcommand\csname LN@kerreductionKeq\endcsname{\href{https://github.com/lsernau/grh-lean/blob/72fead102f663d5c7a407964a13dc0cf9d977207/grh-compact/src/New/Seesaw/GradingsOfY.lean\#L175-L177}{{\leanfont LevelPlaces.\allowbreak{}ker\_reductionK₀\_eq}}}
\expandafter\newcommand\csname LN@cardlevelModulusresRing\endcsname{\href{https://github.com/lsernau/grh-lean/blob/72fead102f663d5c7a407964a13dc0cf9d977207/grh-compact/src/New/Seesaw/GradingsOfY.lean\#L179-L181}{{\leanfont LevelPlaces.\allowbreak{}card\_levelModulus\_resRing}}}
\expandafter\newcommand\csname LN@cardGLlevelModulus\endcsname{\href{https://github.com/lsernau/grh-lean/blob/72fead102f663d5c7a407964a13dc0cf9d977207/grh-compact/src/New/Seesaw/GradingsOfY.lean\#L183-L185}{{\leanfont LevelPlaces.\allowbreak{}card\_GL\_levelModulus}}}
\expandafter\newcommand\csname LN@gradingY\endcsname{\href{https://github.com/lsernau/grh-lean/blob/72fead102f663d5c7a407964a13dc0cf9d977207/grh-compact/src/New/Seesaw/GradingsOfY.lean\#L193-L196}{{\leanfont Data.\allowbreak{}gradingY}}}
\expandafter\newcommand\csname LN@lastGrading\endcsname{\href{https://github.com/lsernau/grh-lean/blob/72fead102f663d5c7a407964a13dc0cf9d977207/grh-compact/src/New/Seesaw/GradingsOfY.lean\#L200-L202}{{\leanfont Data.\allowbreak{}lastGrading}}}
\expandafter\newcommand\csname LN@pairEvalY\endcsname{\href{https://github.com/lsernau/grh-lean/blob/72fead102f663d5c7a407964a13dc0cf9d977207/grh-compact/src/New/Seesaw/GradingsOfY.lean\#L206-L209}{{\leanfont Data.\allowbreak{}pairEvalY}}}
\expandafter\newcommand\csname LN@cutBound\endcsname{\href{https://github.com/lsernau/grh-lean/blob/72fead102f663d5c7a407964a13dc0cf9d977207/grh-compact/src/New/Seesaw/GradingsOfY.lean\#L212-L212}{{\leanfont Data.\allowbreak{}cutBound}}}
\expandafter\newcommand\csname LN@translatedUpsilon\endcsname{\href{https://github.com/lsernau/grh-lean/blob/72fead102f663d5c7a407964a13dc0cf9d977207/grh-compact/src/New/Seesaw/GradingsOfY.lean\#L216-L217}{{\leanfont Data.\allowbreak{}translatedUpsilon}}}
\expandafter\newcommand\csname LN@evaluatedDatum\endcsname{\href{https://github.com/lsernau/grh-lean/blob/72fead102f663d5c7a407964a13dc0cf9d977207/grh-compact/src/New/Seesaw/GradingsOfY.lean\#L221-L222}{{\leanfont Data.\allowbreak{}evaluatedDatum}}}
\expandafter\newcommand\csname LN@coverMultiplicity\endcsname{\href{https://github.com/lsernau/grh-lean/blob/72fead102f663d5c7a407964a13dc0cf9d977207/grh-compact/src/New/Seesaw/IntegralAsCosetSum.lean\#L17-L17}{{\leanfont coverMultiplicity}}}
\expandafter\newcommand\csname LN@coverMultiplicitysmul\endcsname{\href{https://github.com/lsernau/grh-lean/blob/72fead102f663d5c7a407964a13dc0cf9d977207/grh-compact/src/New/Seesaw/IntegralAsCosetSum.lean\#L19-L21}{{\leanfont coverMultiplicity\_smul}}}
\expandafter\newcommand\csname LN@existscutcoverMultiplicitydvd\endcsname{\href{https://github.com/lsernau/grh-lean/blob/72fead102f663d5c7a407964a13dc0cf9d977207/grh-compact/src/New/Seesaw/IntegralAsCosetSum.lean\#L26-L51}{{\leanfont exists\_cut\_coverMultiplicity\_dvd}}}
\expandafter\newcommand\csname LN@measureunioncosets\endcsname{\href{https://github.com/lsernau/grh-lean/blob/72fead102f663d5c7a407964a13dc0cf9d977207/grh-compact/src/New/Seesaw/IntegralAsCosetSum.lean\#L55-L59}{{\leanfont measure\_union\_cosets}}}
\expandafter\newcommand\csname LN@existsintegraleqsumovercut\endcsname{\href{https://github.com/lsernau/grh-lean/blob/72fead102f663d5c7a407964a13dc0cf9d977207/grh-compact/src/New/Seesaw/IntegralAsCosetSum.lean\#L67-L107}{{\leanfont Grading.\allowbreak{}exists\_integral\_eq\_sum\_over\_cut}}}
\expandafter\newcommand\csname LN@ideleLogValaddrevnonpos\endcsname{\href{https://github.com/lsernau/grh-lean/blob/72fead102f663d5c7a407964a13dc0cf9d977207/grh-compact/src/New/Seesaw/LinesStep.lean\#L17-L50}{{\leanfont Grading.\allowbreak{}IsAdaptedAt.\allowbreak{}ideleLogVal\_add\_rev\_nonpos}}}
\expandafter\newcommand\csname LN@step\endcsname{\href{https://github.com/lsernau/grh-lean/blob/72fead102f663d5c7a407964a13dc0cf9d977207/grh-compact/src/New/Seesaw/LinesStep.lean\#L57-L94}{{\leanfont Grading.\allowbreak{}HasNearLines.\allowbreak{}step}}}
\expandafter\newcommand\csname LN@absNormpiGBar\endcsname{\href{https://github.com/lsernau/grh-lean/blob/72fead102f663d5c7a407964a13dc0cf9d977207/grh-compact/src/New/Seesaw/PairTensorData.lean\#L188-L189}{{\leanfont Data.\allowbreak{}absNorm\_piGBar}}}
\expandafter\newcommand\csname LN@vchosenIdele\endcsname{\href{https://github.com/lsernau/grh-lean/blob/72fead102f663d5c7a407964a13dc0cf9d977207/grh-compact/src/New/Seesaw/PairTensorData.lean\#L191-L193}{{\leanfont Data.\allowbreak{}v\_chosenIdele}}}
\expandafter\newcommand\csname LN@valuationshiftconsteqideleAnorm\endcsname{\href{https://github.com/lsernau/grh-lean/blob/72fead102f663d5c7a407964a13dc0cf9d977207/grh-compact/src/New/Seesaw/PairTensorData.lean\#L195-L202}{{\leanfont Data.\allowbreak{}valuation\_shift\_const\_eq\_ideleA\_norm}}}
\expandafter\newcommand\csname LN@spansLatticegradingY\endcsname{\href{https://github.com/lsernau/grh-lean/blob/72fead102f663d5c7a407964a13dc0cf9d977207/grh-compact/src/New/Seesaw/PairTensorData.lean\#L204-L205}{{\leanfont Data.\allowbreak{}spansLattice\_gradingY}}}
\expandafter\newcommand\csname LN@chosenLevelinftrivialOnle\endcsname{\href{https://github.com/lsernau/grh-lean/blob/72fead102f663d5c7a407964a13dc0cf9d977207/grh-compact/src/New/Seesaw/PairTensorData.lean\#L207-L209}{{\leanfont Data.\allowbreak{}chosenLevel\_inf\_trivialOn\_le}}}
\expandafter\newcommand\csname LN@reducedVacuum\endcsname{\href{https://github.com/lsernau/grh-lean/blob/72fead102f663d5c7a407964a13dc0cf9d977207/grh-compact/src/New/Seesaw/PairTensorData.lean\#L215-L216}{{\leanfont Data.\allowbreak{}reducedVacuum}}}
\expandafter\newcommand\csname LN@existsstagecut\endcsname{\href{https://github.com/lsernau/grh-lean/blob/72fead102f663d5c7a407964a13dc0cf9d977207/grh-compact/src/New/Seesaw/SeesawTheorem.lean\#L39-L75}{{\leanfont Data.\allowbreak{}exists\_stage\_cut}}}
\expandafter\newcommand\csname LN@existsisAdmissibleCombunfolding\endcsname{\href{https://github.com/lsernau/grh-lean/blob/72fead102f663d5c7a407964a13dc0cf9d977207/grh-compact/src/New/Seesaw/SeesawTheorem.lean\#L81-L126}{{\leanfont Data.\allowbreak{}exists\_isAdmissibleComb\_unfolding}}}
\expandafter\newcommand\csname LN@seesaw\endcsname{\href{https://github.com/lsernau/grh-lean/blob/72fead102f663d5c7a407964a13dc0cf9d977207/grh-compact/src/New/Seesaw/SeesawTheorem.lean\#L135-L154}{{\leanfont Data.\allowbreak{}seesaw}}}
\expandafter\newcommand\csname LQ@seesaw\endcsname{\leanquoteline{0}{lemma seesaw : ∃ c : ℂˣ, ∀ ϑ, ∃ Ψ, IsGoodDatum ϑ Ψ ∧ ∀ φ,}\leanquoteline{2}{ι ϑ (D.poincare ((c.1 • D.ramifiedColumns.subtype) φ)) = Bμ ϑ Ψ φ}}
\expandafter\newcommand\csname LN@hasNearLineszero\endcsname{\href{https://github.com/lsernau/grh-lean/blob/72fead102f663d5c7a407964a13dc0cf9d977207/grh-compact/src/New/Seesaw/StandardLatticeLevel.lean\#L16-L23}{{\leanfont Grading.\allowbreak{}hasNearLines\_zero}}}
\expandafter\newcommand\csname LN@inclXFinmemKf\endcsname{\href{https://github.com/lsernau/grh-lean/blob/72fead102f663d5c7a407964a13dc0cf9d977207/grh-compact/src/New/Seesaw/StandardLatticeLevel.lean\#L25-L33}{{\leanfont Grading.\allowbreak{}inclXFin\_mem\_Kf}}}
\expandafter\newcommand\csname LN@unfoldStage\endcsname{\href{https://github.com/lsernau/grh-lean/blob/72fead102f663d5c7a407964a13dc0cf9d977207/grh-compact/src/New/Seesaw/UnfoldingStages.lean\#L21-L24}{{\leanfont Data.\allowbreak{}unfoldStage}}}
\expandafter\newcommand\csname LN@unfoldStageupdate\endcsname{\href{https://github.com/lsernau/grh-lean/blob/72fead102f663d5c7a407964a13dc0cf9d977207/grh-compact/src/New/Seesaw/UnfoldingStages.lean\#L26-L27}{{\leanfont Data.\allowbreak{}unfoldStage\_update}}}
\expandafter\newcommand\csname LN@isGenericunfoldStage\endcsname{\href{https://github.com/lsernau/grh-lean/blob/72fead102f663d5c7a407964a13dc0cf9d977207/grh-compact/src/New/Seesaw/UnfoldingStages.lean\#L29-L32}{{\leanfont Data.\allowbreak{}isGeneric\_unfoldStage}}}
\expandafter\newcommand\csname LN@unfoldStageeven\endcsname{\href{https://github.com/lsernau/grh-lean/blob/72fead102f663d5c7a407964a13dc0cf9d977207/grh-compact/src/New/Seesaw/UnfoldingStages.lean\#L38-L60}{{\leanfont Data.\allowbreak{}unfoldStage\_even}}}
\expandafter\newcommand\csname LN@evaluatedDatumout\endcsname{\href{https://github.com/lsernau/grh-lean/blob/72fead102f663d5c7a407964a13dc0cf9d977207/grh-compact/src/New/Seesaw/UnfoldingStages.lean\#L62-L67}{{\leanfont Data.\allowbreak{}evaluatedDatum\_out}}}
\expandafter\newcommand\csname LN@hasNearLinesofunfoldStagenezero\endcsname{\href{https://github.com/lsernau/grh-lean/blob/72fead102f663d5c7a407964a13dc0cf9d977207/grh-compact/src/New/Seesaw/UnfoldingStages.lean\#L72-L112}{{\leanfont Data.\allowbreak{}hasNearLines\_of\_unfoldStage\_ne\_zero}}}
\newcommand{\ConcordanceRows}{\crow{fillingfits}{\href{https://github.com/lsernau/grh-lean/blob/72fead102f663d5c7a407964a13dc0cf9d977207/grh-compact/src/New/Seesaw/BesselGeneric.lean\#L14-L22}{{\leanfont Data.\allowbreak{}filling\_fits}}}{BesselGeneric.lean}{14--22}
\crow{columnCoords}{\href{https://github.com/lsernau/grh-lean/blob/72fead102f663d5c7a407964a13dc0cf9d977207/grh-compact/src/New/Seesaw/BesselGeneric.lean\#L27-L28}{{\leanfont Data.\allowbreak{}columnCoords}}}{BesselGeneric.lean}{27--28}
\crow{goodProjTranslate}{\href{https://github.com/lsernau/grh-lean/blob/72fead102f663d5c7a407964a13dc0cf9d977207/grh-compact/src/New/Seesaw/BesselGeneric.lean\#L32-L33}{{\leanfont Data.\allowbreak{}goodProjTranslate}}}{BesselGeneric.lean}{32--33}
\crow{poincareratmul}{\href{https://github.com/lsernau/grh-lean/blob/72fead102f663d5c7a407964a13dc0cf9d977207/grh-compact/src/New/Seesaw/BesselGeneric.lean\#L40-L41}{{\leanfont Filling.\allowbreak{}Fits.\allowbreak{}poincare\_rat\_mul}}}{BesselGeneric.lean}{40--41}
\crow{existsbesselFinmulmemofnezero}{\href{https://github.com/lsernau/grh-lean/blob/72fead102f663d5c7a407964a13dc0cf9d977207/grh-compact/src/New/Seesaw/BesselGeneric.lean\#L46-L50}{{\leanfont Filling.\allowbreak{}Fits.\allowbreak{}exists\_besselFin\_mul\_mem\_of\_ne\_zero}}}{BesselGeneric.lean}{46--50}
\crow{poincarebesselnormalform}{\href{https://github.com/lsernau/grh-lean/blob/72fead102f663d5c7a407964a13dc0cf9d977207/grh-compact/src/New/Seesaw/BesselGeneric.lean\#L58-L95}{{\leanfont Filling.\allowbreak{}Fits.\allowbreak{}poincare\_bessel\_normal\_form}}}{BesselGeneric.lean}{58--95}
\crow{isGenericpoincare}{\href{https://github.com/lsernau/grh-lean/blob/72fead102f663d5c7a407964a13dc0cf9d977207/grh-compact/src/New/Seesaw/BesselGeneric.lean\#L101-L123}{{\leanfont Filling.\allowbreak{}Fits.\allowbreak{}isGeneric\_poincare}}}{BesselGeneric.lean}{101--123}
\crow{flagOrder}{\href{https://github.com/lsernau/grh-lean/blob/72fead102f663d5c7a407964a13dc0cf9d977207/grh-compact/src/New/Seesaw/EqualLinesOrbit.lean\#L18-L18}{{\leanfont flagOrder}}}{EqualLinesOrbit.lean}{18--18}
\crow{mulapplystringofblockTriangular}{\href{https://github.com/lsernau/grh-lean/blob/72fead102f663d5c7a407964a13dc0cf9d977207/grh-compact/src/New/Seesaw/EqualLinesOrbit.lean\#L24-L27}{{\leanfont mul\_apply\_string\_of\_blockTriangular}}}{EqualLinesOrbit.lean}{24--27}
\crow{IsAdaptedBasis}{\href{https://github.com/lsernau/grh-lean/blob/72fead102f663d5c7a407964a13dc0cf9d977207/grh-compact/src/New/Seesaw/EqualLinesOrbit.lean\#L31-L34}{{\leanfont Grading.\allowbreak{}IsAdaptedBasis}}}{EqualLinesOrbit.lean}{31--34}
\crow{IsAdaptedAt}{\href{https://github.com/lsernau/grh-lean/blob/72fead102f663d5c7a407964a13dc0cf9d977207/grh-compact/src/New/Seesaw/EqualLinesOrbit.lean\#L38-L38}{{\leanfont Grading.\allowbreak{}IsAdaptedAt}}}{EqualLinesOrbit.lean}{38--38}
\crow{HasNearLines}{\href{https://github.com/lsernau/grh-lean/blob/72fead102f663d5c7a407964a13dc0cf9d977207/grh-compact/src/New/Seesaw/EqualLinesOrbit.lean\#L42-L44}{{\leanfont Grading.\allowbreak{}HasNearLines}}}{EqualLinesOrbit.lean}{42--44}
\crow{existsisAdaptedAt}{\href{https://github.com/lsernau/grh-lean/blob/72fead102f663d5c7a407964a13dc0cf9d977207/grh-compact/src/New/Seesaw/EqualLinesOrbit.lean\#L46-L49}{{\leanfont Grading.\allowbreak{}exists\_isAdaptedAt}}}{EqualLinesOrbit.lean}{46--49}
\crow{gramstringentry}{\href{https://github.com/lsernau/grh-lean/blob/72fead102f663d5c7a407964a13dc0cf9d977207/grh-compact/src/New/Seesaw/EqualLinesOrbit.lean\#L54-L60}{{\leanfont Grading.\allowbreak{}IsAdaptedAt.\allowbreak{}gram\_string\_entry}}}{EqualLinesOrbit.lean}{54--60}
\crow{inv}{\href{https://github.com/lsernau/grh-lean/blob/72fead102f663d5c7a407964a13dc0cf9d977207/grh-compact/src/New/Seesaw/EqualLinesOrbit.lean\#L62-L64}{{\leanfont Grading.\allowbreak{}IsAdaptedBasis.\allowbreak{}inv}}}{EqualLinesOrbit.lean}{62--64}
\crow{ideleLogValle}{\href{https://github.com/lsernau/grh-lean/blob/72fead102f663d5c7a407964a13dc0cf9d977207/grh-compact/src/New/Seesaw/EqualLinesOrbit.lean\#L68-L73}{{\leanfont Grading.\allowbreak{}IsAdaptedBasis.\allowbreak{}ideleLogVal\_le}}}{EqualLinesOrbit.lean}{68--73}
\crow{ideleLogValeq}{\href{https://github.com/lsernau/grh-lean/blob/72fead102f663d5c7a407964a13dc0cf9d977207/grh-compact/src/New/Seesaw/EqualLinesOrbit.lean\#L75-L77}{{\leanfont Grading.\allowbreak{}IsAdaptedBasis.\allowbreak{}ideleLogVal\_eq}}}{EqualLinesOrbit.lean}{75--77}
\crow{existsmembesselFin}{\href{https://github.com/lsernau/grh-lean/blob/72fead102f663d5c7a407964a13dc0cf9d977207/grh-compact/src/New/Seesaw/EqualLinesOrbit.lean\#L84-L106}{{\leanfont Grading.\allowbreak{}IsAdaptedAt.\allowbreak{}exists\_mem\_besselFin}}}{EqualLinesOrbit.lean}{84--106}
\crow{isAdmissibleEvalpairEvalY}{\href{https://github.com/lsernau/grh-lean/blob/72fead102f663d5c7a407964a13dc0cf9d977207/grh-compact/src/New/Seesaw/GoodDatumExists.lean\#L21-L38}{{\leanfont Data.\allowbreak{}isAdmissibleEval\_pairEvalY}}}{GoodDatumExists.lean}{21--38}
\crow{existsisAdmissibleCombreducedVacuum}{\href{https://github.com/lsernau/grh-lean/blob/72fead102f663d5c7a407964a13dc0cf9d977207/grh-compact/src/New/Seesaw/GoodDatumExists.lean\#L43-L82}{{\leanfont Data.\allowbreak{}exists\_isAdmissibleComb\_reducedVacuum}}}{GoodDatumExists.lean}{43--82}
\crow{projOnFinPartpairTranslates}{\href{https://github.com/lsernau/grh-lean/blob/72fead102f663d5c7a407964a13dc0cf9d977207/grh-compact/src/New/Seesaw/GoodDatumExists.lean\#L84-L97}{{\leanfont Data.\allowbreak{}projOnFinPart\_pairTranslates}}}{GoodDatumExists.lean}{84--97}
\crow{isAdmissibleCombreducedVacuum}{\href{https://github.com/lsernau/grh-lean/blob/72fead102f663d5c7a407964a13dc0cf9d977207/grh-compact/src/New/Seesaw/GoodDatumExists.lean\#L101-L113}{{\leanfont Data.\allowbreak{}isAdmissibleComb\_reducedVacuum}}}{GoodDatumExists.lean}{101--113}
\crow{existsisGoodDatumofinvariant}{\href{https://github.com/lsernau/grh-lean/blob/72fead102f663d5c7a407964a13dc0cf9d977207/grh-compact/src/New/Seesaw/GoodDatumExists.lean\#L121-L136}{{\leanfont Data.\allowbreak{}exists\_isGoodDatum\_of\_invariant}}}{GoodDatumExists.lean}{121--136}
\crow{inclY}{\href{https://github.com/lsernau/grh-lean/blob/72fead102f663d5c7a407964a13dc0cf9d977207/grh-compact/src/New/Seesaw/GradingsOfY.lean\#L23-L23}{{\leanfont Data.\allowbreak{}inclY}}}{GradingsOfY.lean}{23--23}
\crow{liftTameParteqofmemQ}{\href{https://github.com/lsernau/grh-lean/blob/72fead102f663d5c7a407964a13dc0cf9d977207/grh-compact/src/New/Seesaw/GradingsOfY.lean\#L25-L26}{{\leanfont Data.\allowbreak{}liftTamePart\_eq\_of\_mem\_Q}}}{GradingsOfY.lean}{25--26}
\crow{levelY}{\href{https://github.com/lsernau/grh-lean/blob/72fead102f663d5c7a407964a13dc0cf9d977207/grh-compact/src/New/Seesaw/GradingsOfY.lean\#L30-L30}{{\leanfont Data.\allowbreak{}levelY}}}{GradingsOfY.lean}{30--30}
\crow{levelYlechosenLevel}{\href{https://github.com/lsernau/grh-lean/blob/72fead102f663d5c7a407964a13dc0cf9d977207/grh-compact/src/New/Seesaw/GradingsOfY.lean\#L32-L33}{{\leanfont Data.\allowbreak{}levelY\_le\_chosenLevel}}}{GradingsOfY.lean}{32--33}
\crow{upsilonY}{\href{https://github.com/lsernau/grh-lean/blob/72fead102f663d5c7a407964a13dc0cf9d977207/grh-compact/src/New/Seesaw/GradingsOfY.lean\#L37-L38}{{\leanfont Data.\allowbreak{}upsilonY}}}{GradingsOfY.lean}{37--38}
\crow{thetaKernelY}{\href{https://github.com/lsernau/grh-lean/blob/72fead102f663d5c7a407964a13dc0cf9d977207/grh-compact/src/New/Seesaw/GradingsOfY.lean\#L42-L43}{{\leanfont Data.\allowbreak{}thetaKernelY}}}{GradingsOfY.lean}{42--43}
\crow{liftTame}{\href{https://github.com/lsernau/grh-lean/blob/72fead102f663d5c7a407964a13dc0cf9d977207/grh-compact/src/New/Seesaw/GradingsOfY.lean\#L47-L47}{{\leanfont Data.\allowbreak{}liftTame}}}{GradingsOfY.lean}{47--47}
\crow{thetaKernelYomegaupsilonY}{\href{https://github.com/lsernau/grh-lean/blob/72fead102f663d5c7a407964a13dc0cf9d977207/grh-compact/src/New/Seesaw/GradingsOfY.lean\#L52-L62}{{\leanfont Data.\allowbreak{}thetaKernelY\_omega\_upsilonY}}}{GradingsOfY.lean}{52--62}
\crow{blockNumber}{\href{https://github.com/lsernau/grh-lean/blob/72fead102f663d5c7a407964a13dc0cf9d977207/grh-compact/src/New/Seesaw/GradingsOfY.lean\#L66-L67}{{\leanfont blockNumber}}}{GradingsOfY.lean}{66--67}
\crow{blockNumbereqzero}{\href{https://github.com/lsernau/grh-lean/blob/72fead102f663d5c7a407964a13dc0cf9d977207/grh-compact/src/New/Seesaw/GradingsOfY.lean\#L69-L70}{{\leanfont blockNumber\_eq\_zero}}}{GradingsOfY.lean}{69--70}
\crow{instNonemptyY}{\href{https://github.com/lsernau/grh-lean/blob/72fead102f663d5c7a407964a13dc0cf9d977207/grh-compact/src/New/Seesaw/GradingsOfY.lean\#L72-L72}{{\leanfont Data.\allowbreak{}instNonemptyY}}}{GradingsOfY.lean}{72--72}
\crow{FormY}{\href{https://github.com/lsernau/grh-lean/blob/72fead102f663d5c7a407964a13dc0cf9d977207/grh-compact/src/New/Seesaw/GradingsOfY.lean\#L76-L76}{{\leanfont Data.\allowbreak{}FormY}}}{GradingsOfY.lean}{76--76}
\crow{datumY}{\href{https://github.com/lsernau/grh-lean/blob/72fead102f663d5c7a407964a13dc0cf9d977207/grh-compact/src/New/Seesaw/GradingsOfY.lean\#L80-L81}{{\leanfont Data.\allowbreak{}datumY}}}{GradingsOfY.lean}{80--81}
\crow{isArrowdatumY}{\href{https://github.com/lsernau/grh-lean/blob/72fead102f663d5c7a407964a13dc0cf9d977207/grh-compact/src/New/Seesaw/GradingsOfY.lean\#L86-L89}{{\leanfont Data.\allowbreak{}isArrow\_datumY}}}{GradingsOfY.lean}{86--89}
\crow{liftY}{\href{https://github.com/lsernau/grh-lean/blob/72fead102f663d5c7a407964a13dc0cf9d977207/grh-compact/src/New/Seesaw/GradingsOfY.lean\#L93-L94}{{\leanfont Data.\allowbreak{}liftY}}}{GradingsOfY.lean}{93--94}
\crow{lastResLevel}{\href{https://github.com/lsernau/grh-lean/blob/72fead102f663d5c7a407964a13dc0cf9d977207/grh-compact/src/New/Seesaw/GradingsOfY.lean\#L98-L99}{{\leanfont Data.\allowbreak{}lastResLevel}}}{GradingsOfY.lean}{98--99}
\crow{liftChainY}{\href{https://github.com/lsernau/grh-lean/blob/72fead102f663d5c7a407964a13dc0cf9d977207/grh-compact/src/New/Seesaw/GradingsOfY.lean\#L103-L104}{{\leanfont Data.\allowbreak{}liftChainY}}}{GradingsOfY.lean}{103--104}
\crow{decl}{\href{https://github.com/lsernau/grh-lean/blob/72fead102f663d5c7a407964a13dc0cf9d977207/grh-compact/src/New/Seesaw/GradingsOfY.lean\#L107-L107}{{\leanfont Data.\allowbreak{}ι}}}{GradingsOfY.lean}{107--107}
\crow{algFormRepiota}{\href{https://github.com/lsernau/grh-lean/blob/72fead102f663d5c7a407964a13dc0cf9d977207/grh-compact/src/New/Seesaw/GradingsOfY.lean\#L111-L131}{{\leanfont Data.\allowbreak{}algFormRep\_iota}}}{GradingsOfY.lean}{111--131}
\crow{lastLevelChar}{\href{https://github.com/lsernau/grh-lean/blob/72fead102f663d5c7a407964a13dc0cf9d977207/grh-compact/src/New/Seesaw/GradingsOfY.lean\#L135-L137}{{\leanfont Data.\allowbreak{}lastLevelChar}}}{GradingsOfY.lean}{135--137}
\crow{iotaAtLevel}{\href{https://github.com/lsernau/grh-lean/blob/72fead102f663d5c7a407964a13dc0cf9d977207/grh-compact/src/New/Seesaw/GradingsOfY.lean\#L141-L145}{{\leanfont Data.\allowbreak{}iotaAtLevel}}}{GradingsOfY.lean}{141--145}
\crow{groupLift}{\href{https://github.com/lsernau/grh-lean/blob/72fead102f663d5c7a407964a13dc0cf9d977207/grh-compact/src/New/Seesaw/GradingsOfY.lean\#L149-L150}{{\leanfont Data.\allowbreak{}groupLift}}}{GradingsOfY.lean}{149--150}
\crow{iotaprincipalRep}{\href{https://github.com/lsernau/grh-lean/blob/72fead102f663d5c7a407964a13dc0cf9d977207/grh-compact/src/New/Seesaw/GradingsOfY.lean\#L154-L161}{{\leanfont Data.\allowbreak{}iota\_principalRep}}}{GradingsOfY.lean}{154--161}
\crow{cardparabolicblockNumber}{\href{https://github.com/lsernau/grh-lean/blob/72fead102f663d5c7a407964a13dc0cf9d977207/grh-compact/src/New/Seesaw/GradingsOfY.lean\#L163-L173}{{\leanfont card\_parabolic\_blockNumber}}}{GradingsOfY.lean}{163--173}
\crow{kerreductionKeq}{\href{https://github.com/lsernau/grh-lean/blob/72fead102f663d5c7a407964a13dc0cf9d977207/grh-compact/src/New/Seesaw/GradingsOfY.lean\#L175-L177}{{\leanfont LevelPlaces.\allowbreak{}ker\_reductionK₀\_eq}}}{GradingsOfY.lean}{175--177}
\crow{cardlevelModulusresRing}{\href{https://github.com/lsernau/grh-lean/blob/72fead102f663d5c7a407964a13dc0cf9d977207/grh-compact/src/New/Seesaw/GradingsOfY.lean\#L179-L181}{{\leanfont LevelPlaces.\allowbreak{}card\_levelModulus\_resRing}}}{GradingsOfY.lean}{179--181}
\crow{cardGLlevelModulus}{\href{https://github.com/lsernau/grh-lean/blob/72fead102f663d5c7a407964a13dc0cf9d977207/grh-compact/src/New/Seesaw/GradingsOfY.lean\#L183-L185}{{\leanfont LevelPlaces.\allowbreak{}card\_GL\_levelModulus}}}{GradingsOfY.lean}{183--185}
\crow{gradingY}{\href{https://github.com/lsernau/grh-lean/blob/72fead102f663d5c7a407964a13dc0cf9d977207/grh-compact/src/New/Seesaw/GradingsOfY.lean\#L193-L196}{{\leanfont Data.\allowbreak{}gradingY}}}{GradingsOfY.lean}{193--196}
\crow{lastGrading}{\href{https://github.com/lsernau/grh-lean/blob/72fead102f663d5c7a407964a13dc0cf9d977207/grh-compact/src/New/Seesaw/GradingsOfY.lean\#L200-L202}{{\leanfont Data.\allowbreak{}lastGrading}}}{GradingsOfY.lean}{200--202}
\crow{pairEvalY}{\href{https://github.com/lsernau/grh-lean/blob/72fead102f663d5c7a407964a13dc0cf9d977207/grh-compact/src/New/Seesaw/GradingsOfY.lean\#L206-L209}{{\leanfont Data.\allowbreak{}pairEvalY}}}{GradingsOfY.lean}{206--209}
\crow{cutBound}{\href{https://github.com/lsernau/grh-lean/blob/72fead102f663d5c7a407964a13dc0cf9d977207/grh-compact/src/New/Seesaw/GradingsOfY.lean\#L212-L212}{{\leanfont Data.\allowbreak{}cutBound}}}{GradingsOfY.lean}{212--212}
\crow{translatedUpsilon}{\href{https://github.com/lsernau/grh-lean/blob/72fead102f663d5c7a407964a13dc0cf9d977207/grh-compact/src/New/Seesaw/GradingsOfY.lean\#L216-L217}{{\leanfont Data.\allowbreak{}translatedUpsilon}}}{GradingsOfY.lean}{216--217}
\crow{evaluatedDatum}{\href{https://github.com/lsernau/grh-lean/blob/72fead102f663d5c7a407964a13dc0cf9d977207/grh-compact/src/New/Seesaw/GradingsOfY.lean\#L221-L222}{{\leanfont Data.\allowbreak{}evaluatedDatum}}}{GradingsOfY.lean}{221--222}
\crow{coverMultiplicity}{\href{https://github.com/lsernau/grh-lean/blob/72fead102f663d5c7a407964a13dc0cf9d977207/grh-compact/src/New/Seesaw/IntegralAsCosetSum.lean\#L17-L17}{{\leanfont coverMultiplicity}}}{IntegralAsCosetSum.lean}{17--17}
\crow{coverMultiplicitysmul}{\href{https://github.com/lsernau/grh-lean/blob/72fead102f663d5c7a407964a13dc0cf9d977207/grh-compact/src/New/Seesaw/IntegralAsCosetSum.lean\#L19-L21}{{\leanfont coverMultiplicity\_smul}}}{IntegralAsCosetSum.lean}{19--21}
\crow{existscutcoverMultiplicitydvd}{\href{https://github.com/lsernau/grh-lean/blob/72fead102f663d5c7a407964a13dc0cf9d977207/grh-compact/src/New/Seesaw/IntegralAsCosetSum.lean\#L26-L51}{{\leanfont exists\_cut\_coverMultiplicity\_dvd}}}{IntegralAsCosetSum.lean}{26--51}
\crow{measureunioncosets}{\href{https://github.com/lsernau/grh-lean/blob/72fead102f663d5c7a407964a13dc0cf9d977207/grh-compact/src/New/Seesaw/IntegralAsCosetSum.lean\#L55-L59}{{\leanfont measure\_union\_cosets}}}{IntegralAsCosetSum.lean}{55--59}
\crow{existsintegraleqsumovercut}{\href{https://github.com/lsernau/grh-lean/blob/72fead102f663d5c7a407964a13dc0cf9d977207/grh-compact/src/New/Seesaw/IntegralAsCosetSum.lean\#L67-L107}{{\leanfont Grading.\allowbreak{}exists\_integral\_eq\_sum\_over\_cut}}}{IntegralAsCosetSum.lean}{67--107}
\crow{ideleLogValaddrevnonpos}{\href{https://github.com/lsernau/grh-lean/blob/72fead102f663d5c7a407964a13dc0cf9d977207/grh-compact/src/New/Seesaw/LinesStep.lean\#L17-L50}{{\leanfont Grading.\allowbreak{}IsAdaptedAt.\allowbreak{}ideleLogVal\_add\_rev\_nonpos}}}{LinesStep.lean}{17--50}
\crow{step}{\href{https://github.com/lsernau/grh-lean/blob/72fead102f663d5c7a407964a13dc0cf9d977207/grh-compact/src/New/Seesaw/LinesStep.lean\#L57-L94}{{\leanfont Grading.\allowbreak{}HasNearLines.\allowbreak{}step}}}{LinesStep.lean}{57--94}
\crow{absNormpiGBar}{\href{https://github.com/lsernau/grh-lean/blob/72fead102f663d5c7a407964a13dc0cf9d977207/grh-compact/src/New/Seesaw/PairTensorData.lean\#L188-L189}{{\leanfont Data.\allowbreak{}absNorm\_piGBar}}}{PairTensorData.lean}{188--189}
\crow{vchosenIdele}{\href{https://github.com/lsernau/grh-lean/blob/72fead102f663d5c7a407964a13dc0cf9d977207/grh-compact/src/New/Seesaw/PairTensorData.lean\#L191-L193}{{\leanfont Data.\allowbreak{}v\_chosenIdele}}}{PairTensorData.lean}{191--193}
\crow{valuationshiftconsteqideleAnorm}{\href{https://github.com/lsernau/grh-lean/blob/72fead102f663d5c7a407964a13dc0cf9d977207/grh-compact/src/New/Seesaw/PairTensorData.lean\#L195-L202}{{\leanfont Data.\allowbreak{}valuation\_shift\_const\_eq\_ideleA\_norm}}}{PairTensorData.lean}{195--202}
\crow{spansLatticegradingY}{\href{https://github.com/lsernau/grh-lean/blob/72fead102f663d5c7a407964a13dc0cf9d977207/grh-compact/src/New/Seesaw/PairTensorData.lean\#L204-L205}{{\leanfont Data.\allowbreak{}spansLattice\_gradingY}}}{PairTensorData.lean}{204--205}
\crow{chosenLevelinftrivialOnle}{\href{https://github.com/lsernau/grh-lean/blob/72fead102f663d5c7a407964a13dc0cf9d977207/grh-compact/src/New/Seesaw/PairTensorData.lean\#L207-L209}{{\leanfont Data.\allowbreak{}chosenLevel\_inf\_trivialOn\_le}}}{PairTensorData.lean}{207--209}
\crow{reducedVacuum}{\href{https://github.com/lsernau/grh-lean/blob/72fead102f663d5c7a407964a13dc0cf9d977207/grh-compact/src/New/Seesaw/PairTensorData.lean\#L215-L216}{{\leanfont Data.\allowbreak{}reducedVacuum}}}{PairTensorData.lean}{215--216}
\crow{existsstagecut}{\href{https://github.com/lsernau/grh-lean/blob/72fead102f663d5c7a407964a13dc0cf9d977207/grh-compact/src/New/Seesaw/SeesawTheorem.lean\#L39-L75}{{\leanfont Data.\allowbreak{}exists\_stage\_cut}}}{SeesawTheorem.lean}{39--75}
\crow{existsisAdmissibleCombunfolding}{\href{https://github.com/lsernau/grh-lean/blob/72fead102f663d5c7a407964a13dc0cf9d977207/grh-compact/src/New/Seesaw/SeesawTheorem.lean\#L81-L126}{{\leanfont Data.\allowbreak{}exists\_isAdmissibleComb\_unfolding}}}{SeesawTheorem.lean}{81--126}
\crow{seesaw}{\href{https://github.com/lsernau/grh-lean/blob/72fead102f663d5c7a407964a13dc0cf9d977207/grh-compact/src/New/Seesaw/SeesawTheorem.lean\#L135-L154}{{\leanfont Data.\allowbreak{}seesaw}}}{SeesawTheorem.lean}{135--154}
\crow{hasNearLineszero}{\href{https://github.com/lsernau/grh-lean/blob/72fead102f663d5c7a407964a13dc0cf9d977207/grh-compact/src/New/Seesaw/StandardLatticeLevel.lean\#L16-L23}{{\leanfont Grading.\allowbreak{}hasNearLines\_zero}}}{StandardLatticeLevel.lean}{16--23}
\crow{inclXFinmemKf}{\href{https://github.com/lsernau/grh-lean/blob/72fead102f663d5c7a407964a13dc0cf9d977207/grh-compact/src/New/Seesaw/StandardLatticeLevel.lean\#L25-L33}{{\leanfont Grading.\allowbreak{}inclXFin\_mem\_Kf}}}{StandardLatticeLevel.lean}{25--33}
\crow{unfoldStage}{\href{https://github.com/lsernau/grh-lean/blob/72fead102f663d5c7a407964a13dc0cf9d977207/grh-compact/src/New/Seesaw/UnfoldingStages.lean\#L21-L24}{{\leanfont Data.\allowbreak{}unfoldStage}}}{UnfoldingStages.lean}{21--24}
\crow{unfoldStageupdate}{\href{https://github.com/lsernau/grh-lean/blob/72fead102f663d5c7a407964a13dc0cf9d977207/grh-compact/src/New/Seesaw/UnfoldingStages.lean\#L26-L27}{{\leanfont Data.\allowbreak{}unfoldStage\_update}}}{UnfoldingStages.lean}{26--27}
\crow{isGenericunfoldStage}{\href{https://github.com/lsernau/grh-lean/blob/72fead102f663d5c7a407964a13dc0cf9d977207/grh-compact/src/New/Seesaw/UnfoldingStages.lean\#L29-L32}{{\leanfont Data.\allowbreak{}isGeneric\_unfoldStage}}}{UnfoldingStages.lean}{29--32}
\crow{unfoldStageeven}{\href{https://github.com/lsernau/grh-lean/blob/72fead102f663d5c7a407964a13dc0cf9d977207/grh-compact/src/New/Seesaw/UnfoldingStages.lean\#L38-L60}{{\leanfont Data.\allowbreak{}unfoldStage\_even}}}{UnfoldingStages.lean}{38--60}
\crow{evaluatedDatumout}{\href{https://github.com/lsernau/grh-lean/blob/72fead102f663d5c7a407964a13dc0cf9d977207/grh-compact/src/New/Seesaw/UnfoldingStages.lean\#L62-L67}{{\leanfont Data.\allowbreak{}evaluatedDatum\_out}}}{UnfoldingStages.lean}{62--67}
\crow{hasNearLinesofunfoldStagenezero}{\href{https://github.com/lsernau/grh-lean/blob/72fead102f663d5c7a407964a13dc0cf9d977207/grh-compact/src/New/Seesaw/UnfoldingStages.lean\#L72-L112}{{\leanfont Data.\allowbreak{}hasNearLines\_of\_unfoldStage\_ne\_zero}}}{UnfoldingStages.lean}{72--112}}
% the notations of the defined functions: one slot for each explicit argument of the Lean declaration
\newcommand{\colco}[1]{\operatorname{cc}_{#1}}
\newcommand{\gpt}[2]{\mathrm{pr}_{#1}(#2)}
\newcommand{\flag}[2]{\operatorname{fl}_{#1}(#2)}
\newcommand{\adb}[4]{(#3,#4)\allowbreak\text{ is adapted to }\allowbreak#2\text{ for }#1}
\newcommand{\adat}[2]{\text{adapted at }\allowbreak#2\text{ for }#1}
\newcommand{\near}[6]{#6\allowbreak\text{ has near lines for }\allowbreak(#1;#2,#3;#4,#5)}
\newcommand{\inclY}[1]{\operatorname{in}_{D,#1}}
\newcommand{\levelY}[2]{\mathcal{J}_{#2}(D,#1)}
\newcommand{\ups}[2]{\Upsilon_{#2}(D,#1)}
\newcommand{\kerY}[2]{\mathcal{K}_{#2}(D,#1)}
\newcommand{\ltame}[2]{\psi^{\uparrow}_{#2}(D,#1)}
\newcommand{\blk}[2]{\operatorname{bl}_{#1}(#2)}
\newcommand{\FormY}[1]{\mathcal{F}_{#1}(D)}
\newcommand{\datY}[2]{\mathfrak{d}_{#2}(D,#1)}
\newcommand{\liftY}[3]{\mathcal{L}_{#2}(D,#1)#3}
\newcommand{\lastRes}[1]{\mathsf{P}_{\mathrm{last}}(D,#1)}
\newcommand{\chainY}[3]{\mathcal{L}^{(#2)}(D,#1)#3}
\newcommand{\iot}[1]{\iota_{D,#1}}
\newcommand{\lastChar}[2]{\operatorname{lch}_{D,#1}(#2)}
\newcommand{\iotL}[2]{\iota^{\mathrm{lev}}_{D,#1}#2}
\newcommand{\groupLift}[2]{\operatorname{lift}_{D,#1}(#2)}
\newcommand{\pairEvalY}[1]{\operatorname{ev}_{D,#1}}
\newcommand{\tUps}[3]{\Upsilon_{#3}(D,#1;#2)}
\newcommand{\evD}[3]{\operatorname{ed}_{#3}(D,#1;#2)}
\newcommand{\cov}[4]{\operatorname{mult}_{#1,#2}(#3;#4)}
\newcommand{\redvac}[2]{\mathrm{rv}_{#1}(#2)}
\newcommand{\UnitMembers}{312}
\newcommand{\UnitInNote}{80}
\newcommand{\UnitProofs}{83}
\newcommand{\UnitSteps}{503}
\newcommand{\UnitAllSteps}{550}


\linespread{1.08}                       % 11 pt on 14.7 pt
\setlength{\parindent}{0pt}
\setlength{\parskip}{5.5pt plus 1.5pt minus 0.5pt}
\setlength{\abovedisplayskip}{7.5pt plus 2pt minus 3pt}
\setlength{\belowdisplayskip}{7.5pt plus 2pt minus 3pt}
\setlength{\abovedisplayshortskip}{3pt plus 2pt}
\setlength{\belowdisplayshortskip}{6pt plus 2pt minus 2pt}
\emergencystretch=1.5em
\clubpenalty=10000 \widowpenalty=10000 \displaywidowpenalty=10000
\relpenalty=7000 \binoppenalty=9000
\raggedbottom
\pagestyle{fancy}\fancyhf{}\renewcommand{\headrulewidth}{0pt}
\fancyfoot[C]{\sffamily\fontsize{7.6}{9}\selectfont\color{inkfaint}\thepage\,/\,\pageref*{LastPage}}
\setlist{topsep=2pt,partopsep=0pt,parsep=0pt,itemsep=4pt,leftmargin=2.1em,labelsep=0.55em}

\newcommand{\Q}{\mathbb{Q}}
\newcommand{\Z}{\mathbb{Z}}
\newcommand{\F}{\mathbb{F}}
\newcommand{\len}{\operatorname{len}}
\newcommand{\cok}{\operatorname{cok}}
\newcommand{\End}{\operatorname{End}}
\newcommand{\Hom}{\operatorname{Hom}}
\newcommand{\Gal}{\operatorname{Gal}}
\newcommand{\diag}{\operatorname{diag}}
\newcommand{\kw}[1]{\emph{#1}}
\newcommand{\m}[1]{\mbox{$#1$}}        % a short formula that must not be broken across lines
\newcommand{\head}[1]{\par\addvspace{17pt plus 4pt minus 2pt}\Needspace{6\baselineskip}%
  \noindent{\fontsize{13.2}{16}\selectfont\bfseries\color{ink}#1}\par\nopagebreak\vspace{-1pt}}
\newcommand{\sect}[2]{\refstepcounter{section}\label{#2}\par\addvspace{19pt plus 4pt minus 2pt}\Needspace{7\baselineskip}%
  \noindent\llap{{\fontsize{13.2}{16}\selectfont\bfseries\color{inkfaint}\thesection}\hspace{9pt}}%
  {\fontsize{13.2}{16}\selectfont\bfseries\color{ink}#1}\par\nopagebreak\vspace{-1pt}}
\newcommand{\subhead}[1]{\par\addvspace{9pt plus 2pt}\Needspace{4\baselineskip}\noindent\textbf{#1}\enspace\ignorespaces}
% where a statement lives in the Lean archive
\newcommand{\leantag}{{\sffamily\bfseries\fontsize{6.3}{8}\selectfont\color{leancol}\addfontfeatures{LetterSpace=7}LEAN}}
\newcommand{\sep}{\unskip\nobreak\hspace{5pt}{\color{inkfaint}\textperiodcentered}\hspace{5pt plus 2pt}\ignorespaces}
\newcommand{\leanline}[1]{\par\nopagebreak\vspace{-2.5pt}{\fontsize{8.9}{12.4}\selectfont\color{inksoft}\raggedright
  \noindent\llap{\leantag\hspace{9pt}}#1\par}\vspace{1pt}}
% a displayed statement in a tinted box: #1 colour, #2 text
\newcommand{\statement}[2]{\par\addvspace{7pt}\begingroup\setlength{\fboxsep}{8.5pt}\setlength{\fboxrule}{0.6pt}%
  \noindent\fcolorbox{#1!30!white}{#1!3!white}{\parbox{\dimexpr\linewidth-2\fboxsep-2\fboxrule\relax}{#2}}%
  \endgroup\par\addvspace{7pt}}
\newcommand{\nt}[2]{\par\noindent\hangindent=1.1em\hangafter=1 #1\enspace{\color{inksoft}#2}\par\vspace{1pt}}
\newcommand{\legdot}[1]{\tikz[baseline=-0.62ex]\fill[#1] (0,0) circle (2.5pt);\hspace{3pt}}
\newcommand{\crow}[4]{\csname CS@#1\endcsname & #2 & \textcolor{inksoft}{#3} & \textcolor{inksoft}{#4}\\}
\newcommand{\leanquoteline}[2]{\par\noindent\hangindent=\dimexpr#1\fontdimen2\font+2.2em\relax\hangafter=1
  \hspace*{#1\fontdimen2\font}#2}
\newcommand{\quoted}[1]{\par\vspace{3pt}{\ttfamily\fontsize{7.5}{10.4}\selectfont\color{leancol}\raggedright
  \setlength{\parskip}{0pt}\leanquote{#1}\par}\vspace{4pt}}
\newcommand{\bx}[3]{{\sffamily\bfseries\fontsize{7.3}{9.2}\selectfont\color{#1!82!black}#2}\\[1.5pt]#3}
\newcommand{\cs}[2]{\expandafter\newcommand\csname CS@#1\endcsname{#2}}

% theorems
\newtheoremstyle{plainx}{7pt plus 2pt}{7pt plus 2pt}{\itshape}{}{\bfseries}{.}{0.5em}{}
\newtheoremstyle{uprightx}{7pt plus 2pt}{7pt plus 2pt}{}{}{\bfseries}{.}{0.5em}{}
\theoremstyle{plainx}
\newtheorem{theorem}{Theorem}[section]
\newtheorem{lemma}[theorem]{Lemma}
\newtheorem{proposition}[theorem]{Proposition}
\newtheorem{corollary}[theorem]{Corollary}
\theoremstyle{uprightx}
\newtheorem{definition}[theorem]{Definition}
\newtheorem{setting}[theorem]{Setting}          % hypotheses shared by several statements; each of them cites it by number
\newtheorem{remark}[theorem]{Remark}
\newtheorem{example}[theorem]{Example}
\renewcommand{\qedsymbol}{{\setlength{\fboxsep}{0pt}\setlength{\fboxrule}{0.5pt}\fbox{\rule{0pt}{0.46em}\hspace{0.46em}}}}
% what holds throughout the note: the ambient objects only (implicit and instance arguments), never a hypothesis
\newenvironment{conventions}{\par\addvspace{5pt}\noindent\textbf{Conventions.}\enspace\ignorespaces}{\par\addvspace{3pt}}
\makeatletter
\renewenvironment{proof}[1][Proof]{\par\pushQED{\qed}\normalfont\topsep0pt\partopsep0pt
  \trivlist\item[\hskip\labelsep\itshape #1.]\ignorespaces}{\popQED\endtrivlist\@endpefalse}
\makeatother
% the two boxes of Inputs carry their own numbers: \inputslabel{number}{label}
\makeatletter\newcommand{\inputslabel}[2]{\def\@currentlabel{#1}\label{#2}}\makeatother

% where each declaration of the concordance is treated: \cs{key}{place}, the place naming a \label or a lettered item
\cs{seesaw}{Theorem A}
\cs{existsisAdmissibleCombunfolding}{Theorem B}
\cs{existsstagecut}{Theorem C}
\cs{inclY}{Definition~\ref{d:index}}
\cs{blockNumber}{Definition~\ref{d:index}}
\cs{instNonemptyY}{Definition~\ref{d:index}}
\cs{gradingY}{Definition~\ref{d:gradings}}
\cs{lastGrading}{Definition~\ref{d:gradings}}
\cs{levelY}{Definition~\ref{d:lift}}
\cs{upsilonY}{Definition~\ref{d:lift}}
\cs{thetaKernelY}{Definition~\ref{d:lift}}
\cs{liftTame}{Definition~\ref{d:lift}}
\cs{FormY}{Definition~\ref{d:lift}}
\cs{datumY}{Definition~\ref{d:lift}}
\cs{liftY}{Definition~\ref{d:lift}}
\cs{liftChainY}{Definition~\ref{d:chain}}
\cs{decl}{Definition~\ref{d:chain}}
\cs{lastResLevel}{Definition~\ref{d:lastlevel}}
\cs{lastLevelChar}{Definition~\ref{d:lastlevel}}
\cs{iotaAtLevel}{Definition~\ref{d:lastlevel}}
\cs{groupLift}{Definition~\ref{d:lastlevel}}
\cs{columnCoords}{Definition~\ref{d:poin}}
\cs{goodProjTranslate}{Definition~\ref{d:poin}}
\cs{reducedVacuum}{Definition~\ref{d:redvac}}
\cs{unfoldStage}{Definition~\ref{d:stage}}
\cs{pairEvalY}{Definition~\ref{d:eval}}
\cs{cutBound}{Definition~\ref{d:eval}}
\cs{translatedUpsilon}{Definition~\ref{d:eval}}
\cs{evaluatedDatum}{Definition~\ref{d:eval}}
\cs{flagOrder}{Definition~\ref{d:adapted}}
\cs{IsAdaptedBasis}{Definition~\ref{d:adapted}}
\cs{IsAdaptedAt}{Definition~\ref{d:adapted}}
\cs{HasNearLines}{Definition~\ref{d:near}}
\cs{coverMultiplicity}{Definition~\ref{d:mult}}
\cs{liftTameParteqofmemQ}{Lemma~\ref{l:tame}}
\cs{levelYlechosenLevel}{Lemma~\ref{l:tame}}
\cs{isArrowdatumY}{Lemma~\ref{l:tame}}
\cs{thetaKernelYomegaupsilonY}{Lemma~\ref{l:upsilon}}
\cs{blockNumbereqzero}{Lemma~\ref{l:blocks}}
\cs{cardparabolicblockNumber}{Lemma~\ref{l:blocks}}
\cs{kerreductionKeq}{Lemma~\ref{l:blocks}}
\cs{cardlevelModulusresRing}{Lemma~\ref{l:blocks}}
\cs{cardGLlevelModulus}{Lemma~\ref{l:blocks}}
\cs{algFormRepiota}{Proposition~\ref{p:equiv}}
\cs{iotaprincipalRep}{Corollary~\ref{c:groupG}}
\cs{fillingfits}{Proposition~\ref{p:fits}}
\cs{poincareratmul}{Lemma~\ref{l:poinbasic}}
\cs{existsbesselFinmulmemofnezero}{Lemma~\ref{l:poinbasic}}
\cs{poincarebesselnormalform}{Proposition~\ref{p:normal}}
\cs{isGenericpoincare}{Proposition~\ref{p:genericpoin}}
\cs{unfoldStageupdate}{Lemma~\ref{l:stage}}
\cs{isGenericunfoldStage}{Lemma~\ref{l:stage}}
\cs{evaluatedDatumout}{Lemma~\ref{l:stage}}
\cs{unfoldStageeven}{Proposition~\ref{p:even}}
\cs{mulapplystringofblockTriangular}{Lemma~\ref{l:adapted}}
\cs{existsisAdaptedAt}{Lemma~\ref{l:adapted}}
\cs{gramstringentry}{Lemma~\ref{l:adapted}}
\cs{inv}{Lemma~\ref{l:adapted}}
\cs{ideleLogValle}{Lemma~\ref{l:adapted}}
\cs{ideleLogValeq}{Lemma~\ref{l:adapted}}
\cs{existsmembesselFin}{Proposition~\ref{p:orbit}}
\cs{hasNearLineszero}{Lemma~\ref{l:nearzero}}
\cs{inclXFinmemKf}{Lemma~\ref{l:inclK}}
\cs{ideleLogValaddrevnonpos}{Proposition~\ref{p:rev}}
\cs{step}{Proposition~\ref{p:step}}
\cs{absNormpiGBar}{Lemma~\ref{l:vals}}
\cs{vchosenIdele}{Lemma~\ref{l:vals}}
\cs{valuationshiftconsteqideleAnorm}{Lemma~\ref{l:vals}}
\cs{spansLatticegradingY}{Lemma~\ref{l:vals}}
\cs{chosenLevelinftrivialOnle}{Lemma~\ref{l:vals}}
\cs{hasNearLinesofunfoldStagenezero}{Proposition~\ref{p:nearstage}}
\cs{coverMultiplicitysmul}{Lemma~\ref{l:mult}}
\cs{measureunioncosets}{Lemma~\ref{l:mult}}
\cs{existscutcoverMultiplicitydvd}{Proposition~\ref{p:cut}}
\cs{existsintegraleqsumovercut}{Proposition~\ref{p:sumcut}}
\cs{isAdmissibleEvalpairEvalY}{Proposition~\ref{p:evalpair}}
\cs{existsisAdmissibleCombreducedVacuum}{Proposition~\ref{p:redvac}}
\cs{projOnFinPartpairTranslates}{Lemma~\ref{l:proj}}
\cs{isAdmissibleCombreducedVacuum}{Lemma~\ref{l:redvaccomb}}
\cs{existsisGoodDatumofinvariant}{Proposition~\ref{p:good}}

% ------------------------------------------------------------------ macros of this note that are not notations of
% definitions of the note (those are made by the build, in data.tex, from the spec)
\newcommand{\CC}{\mathbb{C}}
\newcommand{\RR}{\mathbb{R}}
\newcommand{\NN}{\mathbb{N}}
\newcommand{\ZZ}{\mathbb{Z}}
\newcommand{\Ad}{\mathbb{A}}
\newcommand{\Af}{\mathbb{A}_f}
\newcommand{\Oh}{\hat{\mathfrak{o}}}
\newcommand{\cj}[1]{\overline{#1}}
\newcommand{\ord}{\operatorname{ord}}
\newcommand{\Nm}{\operatorname{N}}
\newcommand{\GL}{\operatorname{GL}}
\newcommand{\sgn}{\operatorname{sgn}}
\newcommand{\gr}{\operatorname{gr}}
\newcommand{\tensor}{\otimes}
% objects of the data from other units and from the other notes of this unit; every one shows the data
\newcommand{\ED}[1]{E(#1)}                         % the CM field
\newcommand{\ND}[1]{N(#1)}                         % the level prime
\newcommand{\piN}[1]{\pi_N(#1)}                    % an integer of E(D) of norm N(D)
\newcommand{\piGb}[1]{\bar\pi_g(#1)}               % the conjugate of \pi_g in E(D)
\newcommand{\gD}[1]{g(#1)}                         % g_0 g_1
\newcommand{\fD}[1]{f(#1)}                         % the number of infinite places
\newcommand{\tD}[1]{t(#1)}                         % 5r
\newcommand{\rhoD}[1]{\mathrm{hm}(#1)}                 % m/2 - 1
\newcommand{\nD}[2]{n_{#2}(#1)}                    % k(D) + j(t(D)+1)
\newcommand{\XI}[2]{\underline{X}_{#2}(#1)}        % index set of X_j
\newcommand{\YI}[2]{\underline{Y}_{#2}(#1)}        % index set of Y_j
\newcommand{\Xs}[2]{X_{#2}(#1)}                    % the standard definite space
\newcommand{\Ys}[2]{Y_{#2}(#1)}                    % the hermitian space with frame
\newcommand{\GD}[1]{\mathsf{G}(#1)}                % the finite group
\newcommand{\ezero}[1]{e_0(#1)}                    % the good idempotent
\newcommand{\PiD}[1]{\Pi(#1)}                      % the places dividing \pi_N(D)
\newcommand{\lvn}[1]{\mathrm{n}(#1)}               % the product of the norms of the places of Q
\newcommand{\Qs}[1]{Q(#1)}                         % the set of places of a tame character
\newcommand{\ideA}[1]{\mathfrak{a}(#1)}            % the idele \alpha / conj(chosen idele)
\newcommand{\lsc}[2]{u_{#2}(#1)}                   % the lattice scale
\newcommand{\muc}[3]{\mu_{#1,#3}(#2)}              % \mu^0_{D,j} \vartheta
\newcommand{\psic}[3]{\psi_{#1,#3}(#2)}            % \omega_D^j \vartheta^{\mathrm{t}}
\newcommand{\Col}[1]{\mathrm{Col}(#1)}             % the columns
\newcommand{\Colram}[1]{\mathrm{Col}^{\mathrm{ram}}(#1)}
\newcommand{\Stg}[2]{\mathrm{St}_{#2}(#1)}         % smooth algebraic forms of stage j
\newcommand{\PoinD}[1]{\mathrm{P}_{#1}}            % the Poincare series
\newcommand{\Bmu}[3]{B^{#3}_{#1,#2}}               % the definite chain, restricted
\newcommand{\Tr}[1]{\mathrm{Tr}(#1)}               % tuples of translations
\newcommand{\Tadm}[2]{\mathrm{Tr}_{#2}(#1)}        % admissible translations for q
\newcommand{\tdat}[3]{\mathrm{v}_{#1}(#2)_{#3}}    % the translated vacuum of stage j
\newcommand{\Kst}[2]{K^{\mathrm{col}}_{#2}(#1)}                   % the level of stage j on X_j
\newcommand{\Kch}[2]{K^{\mathrm{ch}}_{#1}(#2)}     % principal level at the chosen places
\newcommand{\Kaw}[2]{K^{(#1)}(#2)}                 % integral at the places not dividing n
\newcommand{\Kz}[1]{K_0(#1)}                       % the stabiliser of the lattice
\newcommand{\Kw}[2]{K_{#2}(#1)}                    % elements of K_0 supported at w and its conjugate
\newcommand{\kerX}[3]{k_{#1,#2,#3}}                % the theta kernel of the definite lift j+1
\newcommand{\chn}[4]{\mathrm{C}^{(#4)}_{#1,#2}(#3)}        % chain of j definite lifts with data
\newcommand{\chnn}[4]{\mathrm{C}^{[#4]}_{#1,#2}(#3)}       % chain of pairs of lifts with paired data
\newcommand{\chcomb}[4]{\mathrm{C}^{(#4)}_{#1,#2}[#3]}     % the combination of chains
\newcommand{\resY}[1]{\operatorname{res}_{#1}}     % restriction to Y_m
\newcommand{\ptr}[3]{\tau_{#1}(#2)_{#3}}           % the pair translates
\newcommand{\Adm}[6]{\operatorname{Adm}(#1,#2;#3,#4;#5,#6)}  % admissible combinations
\newcommand{\Oaw}[1]{O^{(#1)}}                     % complex numbers integral away from n
\newcommand{\gY}[1]{\mathcal{Y}_{#1}(D)}           % the grading of Y_j(D)
\newcommand{\gYlast}{\mathcal{Y}_{\mathrm{last}}(D)}
\newcommand{\cutB}{\mathrm{cut}(D)}                % g(D)^{f(D) n_m(D)^2}
\newcommand{\unf}[6]{\mathfrak{P}^{(#5)}_{D,#1,#2}[#3,#4;#6]} % the stage function
\newcommand{\Bes}[1]{B(#1)}                        % the Bessel group of a grading
\newcommand{\Besf}[1]{B_f(#1)}
\newcommand{\Besr}[1]{\Gamma_B(#1)}                % its rational points
\newcommand{\Uni}[1]{U(#1)}                        % the unipotent subgroup of a grading
\newcommand{\incl}[1]{i_{#1}}                      % U(X) -> U(V)
\newcommand{\inclf}[1]{i_{#1,f}}
\newcommand{\cst}[1]{\mathfrak{c}(#1)}             % the constant of a grading
\newcommand{\lv}[2]{\ell_{#2}(#1)}                 % minus the order at w of an idele
\newcommand{\coef}[2]{\langle #1,#2\rangle}        % the coefficient of an algebraic form at a functional
\newcommand{\cD}[1]{c_{#1}}                         % the constant of Theorem A
\newcommand{\vpi}[2]{\varpi(#1,#2)}                % the generator of the different of E(P,\mu) (Chapter III)
\newcommand{\cgD}[1]{\hat{\pi}_g(#1)}               % the idele of \bar\pi_g(D) (unit Trace)
\newcommand{\piId}[1]{\hat{\pi}_N(#1)}              % the principal idele of \pi_N(D) (unit Trace)
\newcommand{\zw}[2]{z_{#2}(#1)}                     % the central translation at a place of \Pi(D)
\newcommand{\XIo}[2]{\underline{X}^{\circ}_{#2}(#1)} % X_j(D) without its point
\newcommand{\transum}[5]{\mathsf{T}_{#4,#3}(#1,#2,#5)} % the translate sum, as in the note on the Gauss-sum operator
\newcommand{\Tm}[2]{T_{#1,#2}}                      % the string map of two gradings, as in the note on two lifts
\newcommand{\pW}[1]{\omega^{\mathrm{pr}}_{#1}}      % the twisted pair action, as in the note on two lifts

\begin{document}
\color{ink}

{\raggedright
{\sffamily\bfseries\fontsize{7.2}{9}\selectfont\color{leancol}\addfontfeatures{LetterSpace=8}READBACK NOTE 21\par}
\vspace{5pt}
{\fontsize{21.5}{25.5}\selectfont\bfseries The seesaw identity for the chain of theta lifts\par}
\vspace{5pt}
{\fontsize{12.4}{16}\selectfont\itshape\color{inksoft} One unit of the Lean artifact \ArtifactName, read back into
mathematics\par}}
\vspace{8pt}
{\color{rulegrey}\hrule height 0.6pt}
\vspace{7pt}
\thispagestyle{empty}

\textbf{Summary.} The rows of Theorem~A of the write-up are Poincar\'e series on an indefinite unitary group of
rank~$r$, and the forms whose size and divisibility must be controlled are their images under a chain of $m$ theta
lifts through hermitian spaces of growing rank, the last of them definite. This note proves the identity that makes
that control possible: up to one non-zero constant, the chain of lifts of the Poincar\'e series of a ramified column
is a finite combination of chains of $m$ theta lifts of the column itself between \emph{definite} spaces, restricted
at the end, and the coefficients of the combination have denominators only at the conductor of the twist. The proof
unfolds the Poincar\'e series and then one lift at a time against Bessel periods; each unfolded period is a finite
sum over cosets whose weights are reciprocals of numbers dividing a power of the norm of the conductor.

\statement{leancol}{\fontsize{9.6}{13.2}\selectfont\textbf{What this note is.} The artifact \ArtifactName{} is a
body of Lean code whose final theorems state the Riemann hypothesis and its analogue for Dirichlet $L$-functions.
Its authors report that it compiles and that Lean's kernel accepts these theorems from the three standard axioms;
the record of checks of the write-up says who replayed what. That is a fact about a formal text: no
mathematician has read the argument through, and the result is a \emph{claim}. This note is one of four notes of
the unit \emph{Seesaw} (Chapter~V, Theorem~A: short forms) and treats \UnitInNote{} of its \UnitMembers{}
declarations: the indefinite chain of lifts, the stage functions of its unfolding, the lattice estimates and the
identity itself. The proofs of this note's own declarations are written out; the statements taken from other units, from the
other three notes of the unit and from the background library are marked as \emph{Inputs}, with their owners.
Every statement was compared with its Lean text, argument by argument. The Lean text read is \PinText; Appendix~A
is the concordance of \ConcDecls{} declarations.}

\sect{The statement}{s:intro}

\begin{setting}[the data]\label{s:data}
Let $P=(N_0,d)$ be parameters and $D$ admissible data for~$P$ (Chapter~III of the write-up). The components of the
record $D$ are written with bare letters: the scales $x$, $x_1$, $b$, the exponent~$\mu$, the number of lifts~$m$,
the rank~$r$, the thick set, the primes \m{g_0\ne g_1} and the integer $\pi_g$ of \m{K=\mathbb{Q}(\sqrt{-d})} with
\m{\pi_g\cj{\pi_g}=g_0g_1}. Everything else attached to $D$ is written as a function of~$D$:
\begin{enumerate}[label=(\alph*)]
\item from Chapter~III: the CM field $\ED{D}$ with its CM type, the set of its infinite places, and its Weil data
  (an element $\mathbf{j}$ with \m{\cj{\mathbf{j}}=-\mathbf{j}}, a finite idele~$\alpha$ and the splitting
  character~$\upsilon$); the number $\fD{D}$ of infinite places; the product \m{\gD{D}=g_0g_1}; \m{\tD{D}=5r};
  the ranks $\nD{D}{j}$, with \m{\nD{D}{0}=r+1-m};
\item from the note \emph{The characters, the level prime and the index sets of the data} of this unit: the number
  $\rhoD{D}$, for which \m{m=2\rhoD{D}+2}; the integer $\piGb{D}$ of $\ED{D}$, the conjugate of~$\pi_g$, whose prime
  divisors are the \emph{chosen places}; the level prime $\ND{D}$, which splits completely in $\ED{D}$ and is prime
  to $\gD{D}$, and an integer $\piN{D}$ of $\ED{D}$ with \m{\piN{D}\,\cj{\piN{D}}=\ND{D}}; \emph{tame characters}
  $\vartheta$ of~$D$, each with its finite set $\Qs{\vartheta}$ of places and its tame part $\vartheta^{\mathrm{t}}$, the
  number $\lvn{\vartheta}=\prod_{w\in\Qs{\vartheta}}\Nm w$, and the \emph{twists} of~$D$, the tame characters with
  \m{\vartheta^{\mathrm{t}}_w\ne1} for all \m{w\in\Qs{\vartheta}}; the characters $\muc{D}{\vartheta}{j}$ of the norm-one
  adeles and their tame parts $\psic{D}{\vartheta}{j}$; the index sets $\XI{D}{j}$ and $\YI{D}{j}$, where $\YI{D}{j}$
  is $\YI{D}{0}$ followed by $j$ blocks of $\tD{D}$ indices; the finite group
  $\GD{D}=\prod\GL_{\YI{D}{0}}(\mathbb{F}_w)$, the product over the chosen places~$w$, and the element $\ezero{D}$
  of its group algebra;
\item from the unit \emph{Trace} (Chapter~VII): for every $j$ a hermitian space with frame $\Ys{D}{j}$ on the index
  set $\YI{D}{j}$ (Definition~\ref{d:space}), of rank \m{r+j\,\tD{D}}, which for \m{j\ge m} is the standard definite
  space; the space of forms of principal level on $\Ys{D}{0}$ with its representation of $\GD{D}$
  (Definition~\ref{d:forms});
\item from the units \emph{Columns} and \emph{Exchange} (Chapters~V and~VII): the set $\PiD{D}$ of the places of
  $\ED{D}$ that divide $\piN{D}$; the standard definite spaces $\Xs{D}{j}$ on $\XI{D}{j}$, the columns
  $\Col{D}$ and the ramified columns $\Colram{D}$, the definite chains and the maps $\Bmu{D}{\vartheta}{\Psi}$
  (Definitions~\ref{d:columns} and~\ref{d:good});
\item from the unit \emph{Filling} (Chapter~VI): the Poincar\'e series \m{\PoinD{D}} (Definition~\ref{d:poin}).
\end{enumerate}
\end{setting}

The left side of the identity is \emph{indefinite}: $\iot{\vartheta}$ is the composite of $m$ theta lifts
$\Ys{D}{0}\to\Ys{D}{1}\to\dots\to\Ys{D}{m}$, the $j$-th formed with the character of the $j$-th power
of~$\vartheta$, applied to forms of principal level on the indefinite space $\Ys{D}{0}$
(Definition~\ref{d:chain}). The right side is \emph{definite}: $\Bmu{D}{\vartheta}{\Psi}$ is a finite combination,
with coefficients~$\Psi$, of composites of $m$ theta lifts $\Xs{D}{0}\to\dots\to\Xs{D}{m}$ of algebraic modular
forms, followed by the restriction from $\Xs{D}{m}$ to $\Ys{D}{m}$ (Definition~\ref{d:good}).

\statement{arith}{\phantomsection\label{t:A}\textbf{Theorem A (the seesaw identity).}\enspace In
Setting~\ref{s:data} there is one non-zero complex number~$\cD{D}$, depending only on~$D$, such that for every twist $\vartheta$ of~$D$ there is a
finitely supported function $\Psi\colon\Tr{D}\to\CC$ which is a good datum for $\vartheta$
(Definition~\ref{d:good}) and satisfies
\[\iot{\vartheta}\bigl(\PoinD{D}(\cD{D}\,\varphi)\bigr)\;=\;\Bmu{D}{\vartheta}{\Psi}(\varphi)
\qquad\text{for every ramified column }\varphi\in\Colram{D}.\]
The constant $\cD{D}$ is chosen before~$\vartheta$, and $\Psi$ before~$\varphi$: one scalar for the datum~$D$, one good
datum for each twist, and the identity for all ramified columns at once; no condition is put on the column.}
\leanline{\lean{seesaw}}

The proof has two engines. The first is an induction over the lifts. Its statement is abstract on purpose: the
map $S$ below will be specialised, in the proof of Theorem~A, to a tensor of finite theta data.

\statement{arith}{\phantomsection\label{t:B}\textbf{Theorem B (unfolding).}\enspace In Setting~\ref{s:data} let
$j$ be a natural number with \m{j\le 2\rhoD{D}+2}, let $\mathcal{Y}$ be a grading of $\Ys{D}{j}$ with norm $\ND{D}$,
with a finite set $X$ and a string of length~$l$, where \m{l=2\rhoD{D}+1-j} (and \m{l=0} for \m{j=2\rhoD{D}+2}),
and let $\mathsf{a}$ be a complex function of the tuples \m{\xi=(\xi_i)_{i\in\NN}} of real points \m{\xi_i\in G_\infty(\Ys{D}{i})}.
There is a complex number $C$ such that for every tame character $\vartheta$ of~$D$, every complex vector space $M$
and every map $S$ from the tuples \m{q=(q_i)_{i\in\NN}} of cosets $q_i\in G_f(\Ys{D}{i})/\levelY{\vartheta}{i}$ to~$M$
there is a vector \m{\mathsf{w}\in M} such that:
\begin{enumerate}[label=(\roman*)]
\item $\mathsf{w}\in\Adm{v_S}{\mathrm{all}}{\lvn{\vartheta}}{1}{\Qs{\vartheta}=\emptyset}{\cutB^{\,j}}$, where $v_S$ sends a
  tuple \m{y=(y_i)} of finite points $y_i\in\Kaw{\lvn{\vartheta}}{\Ys{D}{i}}$, integral at the places not dividing
  $\lvn{\vartheta}$, to $S\bigl((y_i\,\levelY{\vartheta}{i})_i\bigr)$, and ``all'' is the set of all such tuples;
\item \m{\Omega(\mathsf{w})=\mathsf{w}} for every tuple \m{\mathbf{b}=(\mathbf{b}_i)} with \m{\mathbf{b}_i\in\Kst{D}{i}} and every linear map
  $\Omega\colon M\to M$ such that \m{\Omega(S(q))=S(q')} for all~$q$, where
  $q'_i=\inclf{\gY{i}}(\mathbf{b}_i)\,q_i$ for \m{i<j} and \m{q'_i=q_i} for \m{i\ge j};
\item for every form $z$ on $\Ys{D}{j}$ of the conjugate type \m{\cj{\FormY{j}}}, every ramified column
  \m{\varphi\in\Colram{D}} and every linear form $L\colon M\to\CC$: \emph{if} for all tuples $\xi$ of real points
  and $y$ of finite points with \m{y_{2\rhoD{D}+2}=1} and \m{y_0\in\Kch{D}{\Ys{D}{0}}} the function
  $F=\unf{\vartheta}{\varphi}{\xi}{y}{j}{\mathcal{Y}}\cdot z$, in case it is integrable on
  $\Besr{\mathcal{Y}}\backslash G(\Ys{D}{j})$, satisfies
  \[\int_{\Besr{\mathcal{Y}}\backslash G(\Ys{D}{j})}F\;=\;\mathsf{a}(\xi)\cdot L\bigl(S((y_i\,\levelY{\vartheta}{i})_i)\bigr),\]
  \emph{then}
  \[\int_{[G(\Ys{D}{j})]}z\cdot\chainY{\vartheta}{j}{\bigl(\PoinD{D}\varphi\bigr)}\;=\;C\cdot L(\mathsf{w}).\]
\end{enumerate}
The notation $\Adm{\cdot}{\cdot}{\cdot}{\cdot}{\cdot}{\cdot}$ for admissible combinations is recalled in
Definition~\ref{d:adm}, the number $\cutB$ in Definition~\ref{d:eval}.}
\leanline{\lean{existsisAdmissibleCombunfolding}}

The second engine cuts one period into a finite sum; the induction of Theorem~B uses it once at every lift.

\statement{arith}{\phantomsection\label{t:C}\textbf{Theorem C (finite cuts).}\enspace In Setting~\ref{s:data} let
$j$ be a natural number with \m{j\le 2\rhoD{D}+1}, and let \m{\mathcal{Y}=\gY{j}} be the grading of $\Ys{D}{j}$
(Definition~\ref{d:gradings}). There is a level $A$ of $\Ys{D}{j}$ with
$\inclf{\mathcal{Y}}(\Kst{D}{j})\subseteq A$ such that for every tame character $\vartheta$ of $D$ there is a finite
set $T$ of cosets of $\levelY{\vartheta}{j}$ in the finite adelic group of $\Ys{D}{j}$ with these properties:
\begin{enumerate}[label=(\roman*)]
\item for every \m{p\in\Besf{\mathcal{Y}}\cap A} and every coset $t$: \m{t\in T} if and only if \m{p\,t\in T};
\item if the coset of a finite point $h$ lies in~$T$, then $h$ is integral at every place not dividing $\lvn{\vartheta}$, and for
  \m{j=0} it lies in the principal level \m{\Kch{D}{\Ys{D}{0}}} at the chosen places;
\item let $h$ be a finite adelic point whose coset lies in $T$, or for which there are a ramified column
  \m{\varphi\in\Colram{D}}, tuples $\xi$, $y$ of real and of finite points with \m{y_0\in\Kch{D}{\Ys{D}{0}}} and a
  real point $g_\infty$ with $\unf{\vartheta}{\varphi}{\xi}{y}{j}{\mathcal{Y}}(g_\infty,h)\ne0$; then the
  multiplicity $\cov{\Besf{\mathcal{Y}}}{A}{T}{h\,\levelY{\vartheta}{j}}$ of Definition~\ref{d:mult} divides
  a power of $\lvn{\vartheta}$.
\end{enumerate}
Here $\Kst{D}{j}$ is the level of stage~$j$ (Definition~\ref{d:columns}), $\levelY{\vartheta}{j}$ the level of
Definition~\ref{d:lift}, and $\unf{\vartheta}{\varphi}{\xi}{y}{j}{\mathcal{Y}}$ the stage function of
Definition~\ref{d:stage}.}
\leanline{\lean{existsstagecut}}

\sect{The setting}{s:setting}
\begin{conventions}\phantomsection\label{c:conventions}
$E$ is a CM field with a CM type, with complex conjugation \m{z\mapsto\cj z}; where theta functions occur, Weil
data for~$E$ are fixed, and where integral Witt theory is used $E$ is stably unramified. The letters $I$, $I'$, $X$,
$X'$ and the letter $H$ for an abstract group denote ambient objects: finite index sets and a group. A hermitian
space with frame $V$ over~$E$ has an index set and signs; a grading $\mathcal{Y}$ of $V$ has a norm~$\mathfrak{N}$, a finite
set~$X$ and a string of length~$l$; a second hermitian space with frame is $V'$, with a grading $\mathcal{Y}'$ whose
string has length~$l'$. $P$ denotes parameters and $D$ admissible data for~$P$; a tame character of $D$ is
$\vartheta$. A commutative ring is $R$; index sets of blocks are $R_0$ and $T_0$. Measures are Haar measures or
the measures they induce; a measurable group carries a left invariant measure. Matrices act on column
vectors; $g^{*}$ is the conjugate transpose of a matrix~$g$. All statements about filling data concern filling data
$\mathsf{S}$ of a shape, with a parameter, weights and a central element, as in Setting~\ref{s:filling}.
\end{conventions}

\subhead{Hermitian spaces, their groups and levels.} The following vocabulary belongs to the background library
of the artifact; it is recalled here in the form in which this note uses it.

\begin{definition}[spaces with frame; recalled from the background library]\label{d:space}
Let $\Ad$ be the ring of adeles of~$E$, $\Af$ the ring of finite adeles and $\Oh$ its subring of integral elements.
A \emph{hermitian space with frame} $V$ over $E$ consists of a finite index set~$I$, signs $\sgn_V\colon I\to\{\pm1\}$,
an invertible hermitian matrix \m{J_V\in\GL_I(E)}, the Gram matrix of the basis of~$V$, and a \emph{frame}
\m{x_V\in\GL_I(\Ad)} with \m{x_V^{*}J_Vx_V} equal to the diagonal matrix of the signs at every infinite place and to
the identity at the finite places. Its finite part is \m{x_{V,f}\in\GL_I(\Af)}, and \m{x_{V,f}\Oh^{I}} is the
\emph{lattice} of~$V$; it is unimodular. The groups of $V$ are written in the coordinates of the frame:
\m{G_f(V)=\{g\in\GL_I(\Af):g^{*}g=1\}}, the product $G_\infty(V)$ over the infinite places of the unitary groups of
the signs, and \m{G(V)=G_\infty(V)\times G_f(V)}; a finite point $g$ is identified with \m{(1,g)}. In
\emph{rational coordinates} a point $g$ is the matrix \m{x_Vgx_V^{-1}}, which preserves~$J_V$, and the
\emph{rational points} $\Gamma(V)$ are the $g$ for which it has entries in~$E$. The quotient \m{[G(V)]=\Gamma(V)\backslash G(V)}
carries the invariant measure fixed in the library, and for a subgroup $\Delta$ of $\Gamma(V)$ the quotient
\m{\Delta\backslash G(V)} carries the measure induced by the same Haar measure. $V$ is \emph{standard} if all signs
are $1$ and \m{x_V=1}. The subgroup \m{\Kz{V}=G_f(V)\cap\GL_I(\Oh)} is the stabiliser of the lattice, and a
\emph{level} of $V$ is an open subgroup of $G_f(V)$ contained in~$\Kz{V}$. For a finite place $w$, $\Kw{V}{w}$ is the
subgroup of the elements of $\Kz{V}$ that are $1$ at all places other than $w$ and~$\cj w$. For a natural number
$n$, \m{\Kaw{n}{V}} is the subgroup of the finite points whose component at $w$ is integral with integral inverse
for every finite place $w$ not dividing~$n$. For a finite idele $\beta$ and a finite place~$w$, $\ord_w\beta$ is the
order of the component of $\beta$ at~$w$.
\end{definition}

\begin{definition}[gradings and Bessel groups; recalled from the background library]\label{d:grading}
Let $V$ be a hermitian space with frame, with index set~$I$, let \m{\mathfrak{N}\in E} with \m{\mathfrak{N}\ne0} and \m{\cj{\mathfrak{N}}=\mathfrak{N}}, let
$X$ be a finite set and \m{l\ge0} an integer. A \emph{grading} $\mathcal{Y}$ of~$V$ as \emph{$X$ plus a string of
length~$l$ with norm~$\mathfrak{N}$} consists of a bijection $e_{\mathcal{Y}}\colon X\sqcup\{0,\dots,l-1\}\to I$ and a
constant \m{\cst{\mathcal{Y}}\in E} with \m{\cst{\mathcal{Y}}\ne0} and $\cj{\cst{\mathcal{Y}}}=\cst{\mathcal{Y}}$,
such that in the basis indexed through~$e_{\mathcal{Y}}$ the Gram matrix $J_V$ is the identity on~$X$, pairs the
string vectors $z_a$ and $z_{a'}$ to~$1$ if \m{a+a'=l-1} and \m{a\ne a'}, to \m{-\mathfrak{N}} if \m{a=a'} and \m{2a=l-1}, and to
$0$ otherwise, and pairs $X$ with the string to~$0$. The \emph{grade} of a basis vector is $0$ on~$X$ and
\m{\gr(z_a)=l-1-2a} on the string; $z_{l-1-a}$ is the \emph{reverse} of~$z_a$. With matrices in rational
coordinates:
\begin{itemize}
\item the \emph{unipotent subgroup} $\Uni{\mathcal{Y}}$ of $G(V)$ consists of the points $u$ for which \m{u-1}
  strictly raises grades;
\item $\incl{\mathcal{Y}}\colon G(X)\to G(V)$ is the inclusion of the group of the standard definite space on~$X$,
  acting on the span of $X$ and fixing the string vectors, and $\inclf{\mathcal{Y}}$ is its finite part;
\item the \emph{Bessel group} $\Bes{\mathcal{Y}}$ consists of the points that fix every string vector of grade
  \m{\ge0} up to vectors of strictly higher grade; it contains $\Uni{\mathcal{Y}}$ and the image of
  $\incl{\mathcal{Y}}$. Its finite part is $\Besf{\mathcal{Y}}=\{g\in G_f(V):(1,g)\in\Bes{\mathcal{Y}}\}$, its
  rational points are $\Besr{\mathcal{Y}}=\Gamma(V)\cap\Bes{\mathcal{Y}}$, and \m{[\Bes{\mathcal{Y}}]} is the
  quotient of $\Bes{\mathcal{Y}}$ by~$\Besr{\mathcal{Y}}$, with the finite measure induced by a Haar measure.
\end{itemize}
A function $F\colon G(V)\to\CC$ is \emph{generic for~$\mathcal{Y}$} if it is continuous, left invariant under
$\Besr{\mathcal{Y}}$, and for every \m{u\in\Uni{\mathcal{Y}}} such that \m{u-1} raises grades by at least~$2$ and
every \m{g\in G(V)} satisfies $F(ug)=\psi_E\bigl(\tfrac{\mathbf{j}}2\operatorname{tr}(\mathrm{sh}_{\mathcal{Y}}\,u)\bigr)F(g)$;
here $\psi_E$ is the standard additive character of \m{\Ad/E} and $\mathrm{sh}_{\mathcal{Y}}$ is the rational
matrix that kills $X$ and sends $z_a$ to a multiple of~$z_{a+1}$ (and the last string vector to~$0$), the multiple
being $\cst{\mathcal{Y}}$ times \m{-\mathfrak{N}} if $z_a$ is the middle vector, times \m{-\mathfrak{N}\cst{\mathcal{Y}}} if $l$ is even and
$z_a$ has grade~$1$, and times~$1$ otherwise.
\end{definition}

\begin{definition}[theta data, lifts and the unfolding step; recalled]\label{d:theta}
For hermitian spaces with frame $V$ and $V'$ the space \m{\Hom(V,V')} is a hermitian space with frame. For a
hermitian space with frame $W$ the library provides the space $\mathcal{S}(W)$ of theta data with the Heisenberg
operators \m{\varrho(y)} of the adelic vectors~$y$ (and \m{\varrho_f(y)} of the finite ones), the Weil
representation~$\omega$ of~$G(W)$, the theta functional~$\theta$, the vacuum, the \emph{finite part}
$\mathcal{W}_f(W)$ spanned by the translates \m{\varrho_f(y)} of the vacuum, and the datum $v\tensor f$ made
of a vector $v$ of the Fock model at the infinite places and $f\in\mathcal{S}(W)$. For a character $\nu$ of the
norm-one adeles the pair \m{G(V)\times G(V')} acts on $\mathcal{S}(\Hom(V,V'))$ by the Weil representation
twisted by~$\nu$ and powers of~$\upsilon$, the \emph{twisted pair action} $\pW{\nu}$ of the note \emph{Two
consecutive theta lifts}; its restriction to finite points acts on the finite part. The \emph{theta lift} with
datum $\Phi$ of a function $F$ on $[G(V)]$ is
$h\mapsto\int_{[G(V)]}\theta(\pW{\nu}(g,h)\Phi)\,F(g)\,dg$. Let $\mathcal{Y}$, $\mathcal{Y}'$ be gradings of
$V$, $V'$ with the same norm and strings of lengths $l$ and \m{l-1}. The \emph{string map}
$\Tm{\mathcal{Y}}{\mathcal{Y}'}\colon V\to V'$ is the rational matrix that kills $X$ and sends the string vector
$z_a$ of $\mathcal{Y}$ to a multiple of the string vector of $\mathcal{Y}'$ with the same number~$a$ (and the last
one to~$0$); the multiple is $1$, \m{-\mathfrak{N}\cst{\mathcal{Y}}} or $\cst{\mathcal{Y}}$ for $z_a$ of positive, zero or
negative grade if $l$ is odd, and $\cst{\mathcal{Y}}$ or $1$ for positive or negative grade if $l$ is even. A
matrix \m{n\colon V\to V'} is read as a vector of \m{\Hom(V,V')} in the coordinates of the frames, and the
\emph{degree} of its position $(i',i)$ is \m{\gr(i')-\gr(i)}. For \m{g\in G(V)}, \m{h\in G(V')} the \emph{constant
term at the string map} $\mathrm{ct}_{\mathcal{Y},\mathcal{Y}'}(\nu,\Phi)(g,h)$ is the mean of
\[\theta\bigl(\varrho(n)\,\varrho(\Tm{\mathcal{Y}}{\mathcal{Y}'})\,\pW{\nu}(g,h)\Phi\bigr)\]
over the adelic matrices \m{n\colon V\to V'} with entries only in positions of positive degree, modulo the
rational ones. For a function $F$ on $G(V)$ and \m{x\in G(V)} the \emph{unfolded function} is
\[\bigl(\mathrm{unf}_{\mathcal{Y},\mathcal{Y}'}(\nu,\Phi;F,x)\bigr)(h)=\int_{[\Bes{\mathcal{Y}}]}F(\varsigma x)\;
\mathrm{ct}_{\mathcal{Y},\mathcal{Y}'}(\nu,\Phi)(\varsigma x,h)\,d\varsigma .\]
\end{definition}

\statement{ana}{\textbf{Inputs 1 (the background library).}
\phantomsection\inputslabel{1}{in:lib}\enspace In the notation of Definitions~\ref{d:grading} and~\ref{d:theta}:
\begin{enumerate}[label=(\alph*)]
\item \emph{Genericity is inherited.} If $F$ is generic for $\mathcal{Y}$ and
  $\cst{\mathcal{Y}'}=\cst{\mathcal{Y}}$, then $\mathrm{unf}_{\mathcal{Y},\mathcal{Y}'}(\nu,\Phi;F,x)$ is generic
  for~$\mathcal{Y}'$.
\item \emph{One unfolding step.} If $F$ is generic for~$\mathcal{Y}$ and $z$ is an automorphic, rapidly decreasing
  function on $G(V')$, then $\mathrm{unf}_{\mathcal{Y},\mathcal{Y}'}(\nu,\Phi;F,x)\cdot z$ is integrable on
  $\Besr{\mathcal{Y}'}\backslash G(V')$, and its integral there equals
  $\int_{[\Bes{\mathcal{Y}}]}F(\varsigma x)\int_{[G(V')]}\theta(\pW{\nu}(\varsigma x,h)\Phi)\,z(h)\,dh\,d\varsigma$.
\item \emph{Lift and descent are transposes.} A datum $\Phi$ is an \emph{arrow} from a type $\mathsf{q}$ on $V$ to
  a type $\mathsf{q}'$ on $V'$ if every form of type $\mathsf{q}$ is rapidly decreasing and its theta lift with
  datum $\Phi$ is a form of type~$\mathsf{q}'$. If $\Phi$ is an arrow from $\mathsf{q}$ to $\mathsf{q}'$ and, for the
  pair with the two groups exchanged, an arrow from $\cj{\mathsf{q}'}$ to $\cj{\mathsf{q}}$, then for a form $F$
  of type $\mathsf{q}$ and a form $z$ of type $\cj{\mathsf{q}'}$ the integral over $[G(V')]$ of $z$ times the lift
  of~$F$ equals the integral over $[G(V)]$ of $F$ times the lift of $z$ for the exchanged pair (the \emph{descent}
  of~$z$).
\item \emph{Adapted bases exist} (Hermite reduction over $\Oh$), \emph{and Witt's theorem holds over~$\Oh$}: if
  $E$ is stably unramified and two matrices over $\Oh$ with $n$ rows and equally many columns have left inverses
  over $\Oh$ and equal Gram matrices for the standard form, then a matrix in \m{\GL_n(\Oh)} preserving the standard
  form carries one to the other.
\end{enumerate}}

\statement{ana}{\textbf{Inputs 1, continued (the background library).}\enspace In the same notation:
\begin{enumerate}[label=(\alph*),start=5]
\item \emph{Discrete-series edges are arrows.} If the parameters, weights and signs of $V$ and $V'$ correspond in
  the sense of the unit \emph{Exchange}, then the vector of the Fock model given by the product of the chosen
  joint harmonic polynomials is an \emph{edge} between the two discrete-series types; and an edge, tensored with
  any element of the finite part, is an arrow in the sense of~(c) from the first type to the second and, for the
  pair with the two groups exchanged, from the conjugate of the second to the conjugate of the first, for
  characters of the weights of the two spaces.
\item \emph{Data named modulo a modulus.} For a modulus~$\mathfrak{m}$, integral finite points \m{k\in\Kz{V}},
  \m{k'\in\Kz{V'}} and a residue matrix $y$ modulo~$\mathfrak{m}$, with \m{y'=\cj{k'}^{-1}y\,\cj{k}} for the
  reductions $\cj{k}$, $\cj{k'}$ modulo~$\mathfrak{m}$: the operator \m{\pW{\nu}((1,k),(1,k'))} carries the vector
  named on the vacuum by the indicator of $y'$ to the vector named by the indicator of~$y$, times the value of
  the twisting character of $V$ at \m{\det k'} and the inverse of the value of that of $V'$ at \m{\det k}.
\item \emph{Constant terms are zero terms.} There is a splitting of \m{\Hom(V,V')}, in the sense of the note
  \emph{Two consecutive theta lifts}, whose isotropic frame consists of the basis vectors at the positions of
  positive degree, such that for every datum $\Phi$ and every adelic matrix \m{y\colon V\to V'} the mean of
  \m{\theta(\varrho(n)\varrho(y)\Phi)} over the matrices $n$ of positive degree, as in Definition~\ref{d:theta},
  is the theta function of the zero term of \m{\varrho(y)\Phi} along that splitting.
\item \emph{Kernels of definite pairs give theta lifts.} For standard definite spaces of ranks \m{0<n<n'}, weights
  that fit, a level, a non-zero scale and characters of the two weights, the lifts of algebraic modular forms by
  the theta kernels of the pair form a family of theta lifts in the sense of Input~5 of the note \emph{Two
  consecutive theta lifts}: the theta integral of a coefficient of a form against a datum with a polynomial vector
  at the infinite places is a coefficient of the lift of the form.
\item \emph{Cut-off functions and integration in stages.} A proper action of a group with a right invariant
  measure on a $\sigma$-compact Hausdorff space has a cut-off function, a non-negative function whose integral
  over every orbit is~$1$. If $\Delta$ is a subgroup of a subgroup $H$ of a group~$G$, and a map
  \m{(x,h)\mapsto h\,\mathsf{s}(x)} carries the product of a measure on a space of parameters $x$ and the measure
  of $H$ to the measure of~$G$, then for a measurable, left $\Delta$-invariant function $F$ that is integrable on
  \m{\Delta\backslash G}, the inner integrals \m{\int_{\Delta\backslash H}F(h\,\mathsf{s}(x))\,dh} are integrable
  in~$x$ and their integral is \m{\int_{\Delta\backslash G}F}; the quotients carry the measures induced as in
  Definition~\ref{d:space}.
\end{enumerate}}

\subhead{Status of Inputs~1.} For (a) and (b) the Lean file cites B.~Wang Peng Jun, \emph{Global theta lifting and automorphic
periods associated to nilpotent orbits}, J.~Number Theory 271 (2025), Prop.~3.2 and Thm.~4.5, and for the local
case R.~Gomez and C.-B.~Zhu, GAFA 24 (2014); the authors mark the first as compared, with the place confirmed. Our
catalogue has not confirmed these places: \emph{reference to be confirmed}. For (d) the catalogue lists Witt's
theorem over $\Oh$ in the special case of orthonormal columns, which the library deduces from the statement with
equal Gram matrices used here, and gives for it U.~A.~First, \emph{Witt's extension theorem for quadratic spaces
over semiperfect rings}, Remark~2.1 and Corollary~4.9, applied on each orbit of conjugation on the finite places;
for Hermite reduction it records ``Hermite normal form over a valuation ring'' without a place. For the first
statement of (e) the catalogue cites Ichino, \emph{Theta lifting for discrete series representations of real
unitary groups} (arXiv:2002.09148), Theorem~4.1 and Lemma~7.8, pp.~35--36, for the compact types, and records the
verdict of its second reading: the formal matching condition selects witnesses for a two-sided analytic
statement beyond the published joint-type criterion. So the formal statement is not a special case or a weaker
form of the cited ones: it goes beyond them, and its relation to the literature is unresolved; it is used here
as the library proves it. Statement~(c), the second statement of~(e) and~(f) are classed elementary. For (h) the library's
heading is ``Theta lift for a definite unitary pair as a class sum on algebraic modular forms''; for the
class-sum formula behind it the catalogue names B.~Gross, \emph{Algebraic modular forms}, Israel J. Math.\ 113
(1999), Kashiwara--Vergne (1978) and R.~Howe, \emph{Transcending classical invariant theory}, J.~Amer.\ Math.\
Soc.\ 2 (1989), place not confirmed. For (g) and (i) the catalogue has no entry: \emph{reference to be
confirmed}.

\subhead{The data.}

\begin{definition}[index sets and blocks]\label{d:index}
In Setting~\ref{s:data}, for a natural number $j$ the map
$\inclY{j}\colon\YI{D}{j}\to\YI{D}{j+1}$ is the inclusion of the index set of $\Ys{D}{j}$ as the first summand of
$\YI{D}{j+1}=\YI{D}{j}\sqcup\{0,\dots,\tD{D}-1\}$. Every $\YI{D}{j}$ is non-empty. For sets $R_0$ and $T_0$ and a
natural number~$j$, let $R_0$ followed by $j$ copies of $T_0$ be the disjoint union built by adding one copy of
$T_0$ at a time; for an index $i$ of it the \emph{block number} \m{\blk{j}{i}\in\{0,\dots,j\}} is $0$ if $i$ lies
in~$R_0$ and $\nu$ if $i$ lies in the $\nu$-th copy of~$T_0$. A square matrix indexed by such a set is \emph{block
upper triangular} if its entry at $(i,i')$ vanishes whenever \m{\blk{j}{i}>\blk{j}{i'}}, and for
\m{\nu\in\{0,\dots,j\}} its $\nu$-th \emph{diagonal block} is the square submatrix on the indices of block
number~$\nu$.
\end{definition}
\leanline{\lean{inclY}\sep\lean{instNonemptyY}\sep\lean{blockNumber}}

\begin{definition}[forms of principal level; recalled from the unit \emph{Trace}]\label{d:forms}
In Setting~\ref{s:data}, for a hermitian space with frame $V$ over $\ED{D}$, \m{\Kch{D}{V}} is the level of the
elements of $\Kz{V}$ that are congruent to~$1$ modulo every chosen place, the \emph{principal level at the chosen
places}. The forms on $\Ys{D}{0}$ of the type of the rows that are fixed by \m{\Kch{D}{\Ys{D}{0}}} are the
\emph{forms of principal level}; the group $\GD{D}$ acts on them through the reductions of $\Kz{\Ys{D}{0}}$ at the
chosen places, a form $F$ being sent by \m{\gamma\in\GD{D}} to~$\gamma F$.
\end{definition}

\begin{definition}[the gradings of the spaces $\Ys{D}{j}$]\label{d:gradings}
In Setting~\ref{s:data}:
\begin{enumerate}[label=(\alph*)]
\item for a natural number $j$ with \m{j\le2\rhoD{D}+1}, $\gY{j}$ is the grading of $\Ys{D}{j}$ as
  $\XI{D}{j}$ plus a string of length \m{2\rhoD{D}+1-j=m-1-j} with norm $\ND{D}$, with the bijection onto
  $\YI{D}{j}$ fixed in the unit \emph{Trace} and with the constant
  $\cst{\gY{j}}=-\bigl(\mathbf{j}\,\gD{D}\,\vpi{P}{\mu}\bigr)^{-1}$, where $\vpi{P}{\mu}$ is the generator of the
  different of $\ED{D}$ fixed in Chapter~III;
\item $\gYlast$ is the grading of the standard space $\Ys{D}{m}$ as a set of \m{\nD{D}{m}-1} elements, the set
  $\XI{D}{m}$ without its last element, plus a string of length~$0$, with norm $\ND{D}$ and constant~$1$. In the
  formal text it is given for every number $M$ together with a proof that \m{m=M}.
\end{enumerate}
The last element of $\XI{D}{j}$ is called the \emph{point} of~$\XI{D}{j}$, and $\XIo{D}{j}$ is the set $\XI{D}{j}$
without its point.
\end{definition}
\leanline{\lean{gradingY}\sep\lean{lastGrading}}

For \m{j\le m-1} the space $\Ys{D}{j}$ is thus the orthogonal sum of the definite space $\Xs{D}{j}$ and a string of
length \m{m-1-j}, one vector shorter after each lift. A non-empty string makes the space indefinite (by its
hyperbolic planes and, for odd length, its middle vector of norm \m{-\ND{D}}). For \m{j=m-1} the string is empty
and $\Ys{D}{m-1}$ is definite; $\Ys{D}{m}$ is the standard definite space of rank \m{\nD{D}{m}-1}, whose index
set the grading $\gYlast$ identifies with~$\XIo{D}{m}$.

\begin{definition}[one lift of the indefinite chain]\label{d:lift}
In Setting~\ref{s:data} let $\vartheta$ be a tame character of $D$ and $j$ a natural number.
\begin{enumerate}[label=(\alph*)]
\item $\levelY{\vartheta}{j}$ is the level of $\Ys{D}{j}$ of the elements of $\Kz{\Ys{D}{j}}$ that are congruent
  to~$1$ modulo every chosen place and modulo every place of~$\Qs{\vartheta}$.
\item $\ltame{\vartheta}{j}=\psic{D}{\vartheta^{j+1}}{j+1}$ is the tame part of the character
  \m{\muc{D}{\vartheta^{j+1}}{j+1}} of the lift number \m{j+1}: a family of characters of the residue fields, equal
  to $\omega_{D,w}^{j+1}(\vartheta^{\mathrm{t}}_w)^{j+1}$ at~$w$.
\item \m{\kerY{\vartheta}{j}} is the theta kernel of the pair \m{(\Ys{D}{j},\Ys{D}{j+1})}: the space of theta data
  of the space of homomorphisms from $\Ys{D}{j}$ to $\Ys{D}{j+1}$, with the twisted pair action $\pW{\nu}$ of
  Definition~\ref{d:theta} for \m{\nu=\muc{D}{\vartheta^{j+1}}{j+1}} and with the theta functional.
\item \m{\ups{\vartheta}{j}} is the element of the finite part of this space obtained from the vacuum by the
  operator named by the indicator function of the matrices that are integral and congruent to the block
  inclusion matrix of $\inclY{j}$ modulo the product of the chosen places and of the places of~$\Qs{\vartheta}$.
\item \m{\FormY{j}} is the space of the forms on $\Ys{D}{j}$ of the type given by the parameter of
  $\Ys{D}{j}$ fixed in the unit \emph{Exchange}, with the Petersson product; \m{\cj{\FormY{j}}} is the space of
  the forms of the conjugate type.
\item \m{\datY{\vartheta}{j}} is the datum $v\tensor\ups{\vartheta}{j}$, where $v$ is the vector of the Fock
  model given by the joint harmonic polynomial of the pair fixed in the unit \emph{Exchange}.
\item For \m{F\in\FormY{j}}, $\liftY{\vartheta}{j}{(F)}\in\FormY{j+1}$ is
  $[\Kz{\Ys{D}{j}}:\levelY{\vartheta}{j}]$ times the theta lift of $F$ with the kernel $\kerY{\vartheta}{j}$ and
  the datum $\datY{\vartheta}{j}$ (the lift as a map between spaces of forms, as chosen in the library).
\end{enumerate}
\end{definition}
\leanline{\lean{levelY}\sep\lean{liftTame}\sep\lean{thetaKernelY}\sep\lean{upsilonY}\sep\lean{FormY}\sep
\lean{datumY}\sep\lean{liftY}}

\begin{definition}[the chain of lifts]\label{d:chain}
In Setting~\ref{s:data} let $\vartheta$ be a tame character of~$D$. For a natural number $j$ and a form $F$ of
principal level on~$\Ys{D}{0}$,
\[\chainY{\vartheta}{j}{(F)}\;=\;\liftY{\vartheta}{j-1}{\bigl(\cdots\liftY{\vartheta}{0}{(F)}\cdots\bigr)}\in\FormY{j}\]
is the composite of the first $j$ lifts. The space $\Ys{D}{m}$ is definite and standard, and its forms of the type
of $\FormY{m}$ are identified with algebraic modular forms on $\Ys{D}{m}$; \m{\iot{\vartheta}} is the linear map
that sends a form $F$ of principal level on $\Ys{D}{0}$ to the algebraic modular form that corresponds to
$\chainY{\vartheta}{m}{(F)}$.
\end{definition}
\leanline{\lean{liftChainY}\sep\lean{decl}}

\begin{definition}[the level of the image]\label{d:lastlevel}
In Setting~\ref{s:data} let $\vartheta$ be a tame character of~$D$.
\begin{enumerate}[label=(\alph*)]
\item \m{\lastRes{\vartheta}} is the group of the families \m{\kappa=(\kappa_w)}, indexed by the chosen places and
  the places of $\Qs{\vartheta}$, of matrices $\kappa_w\in\GL_{\YI{D}{m}}(\mathbb{F}_w)$ that are block upper
  triangular for the blocks of $\YI{D}{m}$ and equal to~$1$ at the chosen places.
\item For \m{\kappa\in\lastRes{\vartheta}}, \m{\lastChar{\vartheta}{\kappa}} is the product, over these places $w$
  and over \m{\nu\in\{1,\dots,m\}}, of the values of \m{\ltame{\vartheta}{\nu-1}} at $w$ on the determinant of the
  $\nu$-th diagonal block of~$\kappa_w$.
\item For a form $F$ of principal level on~$\Ys{D}{0}$, \m{\iotL{\vartheta}{(F)}} is \m{\iot{\vartheta}(F)},
  regarded as an element of the space of the algebraic modular forms on $\Ys{D}{m}$ on which every element of
  $\Kz{\Ys{D}{m}}$ with reductions \m{\kappa\in\lastRes{\vartheta}} acts by~\m{\lastChar{\vartheta}{\kappa}}
  (Proposition~\ref{p:equiv} shows that it lies there).
\item For \m{\gamma\in\GD{D}}, \m{\groupLift{\vartheta}{\gamma}} is a chosen element of $\Kz{\Ys{D}{m}}$ whose
  reduction is the block diagonal matrix with blocks $\gamma_w$ and $1$ at every chosen place~$w$, and $1$ at the
  places of~$\Qs{\vartheta}$.
\end{enumerate}
\end{definition}
\leanline{\lean{lastResLevel}\sep\lean{lastLevelChar}\sep\lean{iotaAtLevel}\sep\lean{groupLift}}

\begin{setting}[filling data that fit a space; recalled from the unit \emph{Filling}]\label{s:filling}
Let $\mathsf{S}$ be filling data over $E$ of a shape with a number of hyperbolic planes (Chapter~VI): among its
components are a prime~$N$, two further primes with product~$g_{\mathsf{S}}$, their chosen places, a generator
$\varpi_{\mathsf{S}}$ of the different, a level $K_{\mathsf{S}}$ on the definite part $X_0$, and Bessel systems
at the infinite places with their vectors \m{W(g_\infty)} of Bessel functions. Let $V$ be a hermitian space with
frame, $\mathcal{Y}$ a grading of $V$ as the set of the shape plus a string, and \m{\eta\in G_f(X_0)}; assume
that $\mathsf{S}$ \emph{fits} $V$ with a parameter, weights, the grading $\mathcal{Y}$ and~$\eta$: the Bessel
systems have the types of the parameter and the weights, the bijection of $\mathcal{Y}$ is the standard one of the
shape, the type of the parameter is an idempotent kernel, $\eta$ is central, and $x_{V,f}=u\cdot\eta\cdot\Lambda\cdot k$
for a finite idele~$u$, the lattice basis $\Lambda$ of the filling data, and an integral matrix $k$ with integral
inverse that is congruent to~$1$ at the chosen places. For a vector \m{\phi=(\phi_a)} of coefficients, indexed by
the pairs $a$ of a class of $X_0$ for $K_{\mathsf{S}}$ and a Bessel coordinate, and a subgroup $\Delta$ of
$\Gamma(V)$, the \emph{Poincar\'e series} is
\[\mathrm{P}_{\mathsf{S}}(\phi,\Delta)(g)=\sum_a\phi_a\sum_{\gamma\in\Delta}\mathrm{cb}_a(\gamma g),\]
where the cut-off Bessel function $\mathrm{cb}_a$ is a cut-off times a Bessel function of the real component if a
translate of the finite component, determined by the class of~$a$ and by~$\eta$, lies in the principal level
\m{K^{\mathrm{pr}}_{\mathsf{S}}} at the chosen places, and $0$ otherwise. The \emph{column coordinates}
$\mathrm{cc}_{\mathsf{S}}$ send an algebraic modular form on $X_0$ of level $K_{\mathsf{S}}$ to the vector of the
Bessel coordinates of its values at representatives of the classes, the Weyl coordinates being changed to Bessel
coordinates by a fixed linear isomorphism~$\mathsf{A}$.
\end{setting}

\begin{definition}[columns, the Poincar\'e series and the projection to rows]\label{d:poin}
In Setting~\ref{s:data}, $\Col{D}$ is the space of the algebraic modular forms on $\Xs{D}{0}$ of the weight fixed in
the unit \emph{Columns} that are fixed by the level $\Kst{D}{0}$, and $\Colram{D}$ is its subspace of ramified
columns (Definition~\ref{d:columns}). The filling data of $D$ fit $\Ys{D}{0}$ with the grading $\gY{0}$
(Proposition~\ref{p:fits}).
\begin{enumerate}[label=(\alph*)]
\item \m{\colco{D}} is the linear isomorphism from $\Col{D}$ to the space of coefficient vectors that first
  translates a column by the central element \m{\lsc{D}{0}} and then takes the column coordinates of
  Setting~\ref{s:filling}.
\item \m{\PoinD{D}} is the linear map from $\Col{D}$ to the forms of principal level on $\Ys{D}{0}$ that sends
  $\varphi$ to the Poincar\'e series of \m{\colco{D}(\varphi)} over all rational points (recalled from the unit
  \emph{Filling}).
\item For \m{\gamma\in\GD{D}}, \m{\gpt{D}{\gamma}} is the linear map from the forms of principal level on $\Ys{D}{0}$
  to the image of $\ezero{D}$ that sends $F$ to \m{|\GD{D}|\cdot\ezero{D}(\gamma F)}.
\end{enumerate}
\end{definition}
\leanline{\lean{columnCoords}\sep\lean{goodProjTranslate}}

\subhead{The definite side.}

\begin{definition}[admissible combinations; recalled from the note \emph{Two consecutive theta lifts} of this
unit]\label{d:adm}
For a natural number $n$, \m{\Oaw{n}} is the ring of the complex numbers $z$ with \m{|\iota(z)|_\ell\le1} for every
prime $\ell$ not dividing~$n$ and every ring homomorphism $\iota\colon\CC\to\CC_\ell$. Let $v$ be a family of
vectors of a complex vector space, $\mathcal{I}$ a set of its indices, $n$ and $d$ natural numbers, $p$ a proposition and $B$ a
real number. Then $\Adm{v}{\mathcal{I}}{n}{d}{p}{B}$ is the set of the vectors \m{\sum_xr(x)v(x)}, where $r$ is a finitely
supported complex function with support in~$\mathcal{I}$, with \m{d\cdot r(x)\in\Oaw{n}} for all~$x$ and, if $p$ holds, with
\m{\sum_x|\sigma(r(x))|\le B} for every field automorphism $\sigma$ of~$\CC$. Such a vector is an \emph{admissible
combination} of the $v(x)$, \m{x\in\mathcal{I}}, away from~$n$, with denominator~$d$ and, under~$p$, size~$B$.
\end{definition}

\begin{definition}[columns, stages and definite chains; recalled from the units \emph{Columns} and
\emph{Exchange} and from the note \emph{Chains of definite theta lifts and the Gauss-sum operator} of this
unit]\label{d:columns}
In Setting~\ref{s:data}, for a natural number $j$:
\begin{enumerate}[label=(\alph*)]
\item $\Kst{D}{j}$ is the level of $\Xs{D}{j}$ of the elements that are congruent to~$1$ modulo the chosen places
  and whose last row, the row of the point, is congruent to \m{(0,\dots,0,1)} modulo~$\piN{D}$; $\Stg{D}{j}$ is the
  space of the smooth algebraic modular forms on $\Xs{D}{j}$ of the weight of stage~$j$; the ramified columns are
  the image on $\Col{D}$ of \m{\prod_{w\in\PiD{D}}(1-\zw{D}{w})}, where $\zw{D}{w}$ is the translation by the scalar
  matrix of an integral adele of norm one, fixed in the unit \emph{Columns}, that generates the units modulo~$w$
  and is $1$ at the places other than $w$ and~$\cj w$;
\item for a tame character $\vartheta$ of~$D$, \m{\kerX{D}{\vartheta}{j}} is the theta kernel of the definite pair
  \m{(\Xs{D}{j},\Xs{D}{j+1})} with the character \m{\muc{D}{\vartheta^{j+1}}{j+1}}; its data are the elements of
  the finite part of $\mathcal{S}(\Hom(\Xs{D}{j},\Xs{D}{j+1}))$, and the lift with a datum $\Phi$ is a linear map
  $\Stg{D}{j}\to\Stg{D}{j+1}$;
\item $\Tr{D}$ is the set of the tuples \m{t=(t_j)_{j\in\NN}} of finite adelic matrices $t_j$ with rows
  $\XI{D}{j+1}$ and columns~$\XI{D}{j}$; for a natural number $q$, $\Tadm{D}{q}$ is the set of the $t$ for which
  \m{\pi_g\cdot t_j} is integral at every place not dividing~$q$, for all~$j$ (the weight is~$\pi_g$, the conjugate of
  $\piGb{D}$); \m{\tdat{D}{t}{j}} is the
  normalised vacuum of \m{\Hom(\Xs{D}{j},\Xs{D}{j+1})} translated by \m{\varrho_f(t_j)};
\item for data \m{\Phi=(\Phi_j)}, $\chn{D}{\vartheta}{\Phi}{j}\colon\Stg{D}{0}\to\Stg{D}{j}$ is the composite of
  the first $j$ lifts with the data \m{\Phi_0,\dots,\Phi_{j-1}}; for paired data \m{\Phi=(\Phi_i)} with $\Phi_i$ in the
  tensor product of the data spaces of the lifts \m{2i+1} and \m{2i+2}, $\chnn{D}{\vartheta}{\Phi}{p}$ is the
  composite of the first $p$ pairs of lifts, where a pair with datum $s\tensor s'$ is the lift with $s'$ after the
  lift with~$s$, extended linearly; for a finitely supported $\Psi\colon\Tr{D}\to\CC$,
  $\chcomb{D}{\vartheta}{\Psi}{j}=\sum_t\Psi(t)\,\chn{D}{\vartheta}{\tdat{D}{t}{\cdot}}{j}$;
\item $\resY{D}\colon\Stg{D}{m}\to\{\text{algebraic modular forms on }\Ys{D}{m}\}$ is the restriction of the unit
  \emph{Columns}: restriction to the block diagonal elements \m{\operatorname{diag}(h,1)}, followed by a branching
  matrix and multiplication by a power of the conjugates of~$\piN{D}$.
\end{enumerate}
\end{definition}

\begin{definition}[translate sums, pair translates and reduced vacua]\label{d:redvac}
In Setting~\ref{s:data} let $i$ be a natural number. The \emph{translate sum} of the note on the Gauss-sum
operator, for the index sets $\XIo{D}{2i}$, $\XI{D}{2i+1}$, $\XIo{D}{2i+2}$, the set of places $\PiD{D}$ and the
prime~$\ND{D}$, is the endomorphism
\[\transum{\XIo{D}{2i}}{\XI{D}{2i+1}}{\PiD{D}}{\ND{D}}{\XIo{D}{2i+2}}\;=\;\sum_{y}\varrho_f(\mathrm{c}(y))\tensor\varrho_f(\mathrm{r}(y))\]
of \m{\mathcal{W}_f(\Hom(\Xs{D}{2i},\Xs{D}{2i+1}))\tensor\mathcal{W}_f(\Hom(\Xs{D}{2i+1},\Xs{D}{2i+2}))}: the sum
over the vectors $y$ with entries in \m{\prod_{w\in\PiD{D}}\ZZ/\ND{D}}, indexed by~$\XI{D}{2i+1}$, where
$\mathrm{c}(y)$ is the matrix whose last column, that of the point of $\XI{D}{2i}$, is \m{-y/\ND{D}} at the places
of $\PiD{D}$ and which vanishes elsewhere, and $\mathrm{r}(y)$ the matrix whose last row, that of the point of
$\XI{D}{2i+2}$, is the conjugate of that column. For a family \m{v=(v_j)_{j\in\NN}} of finite data
\m{v_j\in\mathcal{W}_f(\Hom(\Xs{D}{j},\Xs{D}{j+1}))} the \emph{pair translates} are
\[\ptr{D}{v}{i}\;=\;\transum{\XIo{D}{2i}}{\XI{D}{2i+1}}{\PiD{D}}{\ND{D}}{\XIo{D}{2i+2}}\bigl(v_{2i}\tensor v_{2i+1}\bigr),\]
and for \m{t\in\Tr{D}} (Definition~\ref{d:columns}(c)) the \emph{reduced vacuum} of $t$ is the pure tensor
\[\redvac{D}{t}=\ptr{D}{\tdat{D}{t}{\cdot}}{\rhoD{D}}\tensor\dots\tensor\ptr{D}{\tdat{D}{t}{\cdot}}{0},\]
formed with the family \m{\tdat{D}{t}{\cdot}=(\tdat{D}{t}{j})_{j}} of the translated vacua of~$t$.
\end{definition}
\leanline{\lean{reducedVacuum}}

\begin{definition}[good data and the map $\Bmu{D}{\vartheta}{\Psi}$; recalled from the units \emph{Columns} and
\emph{Exchange}]\label{d:good}
In Setting~\ref{s:data} let $\vartheta$ be a tame character of $D$ and $\Psi\colon\Tr{D}\to\CC$ finitely supported.
The map $\Bmu{D}{\vartheta}{\Psi}$ from $\Col{D}$ to the algebraic modular forms on $\Ys{D}{m}$ is
$\varphi\mapsto\resY{D}\bigl(\chcomb{D}{\vartheta}{\Psi}{m}(\varphi)\bigr)$. For a twist $\vartheta$ of~$D$ the
function $\Psi$ is a \emph{good datum for~$\vartheta$} if
\begin{enumerate}[label=(\roman*)]
\item its support lies in \m{\Tadm{D}{\lvn{\vartheta}\ND{D}}}, its values lie in \m{\Oaw{\lvn{\vartheta}}}, and if
  \m{\Qs{\vartheta}=\emptyset}, then $\sum_t|\sigma(\Psi(t))|\le\exp(x^{60})$ for every field automorphism
  $\sigma$ of~$\CC$;
\item it is \emph{jointly invariant}: for every tuple \m{(a_j)} with \m{a_j\in\Kst{D}{j}} and \m{a_j=1} for
  \m{j\ge m}, and every linear map $\beta$ on the tensor product of the data spaces of the first $m$ lifts,
  $\sum_t\Psi(t)\,\beta\bigl(\bigotimes_{j<m}\omega(a_j,a_{j+1})\tdat{D}{t}{j}\bigr)=\sum_t\Psi(t)\,\beta\bigl(\bigotimes_{j<m}\tdat{D}{t}{j}\bigr)$;
\item for every \m{\varphi\in\Stg{D}{0}}:
  \[\chcomb{D}{\vartheta}{\Psi}{m}(\varphi)=\sum_t\Psi(t)\,\chnn{D}{\vartheta}{\ptr{D}{\mathsf{P}_\Pi\tdat{D}{t}{\cdot}}{\cdot}}{\rhoD{D}+1}(\varphi),\]
  where $\mathsf{P}_\Pi$ is the projection of the finite part to the vectors that are vacua at the places of
  $\PiD{D}$ and their conjugates, applied to every member of the family \m{\tdat{D}{t}{\cdot}}, and the paired
  data are the pair translates of Definition~\ref{d:redvac}.
\end{enumerate}
\end{definition}

\subhead{Stage functions.}

\begin{definition}[stage functions]\label{d:stage}
In Setting~\ref{s:data} let $\vartheta$ be a tame character of~$D$, $\varphi$ a column in $\Col{D}$,
\m{\xi=(\xi_i)_{i\in\NN}} a tuple of real points \m{\xi_i\in G_\infty(\Ys{D}{i})} and \m{y=(y_i)_{i\in\NN}} a tuple
of finite points \m{y_i\in G_f(\Ys{D}{i})}. For a natural number $j$ with \m{j\le 2\rhoD{D}+2}, a grading
$\mathcal{Y}$ of $\Ys{D}{j}$ with norm $\ND{D}$ and a point $g$ of \m{G(\Ys{D}{j})}, the complex number
$\unf{\vartheta}{\varphi}{\xi}{y}{j}{\mathcal{Y}}(g)$ is defined by recursion on~$j$:
\begin{itemize}
\item $\unf{\vartheta}{\varphi}{\xi}{y}{0}{\mathcal{Y}}$ is the Poincar\'e series
  $\mathrm{P}_{\mathsf{S}}(\colco{D}(\varphi),\Besr{\mathcal{Y}})$ of Setting~\ref{s:filling} over the rational
  points of the Bessel group of~$\mathcal{Y}$, for the filling data of~$D$;
\item $\unf{\vartheta}{\varphi}{\xi}{y}{j+1}{\mathcal{Y}}$ is the unfolded function
  $\mathrm{unf}_{\gY{j},\mathcal{Y}}(\nu,\datY{\vartheta}{j};F,(\xi_j,y_j))$ of Definition~\ref{d:theta}, with
  \m{\nu=\muc{D}{\vartheta^{j+1}}{j+1}} and $F=\unf{\vartheta}{\varphi}{\xi}{y}{j}{\gY{j}}$.
\end{itemize}
\end{definition}
\leanline{\lean{unfoldStage}}

\begin{definition}[evaluated data]\label{d:eval}
In Setting~\ref{s:data}:
\begin{enumerate}[label=(\alph*)]
\item for a natural number~$i$, \m{\pairEvalY{i}} is the \emph{finite evaluation} of the pair of lifts
  $\Ys{D}{2i}\to\Ys{D}{2i+1}\to\Ys{D}{2i+2}$, a linear map
  \[\begin{gathered}\mathcal{S}\bigl(\Hom(\Ys{D}{2i+1},\Ys{D}{2i}^{-}\oplus\Ys{D}{2i+2})\bigr)\longrightarrow\\
  \mathcal{W}_f\bigl(\Hom(\Xs{D}{2i},\Xs{D}{2i+1})\bigr)\tensor\mathcal{W}_f\bigl(\Hom(\Xs{D}{2i+1},\Xs{D}{2i+2})\bigr),\end{gathered}\]
  where the exponent ${}^{-}$ denotes the space with the opposite form: it is the
  finite evaluation, in the sense of the note \emph{Two consecutive theta lifts}, of the middle arrow pair of the
  gradings $\gY{2i},\gY{2i+1},\gY{2i+2}$ if \m{i<\rhoD{D}}, of the last arrow pair of
  $\gY{2\rhoD{D}},\gY{2\rhoD{D}+1},\gYlast$ formed with $\piN{D}$ if \m{i=\rhoD{D}}, and $0$ if \m{i>\rhoD{D}};
\item $\cutB=\gD{D}^{\,\fD{D}\,\nD{D}{m}^{2}}$;
\item for a tame character~$\vartheta$, a tuple \m{y=(y_j)} of finite points \m{y_j\in G_f(\Ys{D}{j})} and a natural
  number~$j$, $\tUps{\vartheta}{y}{j}=\pW{\nu}(y_j,y_{j+1})\,\ups{\vartheta}{j}$ with
  \m{\nu=\muc{D}{\vartheta^{j+1}}{j+1}};
\item for a natural number~$i$, $\evD{\vartheta}{y}{i}=\pairEvalY{i}\bigl(\tUps{\vartheta}{y}{2i}\boxtimes\tUps{\vartheta}{y}{2i+1}\bigr)$,
  where $\boxtimes$ is the arrow product of two data, a datum on
  $\Hom(\Ys{D}{2i+1},\Ys{D}{2i}^{-}\oplus\Ys{D}{2i+2})$.
\end{enumerate}
\end{definition}
\leanline{\lean{pairEvalY}\sep\lean{cutBound}\sep\lean{translatedUpsilon}\sep\lean{evaluatedDatum}}

\subhead{Lattices along a string.}

\begin{definition}[adapted bases]\label{d:adapted}
Let $e\colon X\sqcup\{0,\dots,l-1\}\to I$ be a bijection. The \emph{flag order} of an index \m{i\in I} is
\m{\flag{e}{i}=2l} if $i$ comes from~$X$ and \m{\flag{e}{i}=4a+1} if $i$ is the $a$-th string index; a square matrix
indexed by $I$ is \emph{triangular for the flag} if its entry at $(i,i')$ vanishes whenever
\m{\flag{e}{i}>\flag{e}{i'}}. Let $\mathcal{Y}$ be a grading of a hermitian space with frame~$V$, let
\m{h,\mathbf{b}\in\GL_I(\Af)} and let \m{\beta=(\beta_a)_{a<l}} be a family of finite ideles. Then $\adb{\mathcal{Y}}{h}{\mathbf{b}}{\beta}$
means: $\mathbf{b}$ is triangular for the flag of $e_{\mathcal{Y}}$; \m{\mathbf{b}^{-1}h\in\GL_I(\Oh)}, so that the columns of $\mathbf{b}$ and of
$h$ span the same lattice; and the diagonal entry of $\mathbf{b}$ at the $a$-th string index is~$\beta_a$, for every~$a$. For a
finite point \m{y\in G_f(V)}, a pair $(\mathbf{b},\beta)$ is $\adat{\mathcal{Y}}{y}$ if
$\adb{\mathcal{Y}}{x_{V,f}\,y}{\mathbf{b}}{\beta}$; then $\mathbf{b}$ is an \emph{adapted basis} of the lattice of~$y$ and $\beta$ its
\emph{string diagonal}.
\end{definition}
\leanline{\lean{flagOrder}\sep\lean{IsAdaptedBasis}\sep\lean{IsAdaptedAt}}

\begin{definition}[near lines]\label{d:near}
Let $\mathcal{Y}$ be a grading of~$V$ with a string of length~$l$, $\mathfrak{a}$ a finite idele, $Q$ a finite set of
finite places, $n$ and $k$ natural numbers and \m{h\in G_f(V)}. Then $\near{\mathcal{Y}}{\mathfrak{a}}{Q}{n}{k}{h}$
means that there is a pair $(\mathbf{b},\beta)$ that is $\adat{\mathcal{Y}}{h}$ such that for every string index $a$ and
every finite place~$w$
\[\bigl|\ord_w\beta_{a}-\ord_w\bigl((\mathfrak{a}/\cj{\mathfrak{a}})^{n}\,\epsilon_{a}\bigr)\bigr|\ \le\ \begin{cases}k&\text{if }w\in Q\text{ or }\cj w\in Q,\\0&\text{otherwise,}\end{cases}\]
where \m{\epsilon_{a}=1} if $l$ is odd and, if $l$ is even, \m{\epsilon_{a}=\mathfrak{a}^{-1}} for \m{2a<l} and
\m{\epsilon_{a}=\cj{\mathfrak{a}}} otherwise. In words: the lattice of $h$ has an adapted basis whose string diagonal
has the orders of the reference ideles exactly outside $Q$ and its conjugate, and within $k$ on them.
\end{definition}
\leanline{\lean{HasNearLines}}

\begin{definition}[multiplicity of a cut]\label{d:mult}
Let $H$ be a group with subgroups $S$, $A$ and~$J$, let $T$ be a finite set of left cosets of $J$ in~$H$ and $t$ a
left coset of~$J$. The \emph{multiplicity} $\cov{S}{A}{T}{t}$ is the number of the left cosets of \m{S\cap A} in $S$
that contain an element $\varsigma$ with \m{\varsigma^{-1}t\in T}; it is $0$ if there are infinitely many. When $T$
is stable under left multiplication by \m{S\cap A}, as in all uses, every element of such a coset has this
property.
\end{definition}
\leanline{\lean{coverMultiplicity}}

\statement{ana}{\textbf{Inputs 2 (other units and the other notes of this unit).}\phantomsection\inputslabel{2}{in:units}
\enspace In Setting~\ref{s:data}:
\begin{enumerate}[label=(\alph*)]
\item (\emph{Trace}, unit~140) for \m{j<m} the frame of $\Ys{D}{j}$ has finite part \m{\Lambda_j\,k_j} with
  $\Lambda_j$ the graded lattice basis of $\gY{j}$ for the ideles \m{\ideA{D}=\alpha/\cj{\cgD{D}}}, the principal
  idele $\piId{D}$ of $\piN{D}$ and the scale $\lsc{D}{j}=(\ideA{D}/\cj{\ideA{D}})^{\lceil(m-j)/2\rceil}$, and $k_j$
  integral with integral inverse and congruent to $1$ at the chosen places; here $\cgD{D}$ is a finite idele
  generating the ideal of $\piGb{D}$, equal to $1$ at the places not dividing it, and $\piId{D}$ is the finite
  idele with the components of~$\piN{D}$. If $w$ or $\cj w$ lies in $\Qs{\vartheta}$ for a tame
  character $\vartheta$, then $w$ divides $\lvn{\vartheta}$ and divides neither $\piGb{D}$ nor $\ND{D}$. Moreover
  \m{m=2\rhoD{D}+2}.
\item (\emph{Columns}, unit~145) $\PiD{D}$ is a split level for $\ND{D}$: it is disjoint from its conjugate, the
  places above $\ND{D}$ are those of $\PiD{D}$ and their conjugates, and $\ND{D}$ is a uniformiser at each of them.
  For every \m{w\in\PiD{D}}, every ramified column has the form \m{\psi-\zw{D}{w}\psi} with \m{\psi\in\Col{D}}.
  If \m{t\in\Tadm{D}{q}} and $S$ is a set of places prime to $q$ whose conjugates are not chosen, then
  $\tdat{D}{t}{j}$ is fixed by the projection to the vectors that are vacua at the places of $S$ and their
  conjugates.
\item (\emph{Exchange}, unit~147) \m{|\YI{D}{j}|\le\nD{D}{m}} for \m{j\le m}. For \m{j<m} the parameters, weights
  and signs of $\Ys{D}{j}$ and $\Ys{D}{j+1}$ correspond, in the sense required by Inputs~\ref{in:lib}(e).
\end{enumerate}}

\statement{ana}{\textbf{Inputs 2, continued: the unit \emph{Filling}.}\enspace In Setting~\ref{s:data}:
\begin{enumerate}[label=(\alph*),start=4]
\item (\emph{Filling}, unit~148) $\Kst{D}{0}$ is the preimage of the principal level at the chosen places under the
  inclusion of the finite points of $\Xs{D}{0}$ into those of~$\Ys{D}{0}$, and the principal level
  \m{K^{\mathrm{pr}}_{\mathsf{S}}} of the filling data of~$D$ is \m{\Kch{D}{\Ys{D}{0}}} (in the formal text, by
  definition). \emph{Terms of the series}, for admissible filling data as in Setting~\ref{s:filling}: (1) if a
  term \m{\mathrm{cb}_a(\gamma'g)}, with $\gamma'$ a rational point of the Bessel group, does not vanish at a
  point $g$ whose finite component, translated by the representative of a class and by~$\eta$, lies in the
  principal level, then $a$ belongs to that class and $\gamma'$ to the group of the rational unipotent elements
  that the class conjugates into the principal level; (2) over that group the cut-offs sum to~$1$, so that the
  terms \m{\mathrm{cb}_a(\gamma'g)} sum to the Bessel function of the real component; (3) \emph{strong
  approximation for the unipotent group}: for a finite point $u$ of the unipotent group of the shape and any
  finite point $h$ there is a rational unipotent $\varsigma$ with \m{h^{-1}u^{-1}\varsigma h} in the principal
  level. \emph{Unfolding of Poincar\'e series}: for
  admissible filling data as in Setting~\ref{s:filling}, a form $z$ on $V$ and a subgroup $\Delta$ of $\Gamma(V)$,
  the product of $z$ and $\mathrm{P}_{\mathsf{S}}(\phi,\Delta)$ is integrable on \m{\Delta\backslash G(V)}, with
  integral $\sum_a\phi_a\int_{G(V)}z\cdot\mathrm{cb}_a$. \emph{The last stage}: there is a linear map $R$ from
  $\Stg{D}{m}$ to $\FormY{m}$ such that (1) for all tame $\vartheta$, finitely supported $\Psi$, forms $F$ of
  principal level and columns~$\varphi$, if
  $\chainY{\vartheta}{m}{(F)}=R\bigl(\chcomb{D}{\vartheta}{\Psi}{m}(\varphi)\bigr)$ then
  $\iot{\vartheta}(F)=\Bmu{D}{\vartheta}{\Psi}(\varphi)$; (2) for every linear functional $\mathsf{t}$ on the
  values of forms of stage~$m$ there is a complex number $C_{\mathsf{t}}$ such that for every \m{\Xi\in\FormY{m}},
  every \m{\psi\in\Stg{D}{m}} and every function $F$ on \m{G(\Ys{D}{m})} with
  $F(\incl{\gYlast}(h))=\coef{\mathsf{t}}{\psi}(\operatorname{diag}(h,1))$ for all~$h$:
  $\int_{[G(\Ys{D}{m})]}F\,\cj\Xi=C_{\mathsf{t}}\cdot\langle\Xi,R\psi\rangle$. Here
  $\coef{\mathsf{t}}{\psi}(g)=\mathsf{t}\bigl(g_\infty^{-1}\psi(g_f)\bigr)$ is the \emph{coefficient} of $\psi$
  at~$\mathsf{t}$, and $\langle\,,\rangle$ the Petersson product.
\end{enumerate}}

\statement{ana}{\textbf{Inputs 2, continued: the unit \emph{Curve} and the other notes of this unit.}\enspace In
Setting~\ref{s:data}:
\begin{enumerate}[label=(\alph*),start=5]
\item (\emph{Curve}, unit~150) \emph{named data are admissible combinations}: let $W$ be a hermitian space with
  frame, with index set~$I$, let $\mathfrak{m}$ be a modulus, \m{g'} an integer of $E$ and $q$ a natural number with
  \m{g'q\in\mathfrak{m}}, let $\mathsf{e}$ be an integral finite adele with \m{\mathsf{e}\,q\equiv0}
  modulo~$\mathfrak{m}$ (when none is named, \m{\mathsf{e}=g'}), let $\Psi_0$ be a function of residue vectors
  modulo~$\mathfrak{m}$ indexed by~$I$ (residue matrices, for a space of homomorphisms) with values in $\Oaw{q}$,
  $p$ a proposition and $B$ a real number such that, if $p$ holds, \m{\sum_x|\sigma(\Psi_0(x))|\le B} for every field
  automorphism $\sigma$ of~$\CC$, and let $u$ be a vector of $\mathcal{S}(W)$ and $u_0$ its projection to the
  vectors that are vacua at the places of~$\mathfrak{m}$ and their conjugates. Then the vector named at the places
  of $\mathfrak{m}$ by $\Psi_0$ on $u$ is an admissible combination of the translates \m{\varrho_f(t)u_0}, over the
  finite vectors $t$ with \m{\cj{\mathsf{e}}\,t} integral at the places not dividing~$q$, away from~$q$, with
  denominator the absolute norm of the ideal of~$g'$ to the power~$|I|$ (which is $0$ for \m{g'=0}) and,
  under~$p$, size~$B$. The weight $\cj{\mathsf{e}}$ of the translates is thus set by $\mathsf{e}$ and not by~$g'$.
\item (the note \emph{Two consecutive theta lifts} of this unit) the calculus of admissible combinations
  (substitution multiplies denominators and sizes; linear and bilinear images; pure tensors; multiplication by a
  root of unity; clearing a constant); the Weil translate of a translated vacuum by a finite point $g$ is an
  admissible combination of translated vacua, with coefficient governed by the finite multiplier of~$g$; the
  finite multiplier and the values of the characters $\muc{D}{\vartheta}{j}$ on determinants have $\ell$-adic
  absolute value $1$ at points integral above~$\ell$; \emph{support of the zero term}, \emph{the double unfolding},
  \emph{the double lift}, \emph{equivariance of the finite evaluation} and \emph{admissible evaluation}, each quoted
  where it is used (§\ref{s:stages}, §\ref{s:good}).
\item (the note \emph{Chains of definite theta lifts and the Gauss-sum operator} of this unit) for finite points
  $a_j$ of the spaces $\Xs{D}{j}$ and paired data $\Phi_i$,
  $\chnn{D}{\vartheta}{(\omega(a_{2i},a_{2i+1})\tensor\omega(a_{2i+1},a_{2i+2}))\Phi_i}{p}(a_0\varphi)=a_{2p}\,\chnn{D}{\vartheta}{\Phi}{p}(\varphi)$,
  the points acting on forms by right translation; the characters \m{\muc{D}{\vartheta^{j}}{j}} are unramified at
  the places above~$\ND{D}$, and a scalar pair acts on the finite part through its quotient.
\end{enumerate}}

\sect{The chain of lifts and its level}{s:chainproofs}

\begin{lemma}\label{l:tame}
In Setting~\ref{s:data} let $\vartheta$ be a tame character of~$D$ and $j$ a natural number.
\begin{enumerate}[label=(\roman*)]
\item For every place \m{w\in\Qs{\vartheta}}, the component of \m{\psic{D}{\vartheta}{j}} at $w$ is
  $\vartheta^{\mathrm{t}}_w$.
\item $\levelY{\vartheta}{j}\subseteq\Kch{D}{\Ys{D}{j}}$.
\item If \m{j<m}, then the datum \m{\datY{\vartheta}{j}} is an arrow, for the kernel \m{\kerY{\vartheta}{j}}, from the
  type of \m{\FormY{j}} to the type of \m{\FormY{j+1}}, and for the pair with the two groups exchanged it is an arrow
  from the type of \m{\cj{\FormY{j+1}}} to the type of \m{\cj{\FormY{j}}} (Inputs~\ref{in:lib}(c)).
\end{enumerate}
\end{lemma}
\begin{proof}
(i) $\psic{D}{\vartheta}{j}=\omega_D^{\,j}\vartheta^{\mathrm{t}}$, and $\omega_{D,w}$ is trivial at the places that are
not chosen; a place of $\Qs{\vartheta}$ does not divide $\piGb{D}$ (Inputs~\ref{in:units}(a)). (ii) The
condition that defines $\levelY{\vartheta}{j}$ contains the condition at the chosen places. (iii) The parameters,
weights and signs of $\Ys{D}{j}$ and $\Ys{D}{j+1}$ correspond (Inputs~\ref{in:units}(c)). By the first statement of
Inputs~\ref{in:lib}(e) the vector $v$ of the Fock model given by the joint harmonic is an edge between the two
types, and by the second statement the datum \m{\datY{\vartheta}{j}=v\tensor\ups{\vartheta}{j}}, an edge tensored
with an element of the finite part, is an arrow in both directions; the character of the kernel has the weights of
the two spaces (the note on the characters of the data). The status of this input, on which every lift of the
chain rests, is given after Inputs~\ref{in:lib}.
\end{proof}
\leanline{\lean{liftTameParteqofmemQ}\sep\lean{levelYlechosenLevel}\sep\lean{isArrowdatumY}}

\begin{lemma}[the finite datum is an eigenvector of the level]\label{l:upsilon}
In Setting~\ref{s:data} let $\vartheta$ be a tame character of~$D$ and $j$ a natural number. Let
\m{k\in\Kz{\Ys{D}{j}}} and \m{k'\in\Kz{\Ys{D}{j+1}}}, and let \m{g=(g_w)} and \m{g'=(g'_w)} be the families of their
reductions at the chosen places and the places of~$\Qs{\vartheta}$. Assume that for every such place $w$ the first
columns of $g'_w$, those indexed by $\YI{D}{j}$ through \m{\inclY{j}}, form the matrix
\m{\binom{g_w}{0}}. Then, for the Weil representation $\omega$ of the kernel \m{\kerY{\vartheta}{j}},
\[\omega\bigl((1,k),(1,k')\bigr)\,\ups{\vartheta}{j}\;=\;\Bigl(\prod_{w}\frac{\ltame{\vartheta}{j}_w(\det g'_w)}{\ltame{\vartheta}{j}_w(\det g_w)}\Bigr)\cdot\ups{\vartheta}{j}.\]
\end{lemma}
\begin{proof}
\m{\ups{\vartheta}{j}} is named on the vacuum by the indicator of one residue matrix, the block inclusion
\m{\binom10} modulo the product $\mathfrak{m}$ of the places in question. The Weil operator of an integral pair
\m{(k,k')} carries the vector named by the indicator of a residue matrix $y$ to the vector named by the indicator
of \m{\cj{k'}\,y\,\cj{k}^{-1}}, the bars denoting reduction modulo~$\mathfrak{m}$, times the values of the two
twisting characters at \m{\det k} and \m{\det k'} (Inputs~\ref{in:lib}(f)). The
hypothesis says exactly that $\cj{k'}\binom10\cj{k}^{-1}=\binom10$ at every place. The twisting character is
\m{\muc{D}{\vartheta^{j+1}}{j+1}} times a power of~$\upsilon$; it has tame part \m{\ltame{\vartheta}{j}}, so its
value at the determinant of an element of $\Kz{\,\cdot\,}$ is the product of the values of the tame part at the
reductions of the determinant, and $\upsilon$ is unramified. This gives the quotient.
\end{proof}
\leanline{\lean{thetaKernelYomegaupsilonY}}

\begin{lemma}[blocks and counts]\label{l:blocks}
\begin{enumerate}[label=(\roman*)]
\item For sets $R_0$, $T_0$, a natural number $j$ and an index $i$ of $R_0$ followed by $j$ copies of~$T_0$:
  \m{\blk{j}{i}=0} if and only if $i$ lies in the first block~$R_0$.
\item For finite sets $R_0$ and $T_0$ and a natural number $j$ there is a natural number $\nu$ such that for every
  commutative ring $R$ the group of the invertible matrices over~$R$, indexed by $R_0$ followed by $j$ copies
  of~$T_0$, that are block upper triangular has $|\GL_{R_0}(R)|\cdot|\GL_{T_0}(R)|^{j}\cdot|R|^{\nu}$ elements
  (the cardinality of an infinite set being counted as~$0$).
\item Let $\mathfrak{X}$ be a system of level places of~$E$: two finite families of finite places, the chosen
  places and the twist places, no two of them equal or conjugate. For a finite set $R_0$, the kernel of the
  reduction of $\Kz{\,\cdot\,}$ on $R_0$ modulo the product $\mathfrak{m}_1$ of all these places equals the kernel of
  the family of the reductions at these places.
\item For $\mathfrak{X}$ as in (iii) and a natural number~$e$, the residue ring of $E$ modulo the $e$-th power of
  the product of the chosen places times the product of the twist places has
  $(\prod_{w}\Nm w)^{e}\cdot\prod_{v}\Nm v$ elements, $w$ running over the chosen and $v$ over the twist places.
\item For a finite set $R_0$ and $\mathfrak{X}$ as in (iii), the order of $\GL_{R_0}$ of the residue ring modulo
  $\mathfrak{m}_1$ equals the index in $\Kz{\,\cdot\,}$ of the level of the elements congruent to $1$ at all places
  of~$\mathfrak{X}$.
\end{enumerate}
\end{lemma}
\begin{proof}
(i) By induction on~$j$: a new copy of $T_0$ gets the new number \m{j+1\ne0}, and old indices keep their number.
(ii) By induction on~$j$, with \m{\nu=0} for \m{j=0}. A block upper triangular matrix for \m{j+1} has the block form
$\bigl(\begin{smallmatrix}A&B\\0&C\end{smallmatrix}\bigr)$ with $A$ block upper triangular for~$j$, and it is
invertible exactly when $A$ and $C$ are; this is a bijection onto triples \m{(A,B,C)} with $B$ arbitrary, so $\nu$
grows by the number of entries of~$B$. (iii) Two integral matrices are congruent modulo $\mathfrak{m}_1$ if and
only if they are congruent modulo each of its prime factors, which are distinct. (iv) The norm is multiplicative.
(v) By (iii) the level is the kernel of the reduction modulo~$\mathfrak{m}_1$, which is surjective onto
\m{\GL_{R_0}} of the residue ring.
\end{proof}
\leanline{\lean{blockNumbereqzero}\sep\lean{cardparabolicblockNumber}\sep\lean{kerreductionKeq}\sep
\lean{cardlevelModulusresRing}\sep\lean{cardGLlevelModulus}}

\begin{proposition}[the chain is equivariant]\label{p:equiv}
In Setting~\ref{s:data} let $\vartheta$ be a tame character of~$D$, \m{\gamma\in\GD{D}}, $F$ a form of principal
level on~$\Ys{D}{0}$ and \m{k\in\Kz{\Ys{D}{m}}}, and let \m{g=(g_w)} be the family of the reductions of $k$ at the
chosen places and the places of~$\Qs{\vartheta}$. Assume that every $g_w$ is block upper triangular for the blocks
of~$\YI{D}{m}$, and that at every chosen place $w$ the diagonal block of number $0$ of $g_w$ is~$\gamma_w$. Then
\[k\cdot\iot{\vartheta}(F)\;=\;\Bigl(\prod_{w}\prod_{\nu=1}^{m}\ltame{\vartheta}{\nu-1}_w\bigl(\det g_w^{(\nu)}\bigr)\Bigr)\cdot\iot{\vartheta}(\gamma F),\]
where \m{k\,\cdot} is right translation of algebraic modular forms and $g_w^{(\nu)}$ the $\nu$-th diagonal block
of~$g_w$.
\end{proposition}
\begin{proof}
The identification of forms on the definite space with algebraic modular forms commutes with right translation,
so the claim is an identity for $\chainY{\vartheta}{m}{(F)}$. We prove it for $\chainY{\vartheta}{j}{(F)}$ and
every $j$ for which the data of the first $j$ lifts are arrows, by induction on~$j$; for \m{j=m} this is
Lemma~\ref{l:tame}(iii). For \m{j=0} the chain is the identity, the product is empty, and $k$ acts on a form of
principal level through its reductions at the chosen places, which are~$\gamma$. For the step let $g'_w$ be the
upper left block of $g_w$, on the indices of~$\YI{D}{j}$. It is block upper triangular, and
$\det g_w=\det g'_w\cdot\det g^{(j+1)}_w$, so $g'_w$ is invertible; choose \m{k'\in\Kz{\Ys{D}{j}}} with the
reductions~$g'_w$. For an arrow, right translation by $k$ of the lift of a form equals the lift, with the datum
moved by \m{\omega((1,k'),(1,k))}, of the form translated by~$k'$. The Weil operator acts on
$\datY{\vartheta}{j}=v\tensor\ups{\vartheta}{j}$ through the finite datum, and by Lemma~\ref{l:upsilon} (its
hypothesis is the block triangularity) it multiplies \m{\ups{\vartheta}{j}} by
$\prod_w\ltame{\vartheta}{j}_w(\det g_w)/\ltame{\vartheta}{j}_w(\det g'_w)=\prod_w\ltame{\vartheta}{j}_w(\det g^{(j+1)}_w)$.
The induction hypothesis for $k'$ supplies the factors of the blocks \m{1,\dots,j}, the diagonal blocks of $g'_w$
being those of~$g_w$.
\end{proof}
\leanline{\lean{algFormRepiota}}

Taking \m{\gamma=1} and $k$ with reductions in \m{\lastRes{\vartheta}} shows that \m{\iot{\vartheta}(F)} lies in the
space of Definition~\ref{d:lastlevel}(c), so that \m{\iotL{\vartheta}{(F)}} is defined.

\begin{corollary}[the chain commutes with the finite group]\label{c:groupG}
In Setting~\ref{s:data}, for a tame character $\vartheta$ of~$D$, \m{\gamma\in\GD{D}} and a form $F$ of principal
level on~$\Ys{D}{0}$: $\iot{\vartheta}(\gamma F)=\groupLift{\vartheta}{\gamma}\cdot\iot{\vartheta}(F)$.
\end{corollary}
\begin{proof}
Apply Proposition~\ref{p:equiv} to \m{k=\groupLift{\vartheta}{\gamma}}. Its reductions are block diagonal with blocks
$\gamma_w$ and~$1$; by Lemma~\ref{l:blocks}(i) the entries outside the block of number~$0$ are those of the
identity, so every diagonal block of number \m{\nu\ge1} is the identity and the product of character values is~$1$.
\end{proof}
\leanline{\lean{iotaprincipalRep}}

\sect{The Poincar\'e series at Bessel points}{s:poincare}

\begin{proposition}\label{p:fits}
In Setting~\ref{s:data} the filling data of $D$ fit the space $\Ys{D}{0}$, with the parameter of the rows, the
weights of the columns, the grading $\gY{0}$ and the central element \m{\eta=\lsc{D}{0}^{-1}}, in the sense of
Setting~\ref{s:filling}.
\end{proposition}
\begin{proof}
Four of the five conditions are properties of the chosen objects: the Bessel systems were chosen of these types,
the bijection of $\gY{0}$ is by definition the standard one of the shape, the type of the parameter of the rows is
an idempotent kernel (unit \emph{Trace}), and a scalar is central. For the frame, Inputs~\ref{in:units}(a) gives
\m{\Lambda_0k_0} with $\Lambda_0$ the graded lattice basis for the scale \m{\lsc{D}{0}}; an entrywise comparison,
using that \m{\lsc{D}{0}} has norm one, identifies $\Lambda_0$ with \m{\lsc{D}{0}\cdot\eta\cdot\Lambda}.
\end{proof}
\leanline{\lean{fillingfits}}

\begin{lemma}\label{l:poinbasic}
In Setting~\ref{s:filling} let \m{\phi} be a vector of coefficients.
\begin{enumerate}[label=(\roman*)]
\item For a subgroup $\Delta$ of $\Gamma(V)$, \m{\gamma\in\Delta} and \m{g\in G(V)}:
  $\mathrm{P}_{\mathsf{S}}(\phi,\Delta)(\gamma g)=\mathrm{P}_{\mathsf{S}}(\phi,\Delta)(g)$.
\item If \m{g\in G(V)} and $\mathrm{P}_{\mathsf{S}}(\phi,\Besr{\mathcal{Y}})(g)\ne0$, then there is
  \m{\varsigma\in\Besf{\mathcal{Y}}} such that \m{\varsigma\,g_f} lies in the principal level
  $K^{\mathrm{pr}}_{\mathsf{S}}$, $g_f$ being the finite component of~$g$.
\end{enumerate}
\end{lemma}
\begin{proof}
(i) Reindex the sum over $\Delta$ by \m{\gamma'\mapsto\gamma'\gamma}. (ii) A term \m{\mathrm{cb}_a(\gamma g)} is
non-zero only if a translate of the finite component of \m{\gamma g} by an element of the image of
$\inclf{\mathcal{Y}}$ lies in $K^{\mathrm{pr}}_{\mathsf{S}}$; this translate is \m{\varsigma g_f} with $\varsigma$ the
finite part of an element of the Bessel group, because $\gamma$ and the image of $\incl{\mathcal{Y}}$ lie in it. If
no such $\varsigma$ existed, all terms would vanish.
\end{proof}
\leanline{\lean{poincareratmul}\sep\lean{existsbesselFinmulmemofnezero}}

\begin{proposition}[normal form]\label{p:normal}
In Setting~\ref{s:filling}, with Weil data for~$E$ and $\mathsf{S}$ admissible in the sense of the unit
\emph{Filling}, let $\psi$ be an algebraic modular form on $X_0$ of level $K_{\mathsf{S}}$, let
\m{u\in\Uni{\mathcal{Y}}}, \m{a\in G(X_0)} and \m{h\in G(V)} with finite component in~$K^{\mathrm{pr}}_{\mathsf{S}}$.
Then
\[\mathrm{P}_{\mathsf{S}}\bigl(\mathrm{cc}_{\mathsf{S}}(\psi),\Besr{\mathcal{Y}}\bigr)\bigl(u\,\incl{\mathcal{Y}}(a)\,h\bigr)
=\psi_E\Bigl(-\frac{\sigma(u)}{g_{\mathsf{S}}\varpi_{\mathsf{S}}}\Bigr)\cdot\Bigl\langle\mathsf{A}\bigl(a_\infty^{-1}\psi(a_f\eta)\bigr),\,W(h_\infty)\Bigr\rangle ,\]
where $\sigma(u)$ is the sum of the entries of $u$ just above the diagonal along the frame of the shape, in the
coordinates of the filling data, $a_\infty^{-1}$ acts on the Weyl coordinates of the weight, and
$\langle\,,\rangle$ is the bilinear dot product of coordinate vectors.
\end{proposition}
\begin{proof}
Write \m{a_f\eta=\gamma\,\mathsf{g}\,k} with $\gamma$ a rational point of~$X_0$, $\mathsf{g}$ the chosen
representative of the class of \m{a_f\eta} and \m{k\in K_{\mathsf{S}}}. \emph{First}, if a point $g$ has finite
component, translated by $\mathsf{g}$ and~$\eta$, in the principal level, then
$\mathrm{P}_{\mathsf{S}}(\mathrm{cc}_{\mathsf{S}}(\psi),\Besr{\mathcal{Y}})(g)=\langle\mathsf{A}(\psi(\mathsf{g})),W(g_\infty)\rangle$:
a term \m{\mathrm{cb}_{a'}(\gamma'g)} with \m{\gamma'\in\Besr{\mathcal{Y}}} can be non-zero only if $a'$ belongs to
the class of $\mathsf{g}$ and $\gamma'$ comes from the group of the rational unipotent elements that this class
conjugates into the principal level; so the sum over the classes reduces to this class, and over that group the
cut-offs sum to~$1$ (Inputs~\ref{in:units}(d), facts (1) and~(2)). \emph{Then}, by strong approximation for the
unipotent group (Inputs~\ref{in:units}(d), fact~(3)), there is a rational unipotent $\varsigma$ such that the finite
part of \m{\varsigma u}, conjugated by~$a_f\eta$, lies in the principal level. By Lemma~\ref{l:poinbasic}(i) the
series is unchanged if its argument is multiplied on the left by
$\varsigma$ and by \m{\incl{\mathcal{Y}}(\gamma^{-1})}, and the first step applies to the new point. The Bessel
vector transforms under the real points of $X_0$ by the Weyl representation, which $\mathsf{A}$ intertwines, and
under the real unipotent part by the character of~$\sigma$; the level of $\psi$ absorbs~$k$; and for a unipotent
element whose finite part lies in the principal level the value of the character at the real part equals
$\psi_E(-\sigma/(g_{\mathsf{S}}\varpi_{\mathsf{S}}))$ of the whole element. Since $\sigma$ is additive on products
of unipotent elements and vanishes on rational points under $\psi_E$, the factor is the one stated. The
remaining level condition holds because $K_{\mathsf{S}}$ is the preimage of the principal level and $h_f$ lies
in~it.
\end{proof}
\leanline{\lean{poincarebesselnormalform}}

\begin{proposition}[the Poincar\'e series is generic]\label{p:genericpoin}
In Setting~\ref{s:filling}, with Weil data for~$E$ and $\mathsf{S}$ admissible, let $\psi$ be an algebraic modular
form on $X_0$ of level $K_{\mathsf{S}}$, and assume
$\mathbf{j}\cdot\cst{\mathcal{Y}}=-(g_{\mathsf{S}}\varpi_{\mathsf{S}})^{-1}$. Then
$\mathrm{P}_{\mathsf{S}}(\mathrm{cc}_{\mathsf{S}}(\psi),\Besr{\mathcal{Y}})$ is generic for~$\mathcal{Y}$.
\end{proposition}
\begin{proof}
Left invariance under $\Besr{\mathcal{Y}}$ is Lemma~\ref{l:poinbasic}(i). For every point $g$ there is a continuous
function $F_g$ with
$\mathrm{P}(u\,h)=\psi_E(-\sigma(u)/(g_{\mathsf{S}}\varpi_{\mathsf{S}}))\,F_g(h)$ for all \m{u\in\Uni{\mathcal{Y}}} and
all $h$ whose finite part differs from that of $g$ by an element of the principal level: if the series vanishes at
all such points, take \m{F_g=0}; otherwise Lemma~\ref{l:poinbasic}(ii) moves one of these points into the principal
level by an element of $\Besf{\mathcal{Y}}$, which is a product of a unipotent element and an element of the image
of $\incl{\mathcal{Y}}$, and Proposition~\ref{p:normal} gives $F_g$ explicitly, a continuous function of the real
component. The sets of such $h$ are open and cover, and \m{F_g=\mathrm{P}} there (take \m{u=1}); so the series is
continuous. For $u$ as in the definition of genericity,
$\mathrm{P}(ug)=\psi_E(-\sigma(u)/(g_{\mathsf{S}}\varpi_{\mathsf{S}}))\mathrm{P}(g)$. Finally
$\mathrm{sh}_{\mathcal{Y}}=\cst{\mathcal{Y}}\cdot\mathsf{C}$ for the matrix $\mathsf{C}$ of the shape with
$\operatorname{tr}(\mathsf{C}u)=\sigma(u)+\cj{\sigma(u^{-1})}$ (an entrywise computation, by cases on the shape).
Moreover \m{\sigma(u^{-1})=-\sigma(u)}, because $\sigma$ is additive on products of unipotent elements; so
\m{\operatorname{tr}(\mathrm{sh}_{\mathcal{Y}}u)=\cst{\mathcal{Y}}\,(\sigma(u)-\cj{\sigma(u)})}. Put
\m{\varkappa=-(g_{\mathsf{S}}\varpi_{\mathsf{S}})^{-1}}; by the hypothesis \m{\varkappa=\mathbf{j}\,\cst{\mathcal{Y}}},
and \m{\cj\varkappa=-\varkappa} because \m{\cj{\mathbf{j}}=-\mathbf{j}} and \m{\cj{\cst{\mathcal{Y}}}=\cst{\mathcal{Y}}}.
Hence
\[\tfrac{\mathbf{j}}{2}\operatorname{tr}(\mathrm{sh}_{\mathcal{Y}}u)=\tfrac12\bigl(\varkappa\,\sigma(u)-\varkappa\,\cj{\sigma(u)}\bigr)
=\tfrac12\bigl(\varkappa\,\sigma(u)+\cj{\varkappa\,\sigma(u)}\bigr),\]
and \m{\psi_E(\tfrac12(x+\cj x))=\psi_E(x)} for every adele~$x$, the standard character being invariant under
conjugation. So the character of the normal form, \m{\psi_E(\varkappa\,\sigma(u))}, is
$\psi_E\bigl(\tfrac{\mathbf{j}}{2}\operatorname{tr}(\mathrm{sh}_{\mathcal{Y}}u)\bigr)$.
\end{proof}
\leanline{\lean{isGenericpoincare}}

\sect{Stage functions}{s:stages}

\begin{lemma}\label{l:stage}
In Setting~\ref{s:data} let $\vartheta$ be a tame character of~$D$, \m{\varphi\in\Col{D}}, and $\xi$, $y$ tuples of
real and of finite points as in Definition~\ref{d:stage}.
\begin{enumerate}[label=(\roman*)]
\item Let $i$ and $j$ be natural numbers with \m{j\le i} and \m{j\le2\rhoD{D}+2}, and let $\xi'$, $y'$ be obtained
  from $\xi$, $y$ by replacing the entries of index $i$ by a real point $\xi'_i$ and a finite point $y'_i$.
  Then $\unf{\vartheta}{\varphi}{\xi'}{y'}{j}{\mathcal{Y}}=\unf{\vartheta}{\varphi}{\xi}{y}{j}{\mathcal{Y}}$ for every
  grading~$\mathcal{Y}$ of $\Ys{D}{j}$.
\item For every natural number $j$ with \m{j\le2\rhoD{D}+1}, the function
  $\unf{\vartheta}{\varphi}{\xi}{y}{j}{\gY{j}}$ is generic for~$\gY{j}$.
\item If $y^{\circ}$ is the tuple of the chosen representatives of the cosets $y_j\,\levelY{\vartheta}{j}$, then
  $\evD{\vartheta}{y^{\circ}}{i}=\evD{\vartheta}{y}{i}$ for every natural number~$i$.
\end{enumerate}
\end{lemma}
\begin{proof}
(i) By the recursion, the function of stage $j$ uses only the points of index \m{<j}. (ii) By the recursion: for
\m{j=0} it is Proposition~\ref{p:genericpoin} for the filling data of~$D$ (Proposition~\ref{p:fits}), whose
hypothesis is the definition of the constant of~$\gY{0}$; the step is Inputs~\ref{in:lib}(a), all gradings $\gY{j}$
having the same constant. (iii) The representative is \m{y_j\,v_j} with $v_j\in\levelY{\vartheta}{j}$, so
$\tUps{\vartheta}{y^{\circ}}{j}=\pW{\nu}(y_j,y_{j+1})\,\pW{\nu}(v_j,v_{j+1})\ups{\vartheta}{j}$, and
$\pW{\nu}(v_j,v_{j+1})\ups{\vartheta}{j}=\ups{\vartheta}{j}$ by Lemma~\ref{l:upsilon}: the reductions of $v_j$ and
$v_{j+1}$ are~$1$, so its hypothesis holds and the character factor is~$1$.
\end{proof}
\leanline{\lean{unfoldStageupdate}\sep\lean{isGenericunfoldStage}\sep\lean{evaluatedDatumout}}

The next statement is the bridge to the definite side. It uses two results of the note \emph{Two consecutive theta
lifts}, stated for an \emph{arrow pair} $\mathcal{R}$ of three gradings $\mathcal{Y}\to\mathcal{Y}'\to\mathcal{Y}''$
with sets $X$, $X'$, $X_0''$: such a pair has a finite set $X''$, a homomorphism
$\mathrm{emb}_{\mathcal{R}}\colon G(X_0'')\to G(X'')$ and a finite evaluation. \emph{The double unfolding}: for a
function $F$ generic for~$\mathcal{Y}$, the iterated integral over \m{[\Bes{\mathcal{Y}'}]} and
\m{[\Bes{\mathcal{Y}}]} of $F$ times the two constant terms at the string maps is a constant times the integral
over \m{[G(X)]\times[G(X')]} of $F$ on the image of \m{\incl{\mathcal{Y}}} times the theta function of the zero
term of the arrow product of the two data. \emph{The double lift}: such an integral of a coefficient of an
algebraic modular form $\varphi$ against the theta kernel of an arrow tensor is a coefficient of the composite of
the two definite lifts of~$\varphi$.

\begin{proposition}[even stages are coefficients of definite chains]\label{p:even}
In Setting~\ref{s:data} let \m{\xi=(\xi_j)} be a tuple of real points as in Definition~\ref{d:stage}. Let
$i$ be a natural number with \m{2i<2\rhoD{D}+1}, let $\mathcal{Y}$ be a grading of $\Ys{D}{2i+2}$ with norm
$\ND{D}$, a finite set $X$ and a string of length~$l$, let $\mathcal{R}$ be an arrow pair, with a finite set~$Q_0$
of hyperbolic planes, of the gradings $\gY{2i}\to\gY{2i+1}\to\mathcal{Y}$ whose finite evaluation is
\m{\pairEvalY{i}}, and let $\xi^{\star}\in G_\infty(\Ys{D}{2i+2})$. Then there is a linear functional
$\mathsf{t}$ such that for every tame character $\vartheta$ of~$D$, every column \m{\varphi\in\Col{D}}, every tuple
$y$ of finite points with \m{y_0\in\Kch{D}{\Ys{D}{0}}} and every \m{g\in G(X)}:
\[\unf{\vartheta}{\varphi}{\xi}{y}{2i+2}{\mathcal{Y}}\bigl(\incl{\mathcal{Y}}(g)\cdot(\xi^{\star},y_{2i+2})\bigr)
=\coef{\mathsf{t}}{\chnn{D}{\vartheta}{\evD{\vartheta}{y}{\cdot}}{i+1}(\varphi)}\bigl(\mathrm{emb}_{\mathcal{R}}(g)\bigr).\]
\end{proposition}
\begin{proof}
By strong induction on~$i$, for all $\mathcal{Y}$, $\mathcal{R}$, $\xi^{\star}$. \emph{Claim}: there is a
functional $\mathsf{t}_0$ such that the function of stage $2i$, on the image of $\incl{\gY{2i}}$ translated by
\m{(\xi_{2i},y_{2i})}, is the coefficient at $\mathsf{t}_0$ of the chain of $i$ pairs. For \m{i=0} the stage function
is the Poincar\'e series, and Proposition~\ref{p:normal} with \m{u=1} (so that the character is~$1$) and the
hypothesis on $y_0$ presents it as $\langle\mathsf{A}(a_\infty^{-1}\varphi(a_f)),W(\xi_0)\rangle$, a linear
functional of the value of $\varphi$; the central translation in \m{\colco{D}} cancels~$\eta$. For \m{i>0} the claim
is the induction hypothesis for \m{i-1}, with the grading $\gY{2i}$ and the middle arrow pair. Now unfold twice:
by Definition~\ref{d:stage} the function of stage \m{2i+2} at the given point is an iterated integral over
$[\Bes{\gY{2i+1}}]$ and $[\Bes{\gY{2i}}]$ of the function of stage~$2i$, generic by Lemma~\ref{l:stage}(ii), times
two constant terms; the real parts of the Weil operators are evaluated at the fixed real points and absorbed in the
Fock vectors. By the double unfolding this is a constant times an integral over
$[G(\Xs{D}{2i})]\times[G(\Xs{D}{2i+1})]$ of the stage function, replaced by its coefficient from the claim, times
the theta function of the zero term of the arrow product $\tUps{\vartheta}{y}{2i}\boxtimes\tUps{\vartheta}{y}{2i+1}$
translated by the stacked string map. The zero term of an integrated datum is an integrated datum whose finite
part is the finite zero term, and by the defining property of the finite evaluation and the hypothesis on
$\mathcal{R}$ this finite part is the arrow tensor of $\evD{\vartheta}{y}{i}$. The double lift, applied to the
lifts \m{2i+1} and \m{2i+2} (theta lifts in the sense required, by Inputs~\ref{in:lib}(h)), turns the integral
into a coefficient of the composite of the two lifts with datum $\evD{\vartheta}{y}{i}$
applied to the chain of $i$ pairs. The functional $\mathsf{t}$ so obtained depends on $\xi$, $\xi^{\star}$,
$\mathcal{R}$ only.
\end{proof}
\leanline{\lean{unfoldStageeven}}

\sect{Lattices along a string}{s:lattices}

Throughout this section, $\ord_w$ of an idele is as in Definition~\ref{d:space}; the formal text works with
\m{-\ord_w}, the logarithm of the valuation.

\begin{lemma}[adapted bases]\label{l:adapted}
Let $\mathcal{Y}$ be a grading of a hermitian space with frame~$V$, with index set~$I$, a finite set $X$ and a
string of length~$l$.
\begin{enumerate}[label=(\roman*)]
\item Let $e\colon X\sqcup\{0,\dots,l-1\}\to I$ be a bijection and $A_1$, $A_2$ square matrices over a commutative
  ring $R$, indexed by~$I$ and triangular for the flag of~$e$. Then for every string index $a$, with
  \m{i=e(a)}: \m{(A_1A_2)_{ii}=(A_1)_{ii}(A_2)_{ii}}.
\item For every \m{y\in G_f(V)} there is a pair $(\mathbf{b},\beta)$ that is $\adat{\mathcal{Y}}{y}$.
\item Let \m{y\in G_f(V)} and let $(\mathbf{b},\beta)$ be $\adat{\mathcal{Y}}{y}$. Put
  $\mathsf{k}=(\mathbf{b}^{-1}\,x_{V,f}\,y)^{-1}$. For string indices $a$, $a'$ with \m{a'\le l-1-a}, the entry of
  \m{\mathsf{k}^{*}\mathsf{k}} at the string indices $(a,a')$ is $J_V(z_a,z_{a'})\cdot\cj{\beta_a}\beta_{a'}$.
\item If $\adb{\mathcal{Y}}{h}{\mathbf{b}}{\beta}$ for some matrix \m{h\in\GL_I(\Af)}, then the inverses satisfy the same
  relation with themselves: $\adb{\mathcal{Y}}{\mathbf{b}^{-1}}{\mathbf{b}^{-1}}{\beta^{-1}}$.
\item Let \m{h,h'\in\GL_I(\Af)} with $\adb{\mathcal{Y}}{h}{\mathbf{b}}{\beta}$ and $\adb{\mathcal{Y}}{h'}{\mathbf{b}'}{\beta'}$, and let
  $w$ be a finite place at which the matrix \m{h^{-1}h'} is integral. Then
  \m{\ord_w\beta'_a\ge\ord_w\beta_a} for every string index~$a$.
\item In the situation of (v), if instead $h$ and $h'$ have the same component at~$w$, then
  \m{\ord_w\beta'_a=\ord_w\beta_a} for every string index~$a$.
\end{enumerate}
\end{lemma}
\begin{proof}
(i) A string index is alone in its class of the flag order, since the numbers \m{4a+1} are distinct and different
from~$2l$; in the sum \m{\sum_{i'}(A_1)_{ii'}(A_2)_{i'i}} triangularity leaves only \m{i'=i}. (ii) Hermite
reduction over~$\Oh$ (Inputs~\ref{in:lib}(d)), applied to the inverse of \m{x_{V,f}\,y}, gives an invertible
$\mathbf{b}$, triangular for the flag, such that \m{(x_{V,f}\,y)^{-1}\mathbf{b}} is integral with integral inverse. The inverse
of a matrix that is triangular for the flag is triangular for the flag, so by~(i), applied to \m{\mathbf{b}\,\mathbf{b}^{-1}=1}, the
diagonal entries of $\mathbf{b}$ at the string indices are units of~$\Af$; let $\beta$ be this diagonal.
(iii) Since $x_{V,f}^{*}J_V\,x_{V,f}=1$ and $y$ is unitary,
\m{\mathsf{k}^{*}\mathsf{k}=\mathbf{b}^{*}J_V\mathbf{b}}. In the flag order the string vectors of grade \m{\ge0} come first, by
decreasing grade, then~$X$, then the other string vectors. By triangularity the column of $\mathbf{b}$ at a string index
$a'$ is \m{\beta_{a'}z_{a'}} plus a combination of basis vectors that precede $z_{a'}$ in the flag order, and $J_V$
pairs a string vector $z_{a''}$ only with its reverse. For \m{a'\le l-1-a} a comparison of flag orders shows that the
only term of \m{\mathbf{b}^{*}J_V\mathbf{b}} at $(a,a')$ is the product of the two diagonal entries with \m{J_V(z_a,z_{a'})}.
(iv) The inverse of a triangular matrix is triangular; \m{(\mathbf{b}^{-1})^{-1}\mathbf{b}^{-1}=1}; and by~(i) the string diagonal of
$\mathbf{b}^{-1}$ is~$\beta^{-1}$. (v) $\mathbf{b}^{-1}\mathbf{b}'=(\mathbf{b}^{-1}h)(h^{-1}h')(h'^{-1}\mathbf{b}')$ is integral at~$w$ and triangular, and by (i)
and~(iv) its entry at the $a$-th string index is \m{\beta_a^{-1}\beta'_a}. (vi) Apply (v) twice.
\end{proof}
\leanline{\lean{mulapplystringofblockTriangular}\sep\lean{existsisAdaptedAt}\sep\lean{gramstringentry}\sep
\lean{inv}\sep\lean{ideleLogValle}\sep\lean{ideleLogValeq}}

\begin{proposition}[lattices with the same lines lie in one Bessel orbit]\label{p:orbit}
Let $E$ be stably unramified, $\mathcal{Y}$ a grading of a hermitian space with frame $V$ with a finite set $X$ and
a string of length~$l$, and let \m{h,h'\in G_f(V)}. Let $(\mathbf{b},\beta)$ be $\adat{\mathcal{Y}}{h}$ and $(\mathbf{b}',\beta')$ be
$\adat{\mathcal{Y}}{h'}$, and assume \m{\ord_w\beta'_a=\ord_w\beta_a} for every finite place $w$ and every string
index $a$ with \m{a\le l-1-a}. Then there is \m{k\in\Kz{V}} with \m{h'\,k\,h^{-1}\in\Besf{\mathcal{Y}}}.
\end{proposition}
\begin{proof}
Multiplying $\mathbf{b}'$ on the right by the diagonal matrix that is $1$ on $X$ and \m{\beta'^{-1}_a\gamma_a} on the string,
with \m{\gamma_a=\beta_a} for \m{a\le l-1-a} and \m{\gamma_a=\beta'_a} otherwise, gives an adapted basis at $h'$ with
string diagonal~$\gamma$: the factor is integral with integral inverse by the hypothesis. So we may assume
\m{\beta'_a=\beta_a} on the first half \m{a\le l-1-a}. Let $\mathsf{k}=(\mathbf{b}^{-1}x_{V,f}h)^{-1}$ and
$\mathsf{k}'$ likewise; both lie in \m{\GL_I(\Oh)}. The columns of $\mathsf{k}$ at the string indices of the first
half form a matrix with an integral left inverse, and so do those of~$\mathsf{k}'$. By Lemma~\ref{l:adapted}(iii)
the two matrices have the same Gram matrix. Witt's theorem over $\Oh$ (Inputs~\ref{in:lib}(d)) gives \m{k\in\Kz{V}}
carrying the first to the second. In rational coordinates this says that \m{h'kh^{-1}}, applied to the column of
$\mathbf{b}$ at a first-half string index, gives the column of $\mathbf{b}'$ there. By the description of the Bessel group it remains
to see that \m{h'kh^{-1}-1}, in rational coordinates, has zero entry $(i,i')$ whenever $i'$ is a string index of
grade \m{\ge0} and \m{\gr(i)\le\gr(i')}; the string indices of grade \m{\ge0} are those of the first half. Writing
the matrix as its product with \m{\mathbf{b}\,\mathbf{b}^{-1}}, the entry is a sum over the columns of $\mathbf{b}$ at first-half string
indices of smaller or equal number (by triangularity of $\mathbf{b}^{-1}$), on which the matrix acts as \m{\mathbf{b}\mapsto \mathbf{b}'}; the
columns of $\mathbf{b}$ and $\mathbf{b}'$ at such an index $a''$ are \m{\beta_{a''}z_{a''}} plus string vectors of higher grade, with the
same diagonal. A check of the cases gives zero.
\end{proof}
\leanline{\lean{existsmembesselFin}}

For the next three statements recall from the note \emph{Two consecutive theta lifts}: for finite ideles $\mathfrak{a}$,
$\pi$ and a finite adele $u$ of norm one, a grading $\mathcal{Y}$ of $V$ as \m{X\sqcup\{\mathrm{pt}\}} plus a string of
length~$l$ \emph{spans the lattice of} \m{(\mathfrak{a},\pi,u)} if \m{x_{V,f}=\Lambda\,k} with
\m{k\in\GL_I(\Oh)} and $\Lambda$ the \emph{graded lattice basis}: for even $l$ the diagonal matrix with $1$ on
\m{X\sqcup\{\mathrm{pt}\}}, \m{u\,\mathfrak{a}^{-1}} at the string indices $a$ with \m{2a<l} and \m{u\,\cj{\mathfrak{a}}}
at the others; for odd $l$ the matrix with $1$ on~$X$, $u$ on the string, and the column of the point equal to
\m{\cj\pi^{-1}} at the point plus \m{u(\pi\cj\pi)^{-1}} at the middle string vector. And $\mathfrak{N}$ is \emph{split and
squarefree along} a finite set $\Sigma$ of finite places with the idele~$\pi$ if \m{\mathfrak{N}=\pi\cj\pi}, $\pi$ is a
uniformiser at the places of $\Sigma$ and a unit elsewhere, and $\Sigma$ is disjoint from its conjugate.

\begin{lemma}[the standard lattice has near lines]\label{l:nearzero}
Let $\mathcal{Y}$ be a grading of $V$ as \m{X\sqcup\{\mathrm{pt}\}} plus a string of length~$l$, let $\mathfrak{a}$,
$\pi$ be finite ideles and $n$ a natural number such that $\mathcal{Y}$ spans the lattice of
$(\mathfrak{a},\pi,(\mathfrak{a}/\cj{\mathfrak{a}})^{n})$, let $K$ be a subgroup of $\Kz{V}$, $Q$ a finite set of
finite places and \m{h\in G_f(V)} such that \m{\varsigma h\in K} for some \m{\varsigma\in\Besf{\mathcal{Y}}}. Then
$\near{\mathcal{Y}}{\mathfrak{a}}{Q}{n}{0}{h}$.
\end{lemma}
\begin{proof}
Let \m{x_{V,f}=\Lambda k}. In rational coordinates \m{\varsigma^{-1}} is triangular for the flag with string
diagonal~$1$, by the description of the Bessel group, and so is $\Lambda$, with string diagonal
$(\mathfrak{a}/\cj{\mathfrak{a}})^{n}\epsilon_a$. Their product $\mathbf{b}$ satisfies
$\mathbf{b}^{-1}x_{V,f}h=k\cdot(\varsigma h)\in\GL_I(\Oh)$ in the coordinates of the frame, and by
Lemma~\ref{l:adapted}(i) its string diagonal is that of~$\Lambda$. So all the differences of orders are~$0$.
\end{proof}
\leanline{\lean{hasNearLineszero}}

\begin{lemma}\label{l:inclK}
Let $\mathcal{Y}$ be a grading of $V$ as \m{X\sqcup\{\mathrm{pt}\}} plus a string of length~$l$ with norm~$\mathfrak{N}$ that
spans the lattice of \m{(\mathfrak{a},\pi,u)}, let $\mathfrak{N}$ be split and squarefree along $\Sigma$ with the idele~$\pi$, and
let \m{h\in G_f(X\sqcup\{\mathrm{pt}\})} be integral with integral inverse and with the row of the point congruent to
\m{(0,\dots,0,1)} modulo~$\pi$. Then $\inclf{\mathcal{Y}}(h)\in\Kz{V}$.
\end{lemma}
\begin{proof}
In the coordinates of the frame \m{x_{V,f}=\Lambda k},
$\inclf{\mathcal{Y}}(h)-1=k^{-1}\Lambda^{-1}(\tilde h-1)\Lambda k$ with $\tilde h$ the block matrix of $h$ and~$1$.
For even $l$ the rows and columns \m{X\sqcup\{\mathrm{pt}\}} of $\Lambda$ are those of the identity, and \m{h-1} is
integral. For odd $l$, \m{\Lambda^{-1}(\tilde h-1)\Lambda} has the entries of \m{h-1}, except that the column of the
point is divided by~$\cj\pi$ and the row of the point is multiplied by~$\cj\pi$ (so that the entry at the point
itself is unchanged); moreover the row of the point of \m{h-1}, divided by~$-\pi$, stands in the row of the middle
string vector, whose entry in the column of the point is the entry of \m{h-1} at the point divided by~$-\mathfrak{N}$. So
three congruences are needed. The row of the point of \m{h-1} is $\equiv0$ modulo~$\pi$ by hypothesis. The
integral matrices with integral inverse and with this property form a group, and $h$ is unitary; so
\m{h^{-1}=h^{*}} has the property too, which says that the column of the point of \m{h-1} is $\equiv0$
modulo~$\cj\pi$. In particular the entry of \m{h-1} at the point is $\equiv0$ modulo $\pi$ and modulo~$\cj\pi$,
hence modulo \m{\mathfrak{N}=\pi\cj\pi}, as $\pi$ is a uniformiser on~$\Sigma$, $\cj\pi$ on~$\cj\Sigma$, both are units
elsewhere and $\Sigma$ is disjoint from~$\cj\Sigma$. So the matrix is integral, and conjugating by
\m{k\in\GL_I(\Oh)} keeps integrality.
\end{proof}
\leanline{\lean{inclXFinmemKf}}

\begin{proposition}[the two halves of a string diagonal]\label{p:rev}
Let $E$ be stably unramified, $\mathcal{Y}$ a grading of $V$ with a finite set $X$, a string of length $l$ and
norm~$\mathfrak{N}$, let \m{h\in G_f(V)} and let $(\mathbf{b},\beta)$ be $\adat{\mathcal{Y}}{h}$. Let $\mathfrak{N}$ be split and squarefree along
$\Sigma$ with the idele~$\pi$. Assume: if $l$ is odd, then for no \m{\varsigma\in\Besf{\mathcal{Y}}} and no
\m{w\in\Sigma} does $(\varsigma h)^{-1}\,\inclf{\mathcal{Y}}(\Kw{X}{w})\,(\varsigma h)\subseteq\Kw{V}{w}$ hold. Then for
every string index $a$ and every finite place~$w$: $\ord_w\beta_a+\ord_{\cj w}\beta_{l-1-a}\ge0$.
\end{proposition}
\begin{proof}
For any adapted pair, Lemma~\ref{l:adapted}(iii) with \m{a'=l-1-a} shows that the number
$\cj{\beta_a}\beta_{l-1-a}J_V(z_a,z_{l-1-a})$ is an entry of the Gram matrix of an element of \m{\GL_I(\Oh)}, hence
integral; and \m{J_V(z_a,z_{l-1-a})} is $1$, or \m{-\mathfrak{N}} at the middle index. This is the claim, except when $l$ is
odd, $a$ is the middle index and $w$ or $\cj w$ lies in~$\Sigma$, where it gives only
$\ord_w\beta_a+\ord_{\cj w}\beta_a\ge-1$. For the middle index the claim is symmetric in $w$ and~$\cj w$, so let
\m{w\in\Sigma}; then \m{\cj w\ne w}. Suppose the claim fails, so that the sum is~$-1$. Let $\gamma$ be the family with \m{\gamma_{a'}=\beta_{a'}}
up to the middle index and $\gamma_{a'}=\cj{\beta_{l-1-a'}}^{-1}$ beyond, and $\mathsf{d}$ the diagonal matrix with
$1$ on $X$ and $\gamma$ on the string; it is an adapted basis of itself. Since \m{\cj w\ne w}, there is a finite
point $y$, equal to $1$ outside $w$ and~$\cj w$, such that \m{x_{V,f}\,y\,h} and $\mathsf{d}$ have the same
component at~$w$. Choose a pair $(\mathbf{b}',\beta')$ adapted at \m{yh} (Lemma~\ref{l:adapted}(ii)). By
Lemma~\ref{l:adapted}(vi) its diagonal has the orders of $\gamma$ at~$w$ and those of $\beta$ away from $w$
and~$\cj w$. At $\cj w$ its orders are at most those of~$\gamma$, by Lemma~\ref{l:adapted}(v) for the matrices
\m{x_{V,f}\,yh} and~$\mathsf{d}$. Indeed, the matrix that compares the two lattices is
\m{(x_{V,f}\,yh)^{-1}\mathsf{d}=(x_{V,f}\,yh)^{*}J_V\,\mathsf{d}}, and since \m{x_{V,f}\,yh}
agrees with $\mathsf{d}$ at~$w$, its component at~$\cj w$ is that of \m{\mathsf{d}^{*}J_V\mathsf{d}}. This matrix
is the identity on~$X$, and its only non-zero entries on the string are
\m{\cj{\gamma_{a'}}\gamma_{l-1-a'}=1} off the middle index and \m{-\mathfrak{N}\,\cj{\gamma_a}\gamma_a} at the middle
index~$a$, whose order at~$\cj w$ is \m{1+\ord_w\beta_a+\ord_{\cj w}\beta_a=0} by the assumed failure. So the
comparison matrix is integral at~$\cj w$. And the orders at $\cj w$ are at least those of $\beta$ on the first
half, by the
Gram condition of the first step for $(\mathbf{b}',\beta')$ and, at the middle index, by the failure of the inequality. As
\m{\gamma=\beta} on the first half, the two diagonals have the same orders there at every place. So
Proposition~\ref{p:orbit} applies to $h$ and \m{yh}: there is \m{k\in\Kz{V}} with
$\varsigma=yhkh^{-1}\in\Besf{\mathcal{Y}}$. For this $\varsigma$ and~$w$ the forbidden inclusion holds: for
\m{a_0\in\Kw{X}{w}}, the conjugate $(\varsigma h)^{-1}\inclf{\mathcal{Y}}(a_0)(\varsigma h)$ is $1$ outside $w$,
$\cj w$, and at $w$ it is $k^{-1}$ times the conjugate of the block matrix of $a_0$ and~$1$ by
\m{\mathsf{d}^{-1}x_{V,f}}, times~$k$; the lattice of $\mathsf{d}$ is stable under such block matrices, so
this is integral. This contradicts the assumption.
\end{proof}
\leanline{\lean{ideleLogValaddrevnonpos}}

\begin{proposition}[one lift moves the lines by one step]\label{p:step}
Let $E$ have Weil data, with idele~$\alpha$. Let $\mathcal{Y}$ be a grading of $V$ with a string of length~$l$ and
$\mathcal{Y}'$ a grading of $V'$ with a string of length~$l'$, with the same norm~$\mathfrak{N}$ and \m{l'+1=l}. Let
$\mathfrak{a}$ be a finite idele, $Q$ a finite set of finite places, $n$, $n'$, $k$ natural numbers with \m{n=n'} if
$l$ is odd and \m{n=n'+1} if $l$ is even, let \m{h\in G(V)}, \m{h'\in G(V')}, $\nu$ a character, $v$ a vector of the
Fock model and $\Phi\in\mathcal{S}(\Hom(V,V'))$, and $\mathfrak{b}$ a finite idele. Assume:
\begin{enumerate}[label=(\alph*)]
\item $\near{\mathcal{Y}}{\mathfrak{a}}{Q}{n}{k}{h_f}$;
\item \m{\cst{\mathcal{Y}}} and \m{\mathfrak{a}\cj{\mathfrak{a}}} have the same order at every finite place;
\item $\mathrm{ct}_{\mathcal{Y},\mathcal{Y}'}(\nu,v\tensor\Phi)(h,h')\ne0$;
\item $\ord_w\mathfrak{b}\le\ord_w\cj{(\alpha/\mathfrak{a})}+[w\in Q]$ for every finite place~$w$;
\item $\varrho_f(\mathfrak{b}\,y)\Phi=\Phi$ for every integral finite vector~$y$;
\item every pair $(\mathbf{b}',\beta')$ that is $\adat{\mathcal{Y}'}{h'_f}$ has
  $\ord_w\beta'_a+\ord_{\cj w}\beta'_{l'-1-a}\ge0$ for all string indices $a$ and all finite places~$w$.
\end{enumerate}
Then $\near{\mathcal{Y}'}{\mathfrak{a}}{Q}{n'}{k+1}{h'_f}$.
\end{proposition}
\begin{proof}
Let $(\mathbf{b},\beta)$ be the adapted pair of (a) and $(\mathbf{b}',\beta')$ any adapted pair at~$h'_f$. \emph{Claim}: for every
string index $a$ of $\mathcal{Y}'$,
\[\cst{\mathcal{Y}}^{\,\delta_a}\cdot\frac{\cj{\mathfrak{b}}}{\cj{\beta_{l'-a}}\;\alpha\;\beta'_a}\in\Oh,\qquad
\delta_a=1\ \text{if ($l'$ odd)}\Leftrightarrow(a\le l'-1-a),\quad\delta_a=0\ \text{otherwise.}\]
Granting the claim, take orders at every place: with (b) the claim bounds \m{\ord_w\beta'_a} from above by
\m{-\ord_{\cj w}\beta_{l'-a}} plus the order of \m{\cj{\mathfrak{b}}/\alpha} and of the power of
\m{\mathfrak{a}\cj{\mathfrak{a}}}; (a) controls \m{\beta_{l'-a}} within~$k$ on \m{Q\cup\cj Q} and exactly elsewhere;
(d) costs one more unit on~$Q$; and (f) turns the upper bound for one half of $\beta'$ into the lower bound for the
other half. Linear arithmetic gives the near-lines bound for $\beta'$ with $n'$ and \m{k+1}.

\emph{Proof of the claim.} By Inputs~\ref{in:lib}(g) there is a splitting of \m{\Hom(V,V')} along the maps of
positive degree for which the constant term at the string map \m{\Tm{\mathcal{Y}}{\mathcal{Y}'}} is the theta
function of the zero term of the Heisenberg translate by this map. By (c)
and the support of the zero term (note \emph{Two consecutive theta lifts}: if a
datum is fixed by the translations \m{\varrho_f(\mathsf{M}y)}, $y$ integral, and the finite zero term of its
translate by a vector $t$ does not vanish, then $t$ lies in \m{\mathsf{M}'\Oh+k^{\perp}}, where
\m{\mathsf{M}\mathsf{M}'^{*}=1} and $k$ is the isotropic frame of the splitting), applied to the datum moved by the
finite parts of \m{(h,h')}, which by (e) is fixed by the translations of \m{\mathfrak{b}} times the lattice moved by
\m{(h,h')}: there is an integral vector $\lambda_0$ such that \m{\Tm{\mathcal{Y}}{\mathcal{Y}'}} and
$\cj{\mathfrak{b}}^{-1}(h,h')\lambda_0$ have the same pairing with the isotropic frame. Reading this relation in
the frames of $V$ and $V'$, for the entries $(i',i)$ with \m{\gr(i')<\gr(i)}, expresses the entries of
\m{\cj{\mathfrak{b}}\,\Tm{\mathcal{Y}}{\mathcal{Y}'}} there through the matrix $\mathsf{W}=x_{V',f}h'_f\cdot\lambda_0\cdot(x_{V,f}h_f)^{-1}$.
Now $\mathbf{b}'^{-1}\,\mathsf{W}\,J_V^{-1}(\mathbf{b}^{-1})^{*}$ is integral, because $\mathbf{b}'$ and $\mathbf{b}$ are adapted and $\lambda_0$ is
integral. Its entry at the $a$-th string index of $\mathcal{Y}'$ and the string index \m{l'-a} of $\mathcal{Y}$ is a
double sum in which, by the triangularity of $\mathbf{b}'^{-1}$ and $\mathbf{b}^{-1}$, only one term survives: the product of the
diagonal entries \m{\beta'^{-1}_a} and \m{\cj{\beta_{l'-a}}^{-1}}, of \m{\cj{\mathfrak{b}}/\alpha}, of the multiple
of the string map at~$z_a$ (which is $1$, \m{\cst{\mathcal{Y}}} or \m{-\mathfrak{N}\cst{\mathcal{Y}}}, by
Definition~\ref{d:theta}) and of the entry of $J_V^{-1}$ at \m{(z_a,z_{l'-a})}, which is $1$, or \m{-\mathfrak{N}^{-1}} when
$z_a$ is the middle vector of~$\mathcal{Y}$. The multiple \m{-\mathfrak{N}\cst{\mathcal{Y}}} occurs exactly at the middle
vector, where the factor \m{-\mathfrak{N}^{-1}} cancels~\m{-\mathfrak{N}}; so the product of the last two factors is
\m{\cst{\mathcal{Y}}} or~$1$, and comparing the cases of Definition~\ref{d:theta} with the definition of
$\delta_a$ shows that it is \m{\cst{\mathcal{Y}}^{\delta_a}}. This is the claim.
\end{proof}
\leanline{\lean{step}}

\begin{lemma}[valuations of the data]\label{l:vals}
In Setting~\ref{s:data}:
\begin{enumerate}[label=(\roman*)]
\item the ideal of $\ED{D}$ generated by $\piGb{D}$ has absolute norm \m{\gD{D}^{\fD{D}}};
\item for every finite place~$w$, the idele $\cgD{D}$ of Inputs~\ref{in:units}(a) and $\piGb{D}$ have the same
  order at~$w$;
\item for every natural number $j$ with \m{j\le2\rhoD{D}+1} and every finite place~$w$, the constant
  \m{\cst{\gY{j}}} and the idele \m{\ideA{D}\,\cj{\ideA{D}}} have the same order at~$w$;
\item for every natural number $j$ with \m{j\le2\rhoD{D}+1}, the grading $\gY{j}$ spans the lattice of
  \m{(\ideA{D},\piId{D},\lsc{D}{j})}, $\piId{D}$ the principal idele of~$\piN{D}$;
\item for a tame character $\vartheta$ of $D$ and a hermitian space with frame $V$ over $\ED{D}$, with a finite
  index set: every element of \m{\Kch{D}{V}} that is $1$ at the places of $\Qs{\vartheta}$ is congruent to $1$
  modulo every chosen place and every place of~$\Qs{\vartheta}$.
\end{enumerate}
\end{lemma}
\begin{proof}
(i) The ideal is the product of the chosen places, and those above $g_i$ have norm $g_i$ and are $\fD{D}$ in
number. (ii) Both orders are the multiplicity of $w$ in the ideal of~$\piGb{D}$. (iii) By the self-duality of the
Weil data, \m{\mathbf{j}\alpha\cj\alpha} has at $w$ the order of the inverse different, which
\m{\vpi{P}{\mu}^{-1}} generates; so \m{(\mathbf{j}\,\gD{D}\,\vpi{P}{\mu})^{-1}} has the order of
\m{\alpha\cj\alpha/\gD{D}}, and
$\ideA{D}\cj{\ideA{D}}=\alpha\cj\alpha/(\cgD{D}\cj{\cgD{D}})$, where \m{\cgD{D}\cj{\cgD{D}}} has the order of
\m{\piGb{D}\cj{\piGb{D}}=\gD{D}} by~(ii). (iv) This is Inputs~\ref{in:units}(a), by the definition of the spaces
and of the gradings. (v) At a chosen place this is the definition of \m{\Kch{D}{V}}; at a place of $\Qs{\vartheta}$
an element with component $1$ reduces to~$1$.
\end{proof}
\leanline{\lean{absNormpiGBar}\sep\lean{vchosenIdele}\sep\lean{valuationshiftconsteqideleAnorm}\sep
\lean{spansLatticegradingY}\sep\lean{chosenLevelinftrivialOnle}}

\begin{proposition}[where a stage function lives]\label{p:nearstage}
In Setting~\ref{s:data} let $\vartheta$ be a tame character of~$D$, $\xi$ and $y$ tuples of real and of finite points
as in Definition~\ref{d:stage} with \m{y_0\in\Kch{D}{\Ys{D}{0}}}, let \m{\varphi\in\Colram{D}} be a ramified column,
$j$ a natural number with \m{j\le2\rhoD{D}+1}, and \m{g\in G(\Ys{D}{j})} a point with
$\unf{\vartheta}{\varphi}{\xi}{y}{j}{\gY{j}}(g)\ne0$. Then
\[\near{\gY{j}}{\ideA{D}}{\Qs{\vartheta}}{\lfloor(m-j+1)/2\rfloor}{j}{g_f}.\]
\end{proposition}
\begin{proof}
By induction on~$j$. For \m{j=0}, Lemma~\ref{l:poinbasic}(ii) moves $g_f$ into the principal level by an element of
the finite Bessel group, and Lemma~\ref{l:nearzero} applies, by Lemma~\ref{l:vals}(iv) with
$\lsc{D}{0}=(\ideA{D}/\cj{\ideA{D}})^{\lfloor(m+1)/2\rfloor}$. For the step from $j$ to \m{j+1}: the value at $g$ is an
integral over \m{\varsigma\in[\Bes{\gY{j}}]} of the stage function at \m{\varsigma\cdot(\xi_j,y_j)} times a constant
term, so both factors are non-zero for some~$\varsigma$. The induction hypothesis at that point is hypothesis (a) of
Proposition~\ref{p:step}, which we apply with \m{\mathfrak{a}=\ideA{D}}, \m{Q=\Qs{\vartheta}}, the datum
$\datY{\vartheta}{j}=v\tensor\ups{\vartheta}{j}$ and $\mathfrak{b}$ a generator of the product $\mathfrak{m}$ of the
chosen places and the places of~$\Qs{\vartheta}$; the parities of the string lengths match the exponents. (b) is
Lemma~\ref{l:vals}(iii). (d): the order of $\mathfrak{b}$ at $w$ is the multiplicity of $w$ in the ideal of
$\piGb{D}$ plus \m{[w\in\Qs{\vartheta}]}, and \m{\alpha/\ideA{D}=\cj{\cgD{D}}}. (e): \m{\ups{\vartheta}{j}} is
named on the vacuum by an indicator modulo~$\mathfrak{m}$; by Inputs~\ref{in:units}(e) with \m{g'=0}, \m{q=1}
and \m{\mathsf{e}=\mathfrak{b}} (so that \m{\mathsf{e}\,q\equiv0} modulo~$\mathfrak{m}$) it is a combination of
translates \m{\varrho_f(t)} of the vacuum with \m{\cj{\mathfrak{b}}t} integral at every place. Only this support
is used: the denominator that the Input gives for \m{g'=0} is~$0$ and says nothing about the coefficients. Now
\m{\varrho_f(\mathfrak{b}y)} fixes the vacuum for integral~$y$ and commutes with \m{\varrho_f(t)} up to the value
of the additive character at the pairing of \m{\mathfrak{b}y} and~$t$, which is~$1$ because \m{\cj{\mathfrak{b}}t}
is integral (Input~1(a) and~(e) of the note \emph{Two consecutive theta lifts}: the standard lattice is self-dual).
(f) follows from Proposition~\ref{p:rev} for $\gY{j+1}$, with \m{\Sigma=\PiD{D}} and the principal idele $\piId{D}$
of $\piN{D}$ (Inputs~\ref{in:units}(b)), once its hypothesis is verified.

So let the string of $\gY{j+1}$ be odd, that is \m{j+1=2i+2} for a natural number~$i$, and suppose that for some
\m{\varsigma\in\Besf{\gY{j+1}}} and \m{w\in\PiD{D}} the inclusion of Proposition~\ref{p:rev} holds at the point~$g_f$.
Write $\varsigma^{-1}=u\,\incl{\gY{j+1}}(a_0)$ with $u$ unipotent, and \m{\varphi=\psi-\zw{D}{w}\psi} with
\m{\psi\in\Col{D}} (Inputs~\ref{in:units}(b)). By Proposition~\ref{p:even} for the pair index~$i$, with the real
point of $g$ and after replacing the entry of $\xi$ and $y$ of index \m{2i+2} (Lemma~\ref{l:stage}(i),(iii)), the
function of stage \m{2i+2} on the middle grading is a coefficient of the chain of \m{i+1} pairs. The central
element of $\zw{D}{w}$ is a scalar of norm one supported at $w$ and~$\cj w$; by Inputs~\ref{in:units}(g) the chain
of pairs of \m{\zw{D}{w}\psi} is the right translate by this scalar of the chain of pairs of~$\psi$, the character
factors being trivial
because the characters are unramified above~$\ND{D}$. Moving $u$ out by genericity (a string of odd length makes
every unipotent element raise grades by~$2$), the stage function at $g$ becomes the difference of two values of
the same coefficient, at a point and at its translate by the scalar; by the assumed inclusion this translate differs
from the point by an element of $\levelY{\vartheta}{j+1}$ — a place of $\PiD{D}$ is neither chosen nor in
$\Qs{\vartheta}$ — and the evaluated data depend only on cosets modulo these levels (Lemma~\ref{l:stage}(iii)). The
two values agree and the stage function vanishes at~$g$, against the assumption.
\end{proof}
\leanline{\lean{hasNearLinesofunfoldStagenezero}}

\sect{Periods as finite sums}{s:cuts}

\begin{lemma}\label{l:mult}
Let $H$ be a group with subgroups $S$, $A$, $J$.
\begin{enumerate}[label=(\roman*)]
\item For a finite set $T$ of left cosets of~$J$, an element \m{p\in S} and a left coset $t$ of~$J$:
  $\cov{S}{A}{T}{p\,t}=\cov{S}{A}{T}{t}$.
\item Let $H$ be a measurable group with a left invariant measure, written \m{\operatorname{vol}}, let $J$ be a
  measurable subgroup and $T_0$ a non-empty finite set of left cosets of~$J$. Then the union of the cosets in $T_0$
  has measure \m{|T_0|\cdot\operatorname{vol}(J)}.
\end{enumerate}
\end{lemma}
\begin{proof}
(i) Left multiplication by $p^{-1}$ is a bijection of the cosets of \m{S\cap A} in $S$ that carries the set counted
for \m{p\,t} to the set counted for~$t$. (ii) The union is disjoint and every coset is a left translate of~$J$.
\end{proof}
\leanline{\lean{coverMultiplicitysmul}\sep\lean{measureunioncosets}}

\begin{proposition}[a cut with multiplicities of prime-power type]\label{p:cut}
Let $V$ be a hermitian space with frame, let $S$, $A$, $J$, $H_0$ be subgroups of $G_f(V)$ with $A$ and $J$ levels,
\m{J\subseteq A\subseteq H_0}, let $\mathsf{K}$ be a compact subset of $G_f(V)$, $q$ a natural number and
$\mathsf{W}$ a set of finite points. Assume: every \m{h\in\mathsf{K}} lies in $H_0$ and has a natural number $n$
with \m{(huh^{-1})^{q^{n}}\in A} for all \m{u\in J}; and for every \m{h\in\mathsf{W}} there is \m{\varsigma\in S} with
\m{\varsigma h\in\mathsf{K}}. Then there is a finite set $T$ of left cosets of $J$ such that
\begin{enumerate}[label=(\roman*)]
\item for every \m{p\in S\cap A} and every coset~$t$: \m{t\in T} if and only if \m{p\,t\in T};
\item every $h$ with \m{hJ\in T} lies in~$H_0$;
\item for every $h$ with \m{hJ\in T} or \m{h\in\mathsf{W}}, the multiplicity $\cov{S}{A}{T}{hJ}$ divides a power
  of~$q$.
\end{enumerate}
\end{proposition}
\begin{proof}
Let $\mathsf{C}$ be a set of cosets \m{hJ}, \m{h\in\mathsf{K}}, that meets every $S$-orbit of such cosets exactly
once, and let $T$ be the set of the cosets \m{p\,t} with \m{p\in S\cap A}, \m{t\in\mathsf{C}}. It is contained in the
image of the compact set \m{A\mathsf{K}} in the discrete space of cosets of the open subgroup~$J$, hence finite.
(i) holds by construction. (ii): if \m{hJ=p\,h_0J} with \m{h_0\in\mathsf{K}}, then \m{h\in A\,H_0\,J\subseteq H_0}.
(iii): a coset \m{hJ} in $T$, and by the second hypothesis also the coset of a point of~$\mathsf{W}$, lies in the
$S$-orbit of some \m{h_0J\in\mathsf{C}}, \m{h_0\in\mathsf{K}}; by Lemma~\ref{l:mult}(i) it has the multiplicity
of~$h_0J$. Let $S_0$ be the stabiliser of $h_0J$ in~$S$. Since $\mathsf{C}$ meets the orbit once, the elements
\m{\varsigma\in S} with \m{\varsigma^{-1}h_0J\in T} are those of \m{S_0\,(S\cap A)}, so the multiplicity is the index
of \m{S_0\cap A} in~$S_0$, which is finite because \m{h_0Ah_0^{-1}} is commensurable with~$A$. This index divides
the order of the quotient of $S_0$ by the normal core of \m{S_0\cap A}. Every element of $S_0$ lies in
\m{h_0Jh_0^{-1}}, so by the first hypothesis a fixed $q^{n}$-th power of it lies in $A$, and the same holds for its
conjugates in~$S_0$: every element
of the finite quotient has order dividing~$q^{n}$. A finite group all of whose elements have order a power of $q$
has order dividing a power of~$q$.
\end{proof}
\leanline{\lean{existscutcoverMultiplicitydvd}}

\begin{proposition}[a period over the Bessel quotient as a finite sum]\label{p:sumcut}
Let $E$ have Weil data, $\mathcal{Y}$ a grading of a hermitian space with frame $V$, with a finite set $X$ and a
string, $A$ a level of~$V$ and $\mathsf{a}$ a complex function on $G_\infty(V)$. There is a complex number $C$ such
that the following holds for every function $F_1$ on $G(V)$ that is generic for~$\mathcal{Y}$, every automorphic
function $F_2$ on~$G(V)$, every level $J$ of~$V$, every finite set $T$ of cosets of $J$ in $G_f(V)$ and every complex
function $\mathsf{b}$ on the cosets of~$J$; write \m{F=F_1F_2}. If
\begin{enumerate}[label=(\alph*)]
\item $F$ is integrable on \m{\Besr{\mathcal{Y}}\backslash G(V)};
\item for every \m{p\in\Besf{\mathcal{Y}}\cap A} and every coset $t$: \m{t\in T} if and only if \m{p\,t\in T};
\item $\cov{\Besf{\mathcal{Y}}}{A}{T}{g_fJ}\ne0$ for every \m{g\in G(V)} with \m{F(g)\ne0};
\item $\int_{[\Bes{\mathcal{Y}}]}F(\varsigma g)\,d\varsigma=[\Kz{V}:J]\cdot\mathsf{a}(g_\infty)\,\mathsf{b}(g_fJ)$ for
  every \m{g\in G(V)} with \m{g_fJ\in T},
\end{enumerate}
then
\[\int_{\Besr{\mathcal{Y}}\backslash G(V)}F\;=\;C\cdot\sum_{t\in T}\frac{\mathsf{b}(t)}{\cov{\Besf{\mathcal{Y}}}{A}{T}{t}}.\]
\end{proposition}
\begin{proof}
Normalise the Haar measure of $\Besf{\mathcal{Y}}$ so that the compact open subgroup \m{\Besf{\mathcal{Y}}\cap A}
has measure~$1$. The Bessel group acts properly on $G(V)$, so there is a cut-off function for this action: a
function $\beta_0$ on $G(V)$ whose integral over each orbit of $\Bes{\mathcal{Y}}$ is~$1$. Let \m{\alpha_0(g_\infty)} be
the integral of $\beta_0$ over the finite part of the orbit, and let $C$ be the integral of $\mathsf{a}$ against
$\alpha_0$ times the measure of $G_\infty(V)$; it depends on $\mathcal{Y}$, $A$, $\mathsf{a}$ only. Let $\psi_T$
be the function on the cosets of $J$ equal to the inverse of the multiplicity on $T$ and to $0$ elsewhere. By (c)
the integral of $F$ may be restricted to the set where the multiplicity does not vanish. Integration in stages
over \m{\Besr{\mathcal{Y}}\backslash G(V)} (a result of the background library on quotient measures, of the type of
Weil's integration formula) writes it as the integral, against the density \m{\alpha_0(g_\infty)\psi_T(g_fJ)}, of
the inner integrals $\int_{[\Bes{\mathcal{Y}}]}F(\varsigma g)\,d\varsigma$; this uses the identity
\[\int_{\Besf{\mathcal{Y}}}\psi_T(\varsigma^{-1}t)\,d\varsigma=1\qquad\text{for every coset }t\text{ of non-zero multiplicity,}\]
whether $t$ lies in $T$ or not. It holds because the multiplicity is invariant under all of $\Besf{\mathcal{Y}}$
(Lemma~\ref{l:mult}(i)), while the indicator of $T$ is invariant only under \m{\Besf{\mathcal{Y}}\cap A}, by~(b): the
integrand is the inverse of the multiplicity of $t$ on the set of the $\varsigma$ with \m{\varsigma^{-1}t\in T}, which
is a union of exactly that many cosets of \m{\Besf{\mathcal{Y}}\cap A} (Lemma~\ref{l:mult}(ii)), and it vanishes
elsewhere. For a coset of multiplicity zero the integrand vanishes identically; by (c) such cosets carry none of
the integral of~$F$. By (d) the inner integrals are
$[\Kz{V}:J]\,\mathsf{a}(g_\infty)\mathsf{b}(g_fJ)$. The density is a product, the real part gives~$C$, and the finite
part is a sum over the finite set $T$ of \m{\psi_T(t)\mathsf{b}(t)} times the measure of a coset of~$J$, which is
\m{[\Kz{V}:J]^{-1}} (Lemma~\ref{l:mult}(ii)); the two indices cancel.
\end{proof}
\leanline{\lean{existsintegraleqsumovercut}}

\begin{proof}[Proof of Theorem~C]
\emph{The case \m{j=0}.} Take \m{A=\Kch{D}{\Ys{D}{0}}}; by Inputs~\ref{in:units}(d) it contains
$\inclf{\gY{0}}(\Kst{D}{0})$. For a tame character $\vartheta$ apply Proposition~\ref{p:cut} with
\m{S=\Besf{\gY{0}}}, $J=\levelY{\vartheta}{0}\subseteq A$ (Lemma~\ref{l:tame}(ii)), \m{\mathsf{K}=A},
$H_0=\Kaw{\lvn{\vartheta}}{\Ys{D}{0}}\cap A$, \m{q=\lvn{\vartheta}} and $\mathsf{W}$ the set of the points where some
function of stage~$0$ does not vanish. For \m{h\in A} and \m{u\in J\subseteq A} the conjugate \m{huh^{-1}} lies in the group~$A$, so the exponent
\m{n=0} works; and Lemma~\ref{l:poinbasic}(ii) moves a point of $\mathsf{W}$ by an element of~$S$ into the
principal level of the filling data of~$D$, which is~$A$ (Inputs~\ref{in:units}(d)).

\emph{The case \m{j\ne0}.} Take \m{A=\Kz{\Ys{D}{j}}}; it contains $\inclf{\gY{j}}(\Kst{D}{j})$ by
Lemma~\ref{l:inclK}, whose hypotheses are Lemma~\ref{l:vals}(iv), Inputs~\ref{in:units}(b) and the congruence
condition in the definition of~$\Kst{D}{j}$. For every finite point $h$ choose an adapted pair \m{(\mathbf{b}(h),\beta(h))}
(Lemma~\ref{l:adapted}(ii)). Fix~$\vartheta$, put \m{Q=\Qs{\vartheta}} and call the \emph{profile} of $h$ the family
of the integers \m{\ord_w\beta(h)_a-\ord_w\beta(1)_a}, for the string indices $a$ and the places $w$ in
\m{Q\cup\cj Q}. For every profile with values between \m{-j} and~$j$ choose a point, with this profile if there is one,
and keep its components at the places of \m{Q\cup\cj Q} only, setting the others to~$1$; let $\mathsf{K}$ be the
union of the finitely many cosets \m{h_0\,\Kz{\Ys{D}{j}}} of these points~$h_0$. It is compact. Apply
Proposition~\ref{p:cut} with \m{S=\Besf{\gY{j}}}, $J=\levelY{\vartheta}{j}$, $H_0=\Kaw{\lvn{\vartheta}}{\Ys{D}{j}}$,
\m{q=\lvn{\vartheta}} and $\mathsf{W}$ the set of the finite parts of the points where a function of stage $j$ of a
ramified column, with \m{y_0\in\Kch{D}{\Ys{D}{0}}}, does not vanish. The hypotheses: a point of $\mathsf{K}$ is
integral at the places not dividing~$\lvn{\vartheta}$, because the places of \m{Q\cup\cj Q} divide it
(Inputs~\ref{in:units}(a)); for \m{u\in J} the conjugate \m{huh^{-1}} agrees with an element of $\Kz{\,\cdot\,}$ outside
\m{Q\cup\cj Q}, and at a place $w$ of $Q$ the component $u_w$ of $u$ is congruent to $1$ modulo~$w$, that is,
\m{u_w\in1+\mathfrak{p}_w\mathrm{M}_I(\mathfrak{o}_w)}, where $\mathfrak{p}_w$ is the maximal ideal of
$\mathfrak{o}_w$ and \m{\mathrm{M}_I(\mathfrak{o}_w)} the ring of the matrices over~$\mathfrak{o}_w$. Let $\ell$ be the
residue characteristic of~$w$. By the binomial expansion, the $\ell$-th power maps
\m{1+\mathfrak{p}_w^{\nu}\mathrm{M}_I(\mathfrak{o}_w)} into \m{1+\mathfrak{p}_w^{\nu+1}\mathrm{M}_I(\mathfrak{o}_w)}
for every \m{\nu\ge1}; so \m{u_w^{\ell^{\nu}}\in1+\mathfrak{p}_w^{\nu+1}\mathrm{M}_I(\mathfrak{o}_w)} for all
\m{u\in J} and all~$\nu$. These subgroups are stable under powers, and $\ell$ divides~$\lvn{\vartheta}$. As they
form a basis of neighbourhoods of~$1$, one exponent~$n$, chosen for the finitely many places of~$Q$, puts
\m{u_w^{\lvn{\vartheta}^{n}}} into the open subgroup \m{h_w^{-1}\GL_I(\mathfrak{o}_w)h_w} for every \m{u\in J} at
once; at the conjugate place the component of \m{(huh^{-1})^{\lvn{\vartheta}^{n}}} is then integral with integral
inverse as well, the element being unitary. Finally let \m{h\in\mathsf{W}}. By
Proposition~\ref{p:nearstage} the lattice of $h$ has near lines; by Lemma~\ref{l:nearzero} for \m{h=1} and
\m{Q=\emptyset} the string diagonal at $1$ has exactly the reference orders; so the profile of $h$ has values in
\m{[-j,j]} and the diagonal of $h$ has the orders of that of $1$ outside \m{Q\cup\cj Q}. Let $h_0$ be the chosen
point with this profile. By Lemma~\ref{l:adapted}(vi), adapted bases at $h_0$ and at $h$ have string diagonals of
the same orders at every place, and Proposition~\ref{p:orbit} gives \m{k\in\Kz{\,\cdot\,}} with
$\varsigma=h_0\,k\,h^{-1}\in\Besf{\gY{j}}$; then \m{\varsigma h=h_0k\in\mathsf{K}}. The three conclusions of
Proposition~\ref{p:cut} are (i), (ii) and (iii) of the theorem; for \m{j=0}, (ii) includes \m{h\in A}.
\end{proof}

\sect{The unfolding}{s:unfolding}

\begin{proof}[Proof of Theorem~B]
By induction on~$j$, for all $\mathcal{Y}$ and all functions~$\mathsf{a}$.

\emph{\m{j=0}.} Take \m{C=\mathsf{a}(1)}, the value at the tuple of trivial points, and \m{\mathsf{w}=S(1)}, the value at the
tuple of trivial cosets: a combination with one coefficient~$1$, which is~(i). (ii) holds because no index is
moved: \m{q'=q} for every~$q$, so that \m{\Omega(S(1))=S(1)}. For (iii), unfold
the Poincar\'e series twice (Inputs~\ref{in:units}(d)): the integral over \m{[G(\Ys{D}{0})]} of $z$ times
\m{\PoinD{D}\varphi}, the series over all rational points, and the integral over
$\Besr{\mathcal{Y}}\backslash G(\Ys{D}{0})$ of $z$ times the series over \m{\Besr{\mathcal{Y}}}, which is the
function of stage~$0$, are both equal to $\sum_a\phi_a\int_{G}z\cdot\mathrm{cb}_a$. The second is integrable, so the
hypothesis at the trivial tuples evaluates it as \m{\mathsf{a}(1)L(S(1))}.

\emph{From $j$ to \m{j+1}.} Theorem~C at~$j$ gives the level $A$ and, for every $\vartheta$, a cut \m{T=T(\vartheta)}.
For every tuple $\xi$ let $C(\xi)$ be the constant of Proposition~\ref{p:sumcut} for the grading $\gY{j}$, the
level~$A$ and the function $\xi_j'\mapsto\mathsf{a}(\xi\text{ with its entry }j\text{ replaced by }\xi_j')$. Apply the
induction hypothesis at $j$ to the grading $\gY{j}$ and the function \m{\xi\mapsto C(\xi)}; it gives a
constant~$C$. Given $\vartheta$, $M$ and~$S$, put
\[S'(q)=\sum_{t\in T}\frac{1}{\cov{\Besf{\gY{j}}}{A}{T}{t}}\;S\bigl(q\text{ with its entry }j\text{ replaced by }t\bigr)\]
and let $\mathsf{w}$ be the vector that the induction hypothesis gives for~$S'$.

(i) Each \m{S'(q)} is a combination of values of $S$ at tuples whose new entry is integral away from
$\lvn{\vartheta}$ (Theorem~C(ii)), with coefficients $1/\cov{\Besf{\gY{j}}}{A}{T}{t}$ that lie in
\m{\Oaw{\lvn{\vartheta}}}, because the multiplicities divide powers of $\lvn{\vartheta}$ (Theorem~C(iii)). If
\m{\Qs{\vartheta}=\emptyset} the sum of the absolute values of their conjugates is at most~$|T|$, and $T$ injects
into the quotient of $\Kz{\Ys{D}{j}}$ by the principal congruence subgroup of~$\piGb{D}$: its elements are integral
everywhere, and two of them that are congruent modulo $\piGb{D}$ differ by an element of
$\levelY{\vartheta}{j}$ (Lemma~\ref{l:vals}(v)). The order of that quotient is at most the norm of $\piGb{D}$ to
the power \m{|\YI{D}{j}|^{2}}, which is at most \m{\cutB} by Lemma~\ref{l:vals}(i) and Inputs~\ref{in:units}(c).
Substituting these combinations into the admissible combination for $S'$ multiplies the sizes: (i) holds with
\m{\cutB^{\,j}\cdot\cutB}.

(ii) Let $\mathbf{b}$ and $\Omega$ be as in (ii) for \m{j+1}. Then \m{\Omega(S'(q))} is the sum over \m{t\in T} of the same
coefficients times $S$ at the moved tuple with entry $\inclf{\gY{j}}(\mathbf{b}_j)\,t$ at~$j$. Since
$\inclf{\gY{j}}(\mathbf{b}_j)\in\Besf{\gY{j}}\cap A$, multiplication by it permutes $T$ (Theorem~C(i)) and keeps the
multiplicity (Lemma~\ref{l:mult}(i)); so \m{\Omega(S'(q))=S'(q')} with $q'$ the tuple moved at the indices \m{<j},
and the induction hypothesis gives \m{\Omega(\mathsf{w})=\mathsf{w}}.

(iii) Let $z$, $\varphi$, $L$ satisfy the hypothesis for \m{j+1} and the grading~$\mathcal{Y}$. By
Lemma~\ref{l:tame}(iii) and Inputs~\ref{in:lib}(c), with $z^{\flat}$ the descent of $z$ times the index
$[\Kz{\Ys{D}{j}}:\levelY{\vartheta}{j}]$,
\[\int_{[G(\Ys{D}{j+1})]}z\cdot\chainY{\vartheta}{j+1}{\bigl(\PoinD{D}\varphi\bigr)}=\int_{[G(\Ys{D}{j})]}z^{\flat}\cdot\chainY{\vartheta}{j}{\bigl(\PoinD{D}\varphi\bigr)},\]
and $z^{\flat}$ is a form of the conjugate type on~$\Ys{D}{j}$. It remains to verify the hypothesis of (iii) at $j$
for $z^{\flat}$, the function $C(\,\cdot\,)$, the map $S'$ and the same~$L$. Fix $\xi$ and $y$ and let $F$ be the
function of stage $j$ on $\gY{j}$ times~$z^{\flat}$, assumed integrable. Proposition~\ref{p:sumcut} applies to it:
the stage function is generic (Lemma~\ref{l:stage}(ii)), $T$ is stable (Theorem~C(i)), and where $F$ does not vanish
the multiplicity divides a power of $\lvn{\vartheta}$, so it is not zero (Theorem~C(iii)). Its hypothesis (d): at a point \m{(\xi_j',t)} with \m{t\in T}, the value of $z^{\flat}$ is the index
$[\Kz{\Ys{D}{j}}:\levelY{\vartheta}{j}]$ times the theta integral of~$z$, and by one unfolding step
(Inputs~\ref{in:lib}(b)) the inner integral over \m{[\Bes{\gY{j}}]} of the stage function times that theta integral
is the integral over
$\Besr{\mathcal{Y}}\backslash G(\Ys{D}{j+1})$ of the function of stage \m{j+1} for the tuples with entry
\m{(\xi_j',t)} at~$j$, times~$z$, and this product is integrable by the same Input, which is the premise that the hypothesis for
\m{j+1} asks for; the stage function at $j$ does not see the new entry (Lemma~\ref{l:stage}(i)). By the
hypothesis for \m{j+1} this integral is $\mathsf{a}$ at the changed tuple times $L$ of $S$ at the changed tuple of cosets; the
conditions \m{y_{2\rhoD{D}+2}=1} and \m{y_0\in\Kch{D}{\Ys{D}{0}}} persist, the latter for \m{j=0} by Theorem~C(ii).
Proposition~\ref{p:sumcut} now gives
$\int F=C(\xi)\sum_{t\in T}L(S(\dots t\dots))/\cov{\Besf{\gY{j}}}{A}{T}{t}=C(\xi)\,L(S'(\text{cosets of }y))$, which
is the hypothesis at~$j$.
\end{proof}

\sect{From an invariant combination to a good datum}{s:good}

In this section the translate sums and the pair translates are those of Definition~\ref{d:redvac}, and
\emph{vac} is the vacuum of the space in question.

\begin{proposition}[the pairs of lifts evaluate admissibly]\label{p:evalpair}
In Setting~\ref{s:data} let $i$ be a natural number with \m{i\le\rhoD{D}}. Then \m{\pairEvalY{i}} is an admissible
evaluation for $\ND{D}$, $\PiD{D}$ and $\piGb{D}$ in the sense of the note \emph{Two consecutive theta lifts}: there is
a complex number $C_i$ such that for every natural number $n$ prime to $\ND{D}$, every proposition $p$ and every
finite vector $t_0$ of $\Hom(\Ys{D}{2i+1},\Ys{D}{2i}^{-}\oplus\Ys{D}{2i+2})$ for which \m{\piGb{D}\,t_0} is integral
on the rows of the first summand and \m{\pi_g\,t_0} on the rows of the second, at all places not
dividing~$n$ (two conjugate weights),
\[\begin{gathered}\pairEvalY{i}\bigl(\varrho_f(t_0)\mathrm{vac}\bigr)\in\Adm{v}{\mathcal{I}_n}{n}{d_i}{p}{1},\\
v(t_1,t_2)=C_i\cdot\transum{\XIo{D}{2i}}{\XI{D}{2i+1}}{\PiD{D}}{\ND{D}}{\XIo{D}{2i+2}}\bigl(\varrho_f(t_1)\mathrm{vac}\tensor\varrho_f(t_2)\mathrm{vac}\bigr),\end{gathered}\]
where $\mathcal{I}_n$ is the set of the pairs of finite matrices \m{(t_1,t_2)} with \m{\pi_g\,t_1} and
\m{\pi_g\,t_2} integral at the places not dividing~$n$, and $d_i$ is the order of \m{\Oh/\piGb{D}} to the
power $(\nD{D}{2i}+\nD{D}{2i+2})\cdot\nD{D}{2i+1}$.
\end{proposition}
\begin{proof}
By the definition of \m{\pairEvalY{i}} this is the theorem \emph{admissible evaluation} of the note \emph{Two
consecutive theta lifts} (Theorem~D there), applied to the middle pair for \m{i<\rhoD{D}} and to the last pair for
\m{i=\rhoD{D}}. Its hypotheses, in the letters of that theorem: $\PiD{D}$ is a split level for $\ND{D}$, and
$\ND{D}$ is split and squarefree along it with the principal idele $\piId{D}$ of $\piN{D}$
(Inputs~\ref{in:units}(b)). The middle grading \m{\gY{2i+1}} spans the lattice of
\m{(\ideA{D},\piId{D},\lsc{D}{2i+1})} (Lemma~\ref{l:vals}(iv)). The integral adele $\beta$ is~$\piGb{D}$, and
\m{\delta=\cj{\cgD{D}}^{-1}\cgD{D}}, of norm one, satisfies
\m{\delta=(\cj\alpha/\alpha)\cdot(\ideA{D}/\cj{\ideA{D}})} because \m{\ideA{D}=\alpha/\cj{\cgD{D}}}. The family
$\kappa$ is $1$ at every index except the point of $\XI{D}{m}$ in the last pair, where it is~$\delta$. At the
places above $\ND{D}$, \m{\delta=1} and $\piGb{D}$ is a unit, because chosen places do not divide~$\ND{D}$
(Lemma~\ref{l:vals}(ii)); and \m{\pi_g\kappa_b} and \m{\piGb{D}\kappa_b} are integral, $\delta$ having at every
place the order of \m{\piGb{D}/\pi_g}. \emph{The hypothesis on weights}: the zero part of the middle pair joins
only indices that carry the same weight; that of the last pair also joins the point of $\XI{D}{m}$ to the only
string vector of~\m{\gY{2\rhoD{D}}}, where the two weights are $\piGb{D}$ and \m{\pi_g\delta}, of the same absolute
value at every place: this is the reason for \m{\kappa=\delta} there. \emph{The hypothesis on lattices}: scaled
neighbour vectors land in the lattices of the outer spaces. In the middle pair this is the lattice lemma of that
note (Lemma~4.3 there) for $\gY{2i}$ and $\gY{2i+2}$, whose scales $\lsc{D}{2i}$ and $\lsc{D}{2i+2}$ are the powers
of \m{\ideA{D}/\cj{\ideA{D}}} with exponents \m{\rhoD{D}-i+1} and \m{\rhoD{D}-i} (Lemma~\ref{l:vals}(iv)). In the
last pair it is the same lemma for \m{\gY{2\rhoD{D}}}, with exponent~$1$, applied to the sum of the two congruences
of a neighbour vector, one of them multiplied by the unit $\varepsilon$ with \m{\cj\varepsilon\varepsilon=-1}; the
frame of the standard space $\Ys{D}{m}$ is~$1$.
\end{proof}
\leanline{\lean{isAdmissibleEvalpairEvalY}}

\begin{proposition}[evaluated data are combinations of reduced vacua]\label{p:redvac}
In Setting~\ref{s:data} there is a complex number $C$ such that for every tame character $\vartheta$ of $D$ and
every tuple \m{y=(y_j)} of finite points with $y_j\in\Kaw{\lvn{\vartheta}}{\Ys{D}{j}}$ for all~$j$:
\[\begin{gathered}\evD{\vartheta}{y}{\rhoD{D}}\tensor\dots\tensor\evD{\vartheta}{y}{0}\\
{}\in\Adm{C\cdot\redvac{D}{\cdot}}{\Tadm{D}{\lvn{\vartheta}}}{\lvn{\vartheta}}{\cutB^{\,2m}}{\Qs{\vartheta}=\emptyset}{1}.\end{gathered}\]
\end{proposition}
\begin{proof}
Write \m{n=\lvn{\vartheta}}. \emph{Each moved datum.} \m{\ups{\vartheta}{j}} is named by an indicator modulo the
product $\mathfrak{m}$ of the chosen places and the places of~$\Qs{\vartheta}$, and \m{\piGb{D}\,n\in\mathfrak{m}}; by
Inputs~\ref{in:units}(e) it is an admissible combination, away from~$n$, of translates \m{\varrho_f(t)\mathrm{vac}}
with \m{\pi_gt} integral at the places not dividing~$n$, with denominator a power of the norm of $\piGb{D}$
and size at most~$1$. The Weil operator of a finite pair $(y_j,y_{j+1})$ is linear, and by Inputs~\ref{in:units}(f) it sends
such a translate to an admissible combination of translates of the same kind: the pair is integral at the places
not dividing~$n$, so the translations stay in the same group; its finite multiplier has $\ell$-adic absolute
value $1$ for $\ell$ prime to~$n$, and is $1$ when \m{\Qs{\vartheta}=\emptyset}, the pair then lying in the
stabiliser of the lattice; and the values of the twisting character are $\ell$-adic units, of complex absolute
value~$1$. So $\tUps{\vartheta}{y}{j}$ is such a combination. \emph{Each pair.} The arrow product is bilinear, and
the arrow product of two translated vacua is a root of unity times a translated vacuum of
$\Hom(\Ys{D}{2i+1},\Ys{D}{2i}^{-}\oplus\Ys{D}{2i+2})$, at the vector with rows \m{-\varepsilon\cj{t}} on the first
summand and $t'$ on the second, which satisfies the integrality of Proposition~\ref{p:evalpair}: the conjugated
block has the weight~$\piGb{D}$, the other block the weight~$\pi_g$. Applying
\m{\pairEvalY{i}} and Proposition~\ref{p:evalpair} (with $n$, prime to $\ND{D}$), $\evD{\vartheta}{y}{i}$ is an
admissible combination of the vectors \m{v(t_1,t_2)} of that proposition, each of them $C_i$ times the translate
sum of the pair~$i$ applied to a tensor of two translated vacua.
\emph{The tensor.} A pure tensor of admissible combinations is an admissible combination of the pure tensors, with
the product of the denominators and of the sizes; the pure tensor of the vectors \m{v(t_1,t_2)} of the pairs is
\m{\prod_iC_i} times a reduced vacuum, up to the normalisation of the vacuum, which is absorbed in~$C$. The
denominators: all are powers of the order of \m{\Oh/\piGb{D}}, which is \m{\gD{D}^{\fD{D}}} (Lemma~\ref{l:vals}(i)),
with exponents bounded, for each pair, by \m{4\,\nD{D}{m}^{2}}; so each pair contributes a divisor of \m{\cutB^{4}}
and the \m{\rhoD{D}+1} pairs a divisor of \m{\cutB^{4(\rhoD{D}+1)}=\cutB^{\,2m}}.
\end{proof}
\leanline{\lean{existsisAdmissibleCombreducedVacuum}}

\begin{lemma}\label{l:proj}
In Setting~\ref{s:data} let $\vartheta$ be a tame character of $D$, \m{t\in\Tadm{D}{\lvn{\vartheta}}} and $i$ a
natural number, and let $\mathsf{P}_\Pi$ be the projection of the finite part to the vectors that are vacua at the
places of $\PiD{D}$ and their conjugates (Definition~\ref{d:good}(iii)). Then
\[(\mathsf{P}_\Pi\tensor\mathsf{P}_\Pi)\bigl(\ptr{D}{\tdat{D}{t}{\cdot}}{i}\bigr)=\tdat{D}{t}{2i}\tensor\tdat{D}{t}{2i+1}.\]
\end{lemma}
\begin{proof}
The places of $\PiD{D}$ and their conjugates do not divide $\lvn{\vartheta}$ and are not chosen, so
\m{\tdat{D}{t}{j}} is fixed by $\mathsf{P}_\Pi$ (Inputs~\ref{in:units}(b)); for such a vector,
\m{\mathsf{P}_\Pi\varrho_f(t')} is \m{\varrho_f(t')} if $t'$ is integral at these places and $0$ otherwise
(Input~1(f) of the note \emph{Two consecutive theta lifts}). In the sum defining
\m{\ptr{D}{\tdat{D}{t}{\cdot}}{i}} the vector \m{-y/\ND{D}} is integral at $\PiD{D}$ only for \m{y=0}, and that
term is the tensor of the two data.
\end{proof}
\leanline{\lean{projOnFinPartpairTranslates}}

\begin{lemma}\label{l:redvaccomb}
In Setting~\ref{s:data} let $q$ be a natural number, $p$ a proposition and \m{t\in\Tadm{D}{q}}. Then
\[\begin{gathered}\redvac{D}{t}\in\Adm{v}{\Tadm{D}{q\,\ND{D}}}{q}{1}{p}{\ND{D}^{\,\fD{D}\nD{D}{m}(\rhoD{D}+1)}},\\
v(t')=\bigotimes_{i\le\rhoD{D}}\bigl(\tdat{D}{t'}{2i}\tensor\tdat{D}{t'}{2i+1}\bigr).\end{gathered}\]
\end{lemma}
\begin{proof}
Factor by factor: \m{\ptr{D}{\tdat{D}{t}{\cdot}}{i}} is a sum, over at most \m{\ND{D}^{\fD{D}\nD{D}{m}}}
vectors~$y$, of terms with
coefficient~$1$, and each term is a root of unity times
$\varrho_f(t_1)\mathrm{vac}\tensor\varrho_f(t_2)\mathrm{vac}$ with $t_1$, $t_2$ the sums of the entries of $t$ and
of the translations by \m{\mp y/\ND{D}}, since Heisenberg operators compose up to values of the additive character.
These translations live at the places above~$\ND{D}$, so the new tuple is admissible for \m{q\ND{D}}. Tensoring the
\m{\rhoD{D}+1} factors multiplies the sizes.
\end{proof}
\leanline{\lean{isAdmissibleCombreducedVacuum}}

\begin{proposition}[an invariant combination of reduced vacua is the chain of a good datum]\label{p:good}
In Setting~\ref{s:data} let $\vartheta$ be a twist of $D$ and let $\mathsf{w}$ be a tensor in the product of the paired data
spaces of the \m{\rhoD{D}+1} pairs of definite lifts. Assume that
\[\mathsf{w}\in\Adm{\redvac{D}{\cdot}}{\Tadm{D}{\lvn{\vartheta}}}{\lvn{\vartheta}}{1}{\Qs{\vartheta}=\emptyset}{\cutB^{\,2m}\cdot\cutB^{\,2\rhoD{D}+2}},\]
and that for every tuple \m{\mathbf{b}=(\mathbf{b}_j)} with \m{\mathbf{b}_j\in\Kst{D}{j}} and \m{\mathbf{b}_{2\rhoD{D}+2}=1} the tensor $\mathsf{w}$ is fixed by the
operator
\[\bigotimes_{i\le\rhoD{D}}\omega(\mathbf{b}_{2i},\mathbf{b}_{2i+1})\tensor\omega(\mathbf{b}_{2i+1},\mathbf{b}_{2i+2}).\]
Then there is a finitely supported
$\Psi\colon\Tr{D}\to\CC$ that is a good datum for $\vartheta$ and such that for every \m{\varphi\in\Stg{D}{0}}
\[\chnn{D}{\vartheta}{\mathsf{w}}{\rhoD{D}+1}(\varphi)=\chcomb{D}{\vartheta}{\Psi}{m}(\varphi),\]
the left side being the chain of pairs extended linearly to tensors.
\end{proposition}
\begin{proof}
Substitute Lemma~\ref{l:redvaccomb} into the combination~$\mathsf{w}$: this gives a finitely supported $\Psi$ on
\m{\Tadm{D}{\lvn{\vartheta}\ND{D}}} with values in \m{\Oaw{\lvn{\vartheta}}} and
$\mathsf{w}=\sum_{t}\Psi(t)\bigotimes_i(\tdat{D}{t}{2i}\tensor\tdat{D}{t}{2i+1})$; if \m{\Qs{\vartheta}=\emptyset} its size is
at most the product of the two sizes, which the comparisons of the data bound by \m{\exp(x^{60})}. This is
condition (i) of Definition~\ref{d:good}. The chain of pairs of a pure tensor of pairs is the chain of the single
data, so $\chnn{D}{\vartheta}{\mathsf{w}}{\rhoD{D}+1}=\chcomb{D}{\vartheta}{\Psi}{m}$, because \m{m=2(\rhoD{D}+1)}. Joint
invariance (ii): regrouping the tensor of the $m$ single data into pairs carries the joint action of a tuple $\mathbf{b}$
to the action in the hypothesis, which fixes~$\mathsf{w}$. Condition (iii): let $\mathsf{R}$ be the endomorphism of the
tensor product of the paired data spaces that is, on the factor of the pair~$i$, the projection
\m{\mathsf{P}_\Pi\tensor\mathsf{P}_\Pi} followed by the translate sum of that pair (Definition~\ref{d:redvac}); on
a pure tensor of single data \m{s=(s_j)_{j<m}} it gives the tensor of the pair translates
\m{\ptr{D}{\mathsf{P}_\Pi s}{i}}, \m{i\le\rhoD{D}}. By Lemma~\ref{l:proj}, $\mathsf{R}$ fixes the reduced vacuum of
every \m{t\in\Tadm{D}{\lvn{\vartheta}}}, hence~$\mathsf{w}$. Applying $\mathsf{R}$ to
$\mathsf{w}=\sum_{t}\Psi(t)\bigotimes_i(\tdat{D}{t}{2i}\tensor\tdat{D}{t}{2i+1})$ and then the chain of pairs gives~(iii).
\end{proof}
\leanline{\lean{existsisGoodDatumofinvariant}}

\sect{Proof of the seesaw identity}{s:proofA}

\begin{proof}[Proof of Theorem~A]
\emph{The constants.} Let $C_{\mathrm{v}}$ be the constant of Proposition~\ref{p:redvac}. Let $R$ and the constants
$C_{\mathsf{t}}$ be as in Inputs~\ref{in:units}(d) (the last stage). For every tuple $\xi$ of real points let
$\mathsf{t}(\xi)$ be the functional of Proposition~\ref{p:even} for the pair index \m{i=\rhoD{D}}, the grading
$\gYlast$ of $\Ys{D}{m}$, the last arrow pair, whose finite evaluation is \m{\pairEvalY{\rhoD{D}}} by definition, and
the trivial real point. Apply Theorem~B with \m{j=m=2\rhoD{D}+2}, the grading $\gYlast$, whose string has
length~$0$, and the function $\mathsf{a}(\xi)=C_{\mathsf{t}(\xi)}$; let $C_{\mathrm{u}}$ be the constant it gives. Put
\[\cD{D}=\Bigl(\frac{C_{\mathrm{u}}\,C_{\mathrm{v}}}{\cutB^{\,2m}}\Bigr)^{-1}\ \text{if this quotient is not zero, and}\ \cD{D}=1\ \text{otherwise.}\]

\emph{The datum.} Fix a twist~$\vartheta$. Let $M$ be the tensor product of the paired data spaces and
\[S(q)=\evD{\vartheta}{q^{\circ}}{\rhoD{D}}\tensor\dots\tensor\evD{\vartheta}{q^{\circ}}{0},\]
where $q^{\circ}$ is the tuple of the chosen representatives of the cosets~$q$; let $\mathsf{w}$ be the vector of
Theorem~B. By Theorem~B(i), $\mathsf{w}$ is an admissible combination of the values \m{S(\text{cosets of }y)}, $y$ integral
away from~$\lvn{\vartheta}$, with denominator $1$ and size \m{\cutB^{\,m}}; by Lemma~\ref{l:stage}(iii) these values
are the tensors of Proposition~\ref{p:redvac}. Substituting, $\mathsf{w}$ is an admissible combination of the vectors
\m{C_{\mathrm{v}}\redvac{D}{t}}, \m{t\in\Tadm{D}{\lvn{\vartheta}}}, with denominator \m{\cutB^{\,2m}}; clearing the
constant and the denominator (Inputs~\ref{in:units}(f)), the vector \m{\mathsf{w}^{\star}=\cD{D}\,C_{\mathrm{u}}\,\mathsf{w}} is an admissible
combination of the reduced vacua themselves with denominator~$1$ and size \m{\cutB^{\,2m}\cdot\cutB^{\,m}}. It is
invariant as Proposition~\ref{p:good} requires: for a tuple $\mathbf{b}$ with \m{\mathbf{b}_j\in\Kst{D}{j}} and \m{\mathbf{b}_m=1}, the operator
$\bigotimes_i\omega(\mathbf{b}_{2i},\mathbf{b}_{2i+1})\tensor\omega(\mathbf{b}_{2i+1},\mathbf{b}_{2i+2})$ sends \m{S(q)} to $S$ at the tuple moved by the
images of the $\mathbf{b}_j$ under the inclusions $\inclf{\gY{j}}$, because the finite evaluation of an arrow pair commutes
with the finite points of the three definite groups (\emph{equivariance of the finite evaluation},
Inputs~\ref{in:units}(f)); so Theorem~B(ii) applies. Proposition~\ref{p:good} gives a good datum $\Psi$ with
$\chnn{D}{\vartheta}{\mathsf{w}^{\star}}{\rhoD{D}+1}=\chcomb{D}{\vartheta}{\Psi}{m}$.

\emph{The identity.} By the criterion (1) of the last stage it suffices to show, for \m{\varphi\in\Colram{D}},
\[\chainY{\vartheta}{m}{\bigl(\PoinD{D}(\cD{D}\,\varphi)\bigr)}=R\bigl(\chcomb{D}{\vartheta}{\Psi}{m}(\varphi)\bigr),\]
and for this to pair both sides with an arbitrary \m{\Xi\in\FormY{m}}. Apply
Theorem~B(iii) to the form \m{z=\cj\Xi}, the column $\varphi$ and the linear form
\[L(\mathsf{w}')=\bigl\langle\Xi,R\bigl(\chnn{D}{\vartheta}{\mathsf{w}'}{\rhoD{D}+1}(\varphi)\bigr)\bigr\rangle .\]
Its hypothesis: for $\xi$ and $y$ with
\m{y_m=1} and \m{y_0\in\Kch{D}{\Ys{D}{0}}}, Proposition~\ref{p:even} says that the function of stage $m$ on $\gYlast$
is, on the image of \m{\incl{\gYlast}}, the coefficient at $\mathsf{t}(\xi)$, on the block diagonal elements, of
\[\psi=\chnn{D}{\vartheta}{\evD{\vartheta}{y}{\cdot}}{\rhoD{D}+1}(\varphi);\]
by (2) of the last stage its integral against \m{\cj\Xi} is
\[C_{\mathsf{t}(\xi)}\langle\Xi,R\psi\rangle=\mathsf{a}(\xi)\,L\bigl(S(\text{cosets of }y)\bigr),\]
again by Lemma~\ref{l:stage}(iii). The conclusion of Theorem~B(iii) is
\[\bigl\langle\Xi,\chainY{\vartheta}{m}{\bigl(\PoinD{D}\varphi\bigr)}\bigr\rangle=C_{\mathrm{u}}\cdot\bigl\langle\Xi,R\bigl(\chnn{D}{\vartheta}{\mathsf{w}}{\rhoD{D}+1}(\varphi)\bigr)\bigr\rangle .\]
Multiplying by $\cD{D}$ and using linearity in the tensor gives the pairing of $\Xi$ with
\[R\bigl(\chnn{D}{\vartheta}{\mathsf{w}^{\star}}{\rhoD{D}+1}(\varphi)\bigr)=R\bigl(\chcomb{D}{\vartheta}{\Psi}{m}(\varphi)\bigr).\]
As $\Xi$ is
arbitrary, the two forms agree, and (1) of the last stage gives the theorem. The number $\cD{D}$ was fixed before
$\vartheta$ and~$\varphi$.
\end{proof}

\sect{Where the results are used}{s:used}

Theorem~A is used once, in the proof of Theorem~A of the write-up (\emph{short forms}, unit \emph{Endgame}): its
constant $\cD{D}$ is the constant in the Poincar\'e maps \m{\gpt{D}{\gamma}\circ\PoinD{D}\circ(\cD{D}\,\cdot)} from
ramified columns to rows: the restriction of this composite to $\Colram{D}$ is the Poincar\'e series
$\mathcal{P}_D(\gamma)$ of Setting~2 of the overview, and the ``fixed non-zero constant'' of that Setting is
$\cD{D}$ (in the formal text the map is the composite of the projection, the Poincar\'e series and the inclusion
of the ramified columns scaled by the constant chosen from Theorem~A). The roman $\PoinD{D}$ of this note is the
series before the projection and the scaling. For every twist the good datum turns the size and the divisibility of
\m{\iot{\vartheta}} of such a row into those of \m{\Bmu{D}{\vartheta}{\Psi}}, which the units \emph{Columns},
\emph{Curve} and \emph{TwistedHecke} estimate. The map \m{\iotL{\vartheta}{(\,\cdot\,)}} restricted to the rows is
the chain of lifts $\Theta_{D,\vartheta}$ of the overview; Corollary~\ref{c:groupG} and Lemma~\ref{l:blocks} are used by
the units \emph{Rows} and \emph{Exchange}, and Proposition~\ref{p:fits} with the maps of Definition~\ref{d:poin} by
the unit \emph{Filling}, which in turn supplies the last stage (Inputs~\ref{in:units}(d)).

\sect{Notes on the formal text}{s:notes}

\subhead{What kind of seesaw.} The identity moves a chain of theta lifts from the tower of indefinite spaces
$\Ys{D}{j}$ to the tower of definite spaces $\Xs{D}{j}$ sitting inside them, at the price of a Bessel period; the
header of the formal file calls it a seesaw identity and names as related work B.~Wang (2025), Gomez--Zhu (2014)
and Gan--Ichino (2016) for the transfer of Bessel and Fourier--Jacobi periods, and Gross (1999), Hsieh--Namikawa
(2015) and X.~Zhang (2022) for periods on definite groups as finite sums. Kudla's seesaw principle in the
form of the background library (which the catalogue marks as a weaker form) is used in the unit \emph{Exchange}, not
here: this unit unfolds.

\subhead{Shape of the statements.} (1) Theorem~A is stated for twists, because a good datum is defined for twists;
Theorems~B and~C and all statements about the chain hold for tame characters. (2) In Theorem~B the grading
$\mathcal{Y}$ is arbitrary with the right string length, and the subtraction in \m{l=2\rhoD{D}+1-j} is truncated; the
proof uses it for $\gY{j}$ and, at \m{j=m}, for $\gYlast$. The integrability in (iii) is a hypothesis \emph{inside}
the hypothesis: the formal text never has to prove that the function of stage \m{j+1} times $z$ is integrable,
because the unfolding step provides it. (3) The constant $\cD{D}$ is Mathlib's normalising unit of
\m{C_{\mathrm{u}}C_{\mathrm{v}}/\cutB^{\,2m}}, which is $1$ when the quotient vanishes; the formal text does not show
that the quotient is non-zero and does not need it. (4) The bound \m{\exp(x^{60})} in the definition of a good
datum is asked only for \m{\Qs{\vartheta}=\emptyset}, and so are all size clauses in this note. (5)
Lemma~\ref{l:blocks}(ii) is stated for every commutative ring, with the cardinality of an infinite set counted
as~$0$. (6) The grading $\gYlast$ is defined for every number $M$ with a proof that \m{m=M}, and the functions of
stage \m{j+1} are defined for an arbitrary grading of~$\Ys{D}{j+1}$, which enters only through the constant term.
(7) The groups of a hermitian space with frame depend on less than the space: $G_f(V)$ on the index set, $G(V)$ on
the signs; the notation of this note records the space whose frame gives the rational structure.

\subhead{Arguments in the census that Lean does not take.} The census of the write-up reads arguments from the
text. For the following declarations it lists section variables that the statement binds again, or that occur
only as the name of a field: the tame character, the column and the tuple of finite points in
Proposition~\ref{p:even} (bound inside the statement; Lean takes only $\xi$); a twist in Proposition~\ref{p:redvac}
(the statement quantifies over tame characters); a Boolean named like the constant of a grading in
Lemma~\ref{l:vals}(iii); a group, a subgroup and a set of cosets in Propositions~\ref{p:cut} and~\ref{p:sumcut};
the two subgroups of Definition~\ref{d:mult}, listed twice. The census ends the statement of Theorem~C at its first
\emph{let} and so lists only the stage and its bound; the theorem is stated here in full, from the source. For five
definitions whose type ends in a linear map or a function (\m{\liftY{\vartheta}{j}{(\,\cdot\,)}},
\m{\chainY{\vartheta}{j}{(\,\cdot\,)}}, \m{\iotL{\vartheta}{(\,\cdot\,)}}, \m{\lastChar{\vartheta}{\cdot}},
\m{\blk{j}{\cdot}}) the census counts the argument of that map as an argument of the definition; the notation shows
it. Unused hypotheses were not determined: no Lean was run for this note.

\subhead{The second reading.} A second reader compared the member texts of the whole unit and a sample of the proof
steps with the Lean source, after the dossier was made. Where her reading and the dossier differ, this note
follows her, having checked each point in the source: one scalar $\cD{D}$ for the datum, chosen before the twist, and a
restriction that acts on the definite chain only, the indefinite chain being paired with a test form directly
(proof of Theorem~A); in Proposition~\ref{p:sumcut} the average over the finite Bessel group is taken for every coset
of non-zero multiplicity, the multiplicity and not the indicator of the cut being invariant under the whole
group; the common weight of the translations is $\pi_g$, and the two blocks of a combined vector carry the two
conjugate weights (Propositions~\ref{p:evalpair} and~\ref{p:redvac}); Proposition~\ref{p:redvac} holds for tame
characters; the normal form of the Poincar\'e series is stated with the coordinates of general filling data, the
translation by the lattice scale entering only through $\colco{D}$; and the constant of Theorem~A keeps the zero
case of the normalisation.

\subhead{What the final theorems use.} The reader's dossier sets nothing of the unit aside. Of this note, the
instance that the index sets $\YI{D}{j}$ are non-empty is used only implicitly.

\subhead{How this note was written.} From the reader's dossier of the unit, after its program check passed, and
from the Lean source. Of the 32 structured proofs that belong to this note all were read. In the Lean text I read
completely the files of the theorem (all three proofs line by line), of the chain and the gradings, of the stage
functions, of the adapted bases, of the two lattice steps, of the standard lattice, of the cuts, of the Poincar\'e
series at Bessel points and of the good datum, and the six lemmas on valuations and levels; among them the four
proofs that the reader asked to sample first (Propositions~\ref{p:redvac}, \ref{p:step}, \ref{p:normal}
and~\ref{p:sumcut}). There the structure, the order of the steps and every cited lemma agree with the dossier; the
long chains of rewriting and of automated case analysis at the end of these four proofs, and of
Propositions~\ref{p:orbit}, \ref{p:rev} and~\ref{p:evalpair}, were not run, and the note describes their effect from
the terms involved rather than reproducing them. The definitions recalled from the background library and from
other units (spaces with frame, gradings, Bessel groups, generic functions, the unfolding step, arrows, chains, good
data, filling data) were read in the Lean source; the Inputs are stated as their Lean statements read, the last
stage, the unfolding of Poincar\'e series, the unfolding step and the transposition of lift and descent after the
source, the others after the dossier. The proofs here follow the formal ones; they are shorter where the formal text
spends lines on coercions, on measurability and on the choice of representatives. The program check of the unit's
dossier covered \UnitSteps{} leaf steps of \UnitProofs{} proofs.

\subhead{Revision.} After an independent referee's report the note was revised. The Inputs were completed: the
discrete-series input on which every lift rests, four further facts of the library and three of the unit
\emph{Filling} are now listed, with their status in the catalogue; Inputs~\ref{in:units}(e) is quoted with the
adele that sets the weight of the translates, read again in the Lean source, so that step~(e) of
Proposition~\ref{p:nearstage} follows from it. The proofs of Lemma~\ref{l:inclK}, of
Propositions~\ref{p:rev}, \ref{p:step} and~\ref{p:evalpair}, of the last step of Proposition~\ref{p:genericpoin}
and of the first hypothesis in the proof of Theorem~C were completed. The two boxes of Inputs have their own
numbers, Theorem~B is stated before Theorem~C, the constant of Theorem~A is written $\cD{D}$, and the objects of
the library carry the names and the arguments that the other notes of the unit give them.

\head{Appendix A. Concordance with the Lean artifact}
{\fontsize{8.6}{11.6}\selectfont\renewcommand{\leanfont}{\ttfamily\fontsize{7.7}{10}\selectfont\color{leancol}}
\setlength{\LTleft}{0pt}\setlength{\LTright}{\fill}\setlength{\tabcolsep}{0pt}
\begin{longtable}{@{}p{26mm}p{63mm}p{40mm}r@{}}
\textbf{In this note} & \textbf{Lean declaration} & \textbf{File} & \textbf{Lines}\\\midrule\endhead
\ConcordanceRows
\end{longtable}}

The statement of the theorem, as the artifact has it:
\quoted{seesaw}

\head{Notation}
\begingroup\fontsize{8.6}{10.9}\selectfont\raggedright
\setlength{\columnsep}{9mm}\setlength{\parskip}{0pt}
\begin{multicols}{2}
\nt{$P=(N_0,d)$, $K$}{parameters; the field \m{\mathbb{Q}(\sqrt{-d})}}
\nt{$D$}{admissible data for~$P$, with components $x$, $x_1$, $b$ (scales), $\mu$ (exponent), $m$ (number of
  lifts), $r$ (rank), $g_0$, $g_1$ (primes), $\pi_g$}
\nt{$\ED{D}$, $\fD{D}$, $\gD{D}$, $\tD{D}$, $\nD{D}{j}$}{the CM field; the number of its infinite places;
  \m{g_0g_1}; \m{5r}; the rank of stage~$j$}
\nt{$\mathbf{j}$, $\alpha$, $\upsilon$, $\vpi{P}{\mu}$}{Weil data: an element of trace zero, an idele, the splitting
  character; the generator of the different of $\ED{D}$}
\nt{$\rhoD{D}$}{\m{m/2-1}}
\nt{$\piGb{D}$, $\cgD{D}$, $\piId{D}$}{the conjugate of $\pi_g$ in $\ED{D}$; an idele generating its ideal; the
  principal idele of $\piN{D}$ (the last two written as in the note of the unit \emph{Trace})}
\nt{$\ND{D}$, $\piN{D}$, $\PiD{D}$}{the level prime; an integer of norm $\ND{D}$; the places dividing it}
\nt{$\vartheta$, $\Qs{\vartheta}$, $\vartheta^{\mathrm{t}}$, $\lvn{\vartheta}$}{a tame character; its places; its tame part;
  the product of the norms of its places}
\nt{$\omega_{D,w}$, $\muc{D}{\vartheta}{j}$, $\psic{D}{\vartheta}{j}$}{residue characters at chosen places; the
  characters of the lifts; their tame parts}
\nt{$\XI{D}{j}$, $\YI{D}{j}$, $\Xs{D}{j}$, $\Ys{D}{j}$}{index sets; the definite and the graded hermitian spaces}
\nt{$\GD{D}$, $\ezero{D}$, $\gamma$}{the finite group; the good idempotent; an element of $\GD{D}$}
\nt{$E$, $\Ad$, $\Af$, $\Oh$, $\mathfrak{o}_w$}{a CM field; adeles; finite adeles; their integers; local integers}
\nt{$w$, $v$, $\cj w$, $\mathbb{F}_w$, $\Nm w$, $\ord_w$}{finite places; conjugate; residue field; norm; order}
\nt{$V$, $V'$, $W$, $I$, $J_V$, $x_V$, $x_{V,f}$}{hermitian spaces with frame; index set; Gram matrix; frame
  and its finite part (written as in the note \emph{Two consecutive theta lifts}); $W$ also a Bessel vector in Section~\ref{s:poincare} and Proposition~\ref{p:even}}
\nt{$G(V)$, $G_\infty(V)$, $G_f(V)$, $\Gamma(V)$}{adelic, real, finite and rational points}
\nt{$\Kz{V}$, $\Kw{V}{w}$, $\Kaw{n}{V}$, $\Kch{D}{V}$}{lattice stabiliser; its part at $w$; points integral away
  from~$n$; principal level at the chosen places}
\nt{$\mathcal{Y}$, $\mathcal{Y}'$, $e_{\mathcal{Y}}$, $\cst{\mathcal{Y}}$, $\mathfrak{N}$, $X$, $l$}{gradings; bijection;
  constant; norm; set; string length}
\nt{$z_a$, $\gr$, $a$, $a'$}{string vectors; grade; string indices}
\nt{$\Uni{\mathcal{Y}}$, $\Bes{\mathcal{Y}}$, $\Besf{\mathcal{Y}}$, $\Besr{\mathcal{Y}}$}{unipotent and Bessel
  groups; finite part; rational points}
\nt{$\incl{\mathcal{Y}}$, $\inclf{\mathcal{Y}}$}{the inclusion of the group of $X$}
\nt{$\psi_E$, $\mathrm{sh}_{\mathcal{Y}}$}{the standard additive character; the shift matrix}
\nt{$\mathcal{S}(W)$, $\mathcal{W}_f(W)$, $\varrho(y)$, $\varrho_f(y)$, $\omega$, $\pW{\nu}$, $\theta$}{theta data;
  finite part; Heisenberg operators; Weil representation; twisted pair action; theta functional}
\nt{$\nu$, $\Phi$, $v\tensor f$, $\boxtimes$}{a character; a datum; integrated datum; arrow product}
\nt{$\Tm{\mathcal{Y}}{\mathcal{Y}'}$, $\mathrm{ct}$, $\mathrm{unf}$}{the string map; constant term at the string
  map; unfolded function}
\nt{$\gY{j}$, $\gYlast$}{the gradings of the spaces $\Ys{D}{j}$}
\nt{$\inclY{j}$, $\blk{j}{i}$, $R_0$, $T_0$}{block inclusion; block number; sets of blocks}
\nt{$\levelY{\vartheta}{j}$, $\ltame{\vartheta}{j}$, $\kerY{\vartheta}{j}$, $\ups{\vartheta}{j}$}{level, tame part,
  kernel and finite datum of lift \m{j+1}}
\nt{$\FormY{j}$, $\datY{\vartheta}{j}$, $\liftY{\vartheta}{j}{(F)}$, $\chainY{\vartheta}{j}{(F)}$}{forms; datum; one
  lift; $j$ lifts}
\nt{$\iot{\vartheta}$, $\iotL{\vartheta}{(F)}$, $\lastRes{\vartheta}$, $\lastChar{\vartheta}{\kappa}$,
  $\groupLift{\vartheta}{\gamma}$}{the chain of $m$ lifts; with its level; the level group and its character; a lift
  of~$\gamma$}
\nt{$\mathsf{S}$, $g_{\mathsf{S}}$, $\varpi_{\mathsf{S}}$, $K_{\mathsf{S}}$, $K^{\mathrm{pr}}_{\mathsf{S}}$, $X_0$,
  $\eta$}{filling data and their components; a central element}
\nt{$\mathrm{P}_{\mathsf{S}}(\phi,\Delta)$, $\mathrm{cb}_a$, $\mathrm{cc}_{\mathsf{S}}$, $\mathsf{A}$,
  $\sigma(u)$}{Poincar\'e series; cut-off Bessel functions; column coordinates; change of coordinates; simple root sum}
\nt{$\Col{D}$, $\Colram{D}$, $\colco{D}$, $\PoinD{D}$, $\gpt{D}{\gamma}$}{columns; ramified columns; their
  coordinates; Poincar\'e series; projection to rows}
\nt{$\Oaw{n}$, $\Adm{v}{\mathcal{I}}{n}{d}{p}{B}$}{numbers integral away from~$n$; admissible combinations}
\nt{$\ell$, $\iota$, $\sigma$}{a prime; a ring homomorphism \m{\CC\to\CC_\ell}; an automorphism of~$\CC$}
\nt{$\Kst{D}{j}$, $\Stg{D}{j}$, $\zw{D}{w}$}{level and forms of stage~$j$; the central translation at~$w$}
\nt{$\Tr{D}$, $\Tadm{D}{q}$, $\tdat{D}{t}{j}$}{translations; admissible ones; translated vacua}
\nt{$\chn{D}{\vartheta}{\Phi}{j}$, $\chnn{D}{\vartheta}{\Phi}{p}$, $\chcomb{D}{\vartheta}{\Psi}{j}$}{chains of
  definite lifts: single, paired, combined}
\nt{$\resY{D}$, $\Bmu{D}{\vartheta}{\Psi}$, $\Psi$}{restriction to $\Ys{D}{m}$; the definite side; a good datum}
\nt{$\XIo{D}{j}$, $\transum{I}{I'}{\Sigma}{N}{I''}$, $\ptr{D}{v}{i}$, $\redvac{D}{t}$}{$\XI{D}{j}$ without its point;
  translate sum of three index sets, a set of places and a prime; pair translates of a family~$v$; reduced vacuum}
\nt{$\mathsf{P}_\Pi$, $\mathsf{R}$, $I'$, $I''$}{projection at $\PiD{D}$; an endomorphism (Proposition~\ref{p:good});
  index sets}
\nt{$\varphi$, $\psi$, $\xi$, $y$, $\unf{\vartheta}{\varphi}{\xi}{y}{j}{\mathcal{Y}}$}{columns; tuples of real and
  finite points; stage functions}
\nt{$\pairEvalY{i}$, $\cutB$, $\tUps{\vartheta}{y}{j}$, $\evD{\vartheta}{y}{i}$}{finite evaluation of a pair; the cut
  bound; moved and evaluated data}
\nt{$\flag{e}{i}$, $\mathbf{b}$, $\beta$, $\mathsf{k}$, $\mathsf{d}$}{flag order; adapted basis; string diagonal;
  auxiliary matrices}
\nt{$\mathfrak{a}$, $\mathfrak{b}$, $\ideA{D}$, $\lsc{D}{j}$, $\epsilon_a$, $\Lambda$, $\pi$, $\Sigma$}{ideles; the
  idele and the scale of the data; reference factors; graded lattice basis; an idele of norm~$\mathfrak{N}$; a split set of
  places}
\nt{$H$, $S$, $A$, $J$, $T$, $t$, $H_0$, $\mathsf{W}$, $\mathsf{C}$}{a group, subgroups, a set of cosets, a coset; a subgroup, a set of points, representatives; $\mathsf{W}$ also a matrix in Proposition~\ref{p:step}}
\nt{$\mathsf{K}$}{a compact subset of $G_f(V)$ in Section~\ref{s:cuts} (a compact open subgroup in the applications)}
\nt{$\cov{S}{A}{T}{t}$}{multiplicity of a cut}
\nt{$C$, $C_i$, $C_{\mathsf{t}}$, $C_{\mathrm{u}}$, $C_{\mathrm{v}}$, $\cD{D}$}{constants; the constant of Theorem~A}
\nt{$M$, $L$, $\Omega$, $\mathsf{a}$, $\mathsf{b}$, $\mathsf{w}$, $\mathsf{w}^{\star}$, $q$, $q'$}{a vector space, a
  linear form, an endomorphism; functions; vectors of~$M$; tuples of cosets}
\nt{$\mathsf{t}$, $\coef{\mathsf{t}}{\psi}$, $R$, $\Xi$, $z$, $z^{\flat}$, $F$}{a functional; a coefficient; the
  last-stage map; forms; functions}
\nt{$g$, $h$, $h_0$, $k$, $u$, $p$, $n$, $i$, $j$, $e$}{points, integral points, unipotent and
  Bessel elements; natural numbers and indices; a bijection}
\nt{$\varsigma$}{in Sections~\ref{s:setting}, \ref{s:poincare}, \ref{s:lattices} and~\ref{s:cuts} also an element of the Bessel group $\Besf{\mathcal{Y}}$ (the chosen root of unity of Convention~0.2 does not occur in this note)}
\nt{$f$}{in Sections~\ref{s:setting}, \ref{s:chainproofs} and~\ref{s:lattices} also a vector of a Weil representation, a datum (the number of infinite places is $\fD{D}$, Convention~0.3)}
\nt{$\delta$, $\delta_a$, $\lambda_0$, $\mathsf{M}$, $\varepsilon$, $\kappa$}{a norm-one idele;
  exponents; an integral vector; matrices; a unit with \m{\cj\varepsilon\varepsilon=-1}; a family of weights}
\nt{$\mathsf{e}$, $u_0$, $\mathsf{g}$, $\mathsf{s}$, $S_0$, $\varkappa$, $\mathfrak{p}_w$, $\mathrm{M}_I(\mathfrak{o}_w)$}{in
  the Inputs and in single proofs: an integral adele, a projected vector, a representative, a map of parameters, a
  stabiliser, a number, the maximal ideal at~$w$, a matrix ring}
\nt{letters with two meanings}{$b$: the third scale of~$D$ (Setting~\ref{s:data} only); $\mathbf{b}$: an adapted basis
  (Section~\ref{s:lattices}), and with an index a member of a tuple of level elements; $d$: the prime of~$P$, and a
  denominator of an admissible combination; $t$: a coset, and a tuple of translations in $\Tr{D}$; $T$: a set of
  cosets, and with two gradings as indices the string map; $x$: the first scale of~$D$, and a point or a
  parameter; $\varrho$: the Heisenberg operators, and with $D$ the number $\rhoD{D}$; $P$: the parameters, the
  roman $\PoinD{D}$ is the Poincar\'e series and the fraktur one the stage functions; $w$: a finite place, the
  sans-serif $\mathsf{w}$ a vector; $\nu$: a character, and an index or exponent}
\nt{$\Theta_{D,\vartheta}$}{the chain of lifts on the rows, as in the overview}
\nt{$\mathfrak{X}$, $\mathfrak{m}$, $\mathfrak{m}_1$, $\Delta$, $\Gamma$, $\phi$, $Q$, $Q_0$, $U$, $B$, $G$,
  $Y$, $Z$}{level places; moduli; a datum; a subgroup; coefficients; sets of places and of planes; letters
  occurring only inside the names above}
\nt{$s$}{in Definition~\ref{d:poin} and Proposition~\ref{p:good} also a datum of a lift (the complex variable and the exponent of Convention~0.6 are not meant)}
\end{multicols}
\endgroup

\end{document}
